\chapter{Measure and Integration}\label{measure-theory}
\chapterstatus{complete}

\section{Introduction}\label{measure-theory:section-introduction}

The Lebesgue integral extends the Riemann integral to a class of functions that is closed under the limits needed in analysis: the
monotone and dominated convergence theorems (Theorems~\ref{measure-theory:theorem-monotone-convergence}
and~\ref{measure-theory:theorem-dominated-convergence}) and the completeness of the spaces $L^p$
(Theorem~\ref{measure-theory:theorem-riesz-fischer}) are the basis of modern analysis and probability. This chapter is the first of the part on measure theory. It starts with $\sigma$-algebras and Borel sets, measures and their
continuity properties, measurable functions and image measures. It then constructs measures from outer
measures (Carath\'eodory's theorem) and Lebesgue measure in particular, shows why not every set can be measurable (Vitali, and the paradoxes of Hausdorff and Banach--Tarski), proves the uniqueness
theorem for measures (Dynkin's $\pi$-$\lambda$ lemma) and the Carath\'eodory--Hahn extension theorem, completes measures, and constructs all locally finite Borel measures on the line from their
distribution functions (the Lebesgue--Stieltjes measures, Theorem~\ref{measure-theory:theorem-lebesgue-stieltjes}). It then defines the integral and
proves the convergence theorems, Jensen's inequality and the comparison with the Riemann integral; it studies the spaces $L^p$, the modes of convergence and Egorov's theorem, independence and the two
Borel--Cantelli lemmas with Borel's normal number theorem (Theorem~\ref{measure-theory:theorem-borel-normal}), Littlewood's principles and Lusin's theorem, product measures and the theorems of Tonelli and
Fubini with the layer cake formula, and the change of variables formula with Cavalieri's principle, volumes of solids of revolution and the theorems of Pappus--Guldin.

The next chapters continue the theory: Chapter~\ref{lebesgue-differentiation} differentiates integrals (Lebesgue's differentiation theorem, the Radon--Nikodym theorem, absolutely continuous
functions), Chapter~\ref{signed-measures} treats signed and complex measures, the duality of the spaces $L^p$ and uniform integrability, Chapter~\ref{radon-measures} studies measures on locally compact spaces (the Riesz--Markov theorem, weak convergence, differentiation of measures), and
Chapter~\ref{geometric-measure-theory} measures lower-dimensional sets (Hausdorff measures, fractals, the area and coarea formulas, rectifiability).
References: \cite{rudin-rca} and \cite{folland-real}.

\noindent\textbf{Prerequisites.} Chapter~\ref{real-analysis} (Real Analysis).

\section{$\sigma$-algebras}\label{measure-theory:section-sigma-algebras}

\emph{Why $\sigma$-algebras.} We want to assign a size to subsets of $\RR$ extending the length of intervals, additive on disjoint sets, and invariant under translation. Two features of the
definitions below need explanation. First, additivity is required for \emph{countable} disjoint unions, not only finite ones: this is what makes limits work (the convergence theorems of
Section~\ref{measure-theory:section-integral}), and it is natural, since for instance $[0, 1) = \bigsqcup_n[1 - 2^{-n}, 1 - 2^{-n-1})$ should have length $\sum_n2^{-n-1} = 1$. Second, the size is defined only
on a family of ``measurable'' sets, not on all subsets: by the Vitali construction (Counterexample~\ref{measure-theory:counterexample-vitali}) there is no countably additive, translation invariant length
defined on all subsets of $\RR$ and finite on bounded intervals. The family of measurable sets must be closed under the operations used in analysis, complements and countable unions; this is a
$\sigma$-algebra.

\begin{definition}\label{measure-theory:definition-sigma-algebra}
A \emph{$\sigma$-algebra} on a set $X$ is a family $\mathcal{A}$ of subsets of $X$ such that
\begin{enumerate}
\item $X \in \mathcal{A}$;
\item if $A \in \mathcal{A}$, then $X \setminus A \in \mathcal{A}$;
\item if $A_1, A_2, \ldots \in \mathcal{A}$, then $\bigcup_kA_k \in \mathcal{A}$.
\end{enumerate}
The pair $(X, \mathcal{A})$ is a \emph{measurable space}, and the members of $\mathcal{A}$ are its \emph{measurable sets}.
\end{definition}

\begin{lemma}\label{measure-theory:lemma-sigma-algebra-closure}
A $\sigma$-algebra $\mathcal{A}$ contains $\emptyset$ and is closed under finite unions, countable and finite intersections, differences $A \setminus B$ and symmetric differences
$A \mathbin{\triangle} B = (A \setminus B) \cup (B \setminus A)$. For a sequence $A_k \in \mathcal{A}$ it contains
\[
\limsup_kA_k = \bigcap_{n \geq 1}\bigcup_{k \geq n}A_k, \qquad \liminf_kA_k = \bigcup_{n \geq 1}\bigcap_{k \geq n}A_k,
\]
the sets of points lying in infinitely many $A_k$, respectively in all but finitely many $A_k$.
\end{lemma}

\begin{proof}
$\emptyset = X \setminus X$. A finite union $A_1 \cup \cdots \cup A_m$ is the countable union with $A_k = \emptyset$ for $k > m$. By De Morgan, $\bigcap_kA_k = X \setminus \bigcup_k(X \setminus A_k)$.
Then $A \setminus B = A \cap (X \setminus B)$, and the symmetric difference is a union of two differences. The formulas for $\limsup$ and $\liminf$ combine countable unions and intersections. A point
lies in $\bigcup_{k \geq n}A_k$ for every $n$ if and only if it lies in $A_k$ for arbitrarily large $k$; and it lies in some $\bigcap_{k \geq n}A_k$ if and only if it lies in all $A_k$ from some
index on.
\end{proof}

\begin{example}\label{measure-theory:example-sigma-algebras}
\begin{enumerate}
\item On any set $X$, $\{\emptyset, X\}$ and the power set $\mathcal{P}(X)$ are $\sigma$-algebras, the smallest and the largest. If $X = A_1 \sqcup \cdots \sqcup A_k$ is a finite partition, the unions of some of
the $A_i$ form a $\sigma$-algebra with $2^k$ elements; on a finite set every $\sigma$-algebra is of this form, generated by its minimal nonempty members. The same holds for a countable partition
$X = \bigsqcup_kA_k$: the unions $\bigcup_{k \in S}A_k$, $S \subseteq \NN$, form a $\sigma$-algebra, in bijection with $\mathcal{P}(\NN)$.
\item The sets that are countable or have countable complement form a $\sigma$-algebra on any set: complements exchange the two kinds, and a countable union is countable unless one of the sets is
cocountable, in which case the union is cocountable. On $\RR$ it is strictly smaller than the Borel $\sigma$-algebra (defined below), since $[0, 1]$ is neither countable nor cocountable.
\item \emph{Trace.} If $\mathcal{A}$ is a $\sigma$-algebra on $X$ and $Y \subseteq X$, then $\mathcal{A}|_Y = \{A \cap Y : A \in \mathcal{A}\}$ is a $\sigma$-algebra on $Y$, since
$Y \setminus (A \cap Y) = (X \setminus A) \cap Y$ and unions commute with intersecting with $Y$. If $Y \in \mathcal{A}$, it consists of the measurable subsets of $Y$.
\item \emph{Preimages and images.} If $f \colon X \to Y$ is a map and $\mathcal{B}$ a $\sigma$-algebra on $Y$, then $f^{-1}(\mathcal{B}) = \{f^{-1}(B) : B \in \mathcal{B}\}$ is a $\sigma$-algebra on $X$. If
$\mathcal{A}$ is a $\sigma$-algebra on $X$, then $\{B \subseteq Y : f^{-1}(B) \in \mathcal{A}\}$ is a $\sigma$-algebra on $Y$. Both statements hold because preimages commute with complements and
arbitrary unions: $f^{-1}(Y \setminus B) = X \setminus f^{-1}(B)$ and $f^{-1}(\bigcup_kB_k) = \bigcup_kf^{-1}(B_k)$. Images do not commute with complements, and $\{f(A) : A \in \mathcal{A}\}$ is not a
$\sigma$-algebra in general (take $f$ constant and $Y$ with two points).
\end{enumerate}
\end{example}

\begin{counterexample}\label{measure-theory:counterexample-union-sigma-algebras}
An increasing union of $\sigma$-algebras need not be a $\sigma$-algebra. On $\NN = \{1, 2, \ldots\}$ let $\mathcal{A}_n$ consist of the sets $A$ such that $A \cap \{n + 1, n + 2, \ldots\}$ is either empty or
all of $\{n + 1, n + 2, \ldots\}$; it is a $\sigma$-algebra (the unions of members of the partition $\{1\}, \ldots, \{n\}, \{n + 1, n + 2, \ldots\}$), and $\mathcal{A}_n \subseteq \mathcal{A}_{n+1}$. The
union $\bigcup_n\mathcal{A}_n$ contains every singleton, but not the set of even numbers, which is a countable union of singletons. The union is still closed under complements and finite unions: it is an
algebra (Definition~\ref{measure-theory:definition-premeasure}).
\end{counterexample}

Usually a $\sigma$-algebra is not described by listing its members but by naming some sets that must be measurable.

\begin{definition}\label{measure-theory:definition-generated-sigma-algebra}
Let $\mathcal{E}$ be a family of subsets of $X$. The \emph{$\sigma$-algebra generated by $\mathcal{E}$}, written $\sigma(\mathcal{E})$, is the intersection of all $\sigma$-algebras on $X$ containing
$\mathcal{E}$. For a topological space $X$, the \emph{Borel $\sigma$-algebra} $\mathcal{B}(X)$ is the $\sigma$-algebra generated by the open sets; its members are the \emph{Borel sets}.
\end{definition}

The intersection is over a nonempty family, since $\mathcal{P}(X)$ is one of its members, and an intersection of $\sigma$-algebras is a $\sigma$-algebra (each of the three conditions is preserved). So
$\sigma(\mathcal{E})$ is the smallest $\sigma$-algebra containing $\mathcal{E}$: it is contained in every $\sigma$-algebra that contains $\mathcal{E}$. This minimality is the usual tool. To show that all
sets of $\sigma(\mathcal{E})$ have a property, one shows that the sets with the property form a $\sigma$-algebra containing $\mathcal{E}$; and if $\mathcal{E} \subseteq \sigma(\mathcal{E}^{\prime})$ then
$\sigma(\mathcal{E}) \subseteq \sigma(\mathcal{E}^{\prime})$. The Borel $\sigma$-algebra contains the open and the closed sets, the countable unions of closed sets (called $F_\sigma$ sets) and the countable
intersections of open sets ($G_\delta$ sets), countable intersections of $F_\sigma$ sets, and so on.

\begin{proposition}\label{measure-theory:proposition-borel-generators}
\begin{enumerate}
\item The Borel $\sigma$-algebra of $\RR$ is generated by each of the following families: the open intervals $(a, b)$; the closed intervals $[a, b]$; the half-open intervals $(a, b]$; the rays
$(a, \infty)$, $a \in \RR$; the rays $(-\infty, a]$, $a \in \QQ$.
\item The Borel $\sigma$-algebra of $\RR^n$ is generated by the open boxes $\prod_i(a_i, b_i)$ with rational $a_i < b_i$, and also by the half-spaces $\{x : x_i \leq c\}$, $1 \leq i \leq n$, $c \in \RR$.
\item The Borel $\sigma$-algebra of $[-\infty, \infty]$ (with the order topology) is generated by the rays $(a, \infty]$, $a \in \RR$.
\end{enumerate}
\end{proposition}

\begin{proof}
\emph{Step 1: open sets are countable unions of rational boxes.} Let $U \subseteq \RR^n$ be open and $x \in U$. Some ball around $x$ lies in $U$, and inside it there is an open box with rational vertices
containing $x$. There are countably many rational boxes, so $U$ is the countable union of those contained in it. Hence the open sets, and so all Borel sets, lie in the $\sigma$-algebra generated by the
rational boxes; conversely the boxes are open. This proves the first statement of (2), and for $n = 1$ that the open intervals generate $\mathcal{B}(\RR)$.

\emph{Step 2: the other families on $\RR$.} Every family in (1) consists of Borel sets, so it generates a $\sigma$-algebra contained in $\mathcal{B}(\RR)$; it suffices to show that the generated
$\sigma$-algebra contains the open intervals. We have $(a, b) = \bigcup_k[a + 1/k, b - 1/k]$ and $(a, b) = \bigcup_k(a, b - 1/k]$; also $(a, b] = (a, \infty) \setminus (b, \infty)$, so the rays
$(a, \infty)$ generate the half-open intervals; and for real $a$, $(-\infty, a] = \bigcap_k(-\infty, q_k]$ for rationals $q_k \downarrow a$, and $(a, \infty) = \RR \setminus (-\infty, a]$.

\emph{Step 3: half-spaces.} A half-space $\{x_i \leq c\}$ is closed, hence Borel. Conversely $\{a_i < x_i < b_i\} = \bigcup_k\{x_i \leq b_i - 1/k\} \setminus \{x_i \leq a_i\}$, and an open box is the
intersection of $n$ such slabs.

\emph{Step 4: $[-\infty, \infty]$.} The open sets of $[-\infty, \infty]$ are countable unions of intervals $(a, b)$, $[-\infty, b)$ and $(a, \infty]$ with rational $a$, $b$, by the argument of Step 1.
Each of these lies in the $\sigma$-algebra generated by the rays $(a, \infty]$: $[-\infty, b) = \bigcup_k([-\infty, \infty] \setminus (b - 1/k, \infty])$ and $(a, b) = (a, \infty] \cap [-\infty, b)$.
\end{proof}

\begin{remark}\label{measure-theory:remark-borel-hierarchy}
The Borel $\sigma$-algebra is much larger than the open and closed sets, and building it ``from below'' requires transfinitely many steps. Already the first steps are distinct: $\QQ$ is an $F_\sigma$ set
(a countable union of points), but not a $G_\delta$ set. Indeed, if $\QQ = \bigcap_kU_k$ with $U_k$ open, each $U_k$ is dense, and so are the open sets $U_k \setminus \{q_k\}$ for an enumeration
$(q_k)$ of $\QQ$; their intersection is empty, contradicting the Baire category theorem (Theorem~\ref{general-topology:theorem-baire}) in the complete metric space $\RR$. Consequently the irrationals form a
$G_\delta$ set that is not $F_\sigma$.
\end{remark}

\section{Measures}\label{measure-theory:section-measures}

\begin{definition}\label{measure-theory:definition-measure}
A \emph{measure} on a measurable space $(X, \mathcal{A})$ is a map $\mu \colon \mathcal{A} \to [0, \infty]$ with $\mu(\emptyset) = 0$ which is \emph{countably additive}: $\mu(\bigsqcup_kA_k) = \sum_k\mu(A_k)$
for every sequence of pairwise disjoint $A_k \in \mathcal{A}$. The triple $(X, \mathcal{A}, \mu)$ is a \emph{measure space}. The measure is
\begin{enumerate}
\item \emph{finite} if $\mu(X) < \infty$, and a \emph{probability measure} if $\mu(X) = 1$;
\item \emph{$\sigma$-finite} if $X$ is a countable union of measurable sets of finite measure;
\item \emph{complete} if every subset of a measurable set of measure zero is measurable.
\end{enumerate}
A measurable set of measure zero is a \emph{null set}. A property of points of $X$ holds \emph{almost everywhere} (a.e., or $\mu$-a.e.) if the set of points where it fails is contained in a null set.
\end{definition}

\begin{notation}\label{measure-theory:notation-measures}
We use the extended half-line $[0, \infty]$ with $a + \infty = \infty$ for all $a$ and $0 \cdot \infty = 0$; a series of terms in $[0, \infty]$ always has a sum in $[0, \infty]$, independent of the
order of the terms. We write $\bigsqcup$ for a union of pairwise disjoint sets. When the $\sigma$-algebra is clear, we speak of the measure space $(X, \mu)$.
\end{notation}

Taking $A_k = \emptyset$ for $k > m$ shows that a measure is \emph{finitely additive}. In a probability space, the measurable sets are \emph{events} and $\mu(A)$ is the probability of $A$; countable additivity
lets one compute probabilities of events described by infinitely many conditions, such as ``a six eventually appears'' in an infinite sequence of throws of a die.

\begin{example}\label{measure-theory:example-measures}
\begin{enumerate}
\item \emph{Counting measure} on $\mathcal{P}(X)$: $\mu(A)$ is the number of elements of $A$, or $\infty$. It is $\sigma$-finite if and only if $X$ is countable.
\item \emph{Dirac measure} at $x$: $\delta_x(A) = 1$ if $x \in A$ and $0$ otherwise; more generally $\sum_ic_i\delta_{x_i}$ with $c_i \geq 0$, a \emph{discrete} measure. A fair die is the probability measure
$\frac16\sum_{k=1}^6\delta_k$ on $\{1, \ldots, 6\}$, and the probability of an even result is $\mu(\{2, 4, 6\}) = \frac12$.
\item If $\mu$ is a measure and $E$ a measurable set, $\mu_E(A) = \mu(A \cap E)$ is a measure, the restriction of $\mu$ to $E$; and if $f \geq 0$ is measurable, $A \mapsto \int_Af\,d\mu$ is a measure
(Section~\ref{measure-theory:section-integral}), the measure with \emph{density} $f$.
\item Lebesgue measure on $\RR^n$ (Section~\ref{measure-theory:section-lebesgue}), which assigns to boxes their volume.
\item \emph{The cocountable measure.} On an uncountable set $X$ with the $\sigma$-algebra of countable and cocountable sets (Example~\ref{measure-theory:example-sigma-algebras}(2)), let $\mu(A) = 0$ for
countable $A$ and $\mu(A) = 1$ for cocountable $A$. Two cocountable sets intersect (their union of complements is countable, so not all of $X$), so a disjoint sequence contains at most one cocountable
set, and countable additivity follows: the union is cocountable exactly when one of the sets is.
\item The \emph{trivial infinite measure} $\mu(A) = \infty$ for $A \neq \emptyset$, $\mu(\emptyset) = 0$, is a measure on any $\sigma$-algebra; it is not $\sigma$-finite unless $X = \emptyset$.
\item Sums and multiples: if $\mu_1, \mu_2, \ldots$ are measures on $\mathcal{A}$ and $c_k \geq 0$, then $\sum_kc_k\mu_k$ is a measure, since sums of nonnegative double series may be rearranged.
\end{enumerate}
\end{example}

\begin{proposition}\label{measure-theory:proposition-measure-properties}
Let $\mu$ be a measure. If $A \subseteq B$, then $\mu(A) \leq \mu(B)$; $\mu(\bigcup_kA_k) \leq \sum_k\mu(A_k)$; if $A_1 \subseteq A_2 \subseteq \cdots$, then $\mu(\bigcup_kA_k) = \lim_k\mu(A_k)$; and if
$A_1 \supseteq A_2 \supseteq \cdots$ with $\mu(A_1) < \infty$, then $\mu(\bigcap_kA_k) = \lim_k\mu(A_k)$.
\end{proposition}

The four statements are called \emph{monotonicity}, \emph{countable subadditivity}, \emph{continuity from below} and \emph{continuity from above}.

\begin{proof}
\emph{Monotonicity.} $B = A \sqcup (B \setminus A)$, so $\mu(B) = \mu(A) + \mu(B \setminus A) \geq \mu(A)$. If $\mu(A) < \infty$ this also gives $\mu(B \setminus A) = \mu(B) - \mu(A)$.

\emph{Subadditivity.} The sets $B_k = A_k \setminus \bigcup_{j<k}A_j$ are measurable, pairwise disjoint, contained in $A_k$, and have the same union as the $A_k$ (a point of the union lies in $B_k$ for
the least $k$ with $x \in A_k$). So $\mu(\bigcup_kA_k) = \sum_k\mu(B_k) \leq \sum_k\mu(A_k)$.

\emph{Continuity from below.} With $A_0 = \emptyset$, $\bigcup_kA_k = \bigsqcup_k(A_k \setminus A_{k-1})$, and by finite additivity the partial sums of $\sum_k\mu(A_k \setminus A_{k-1})$ are
$\mu(A_k)$.

\emph{Continuity from above.} The sets $A_1 \setminus A_k$ increase to $A_1 \setminus \bigcap_kA_k$. By continuity from below and the formula for differences (all measures being finite),
$\mu(A_1) - \mu(A_k) \to \mu(A_1) - \mu(\bigcap_kA_k)$.
\end{proof}

Countable additivity forces continuity properties that finite additivity does not give. For instance, for counting measure on $\NN$ the sets $A_n = \{n, n + 1, \ldots\}$ decrease to $\emptyset$ while
$\mu(A_n) = \infty$ for all $n$: the hypothesis $\mu(A_1) < \infty$ in continuity from above cannot be dropped. Conversely, continuity characterises countable additivity: a finitely additive
$\mu \colon \mathcal{A} \to [0, \infty]$ that is continuous from below is countably additive, since $\mu(\bigsqcup_kA_k) = \lim_m\mu(A_1 \sqcup \cdots \sqcup A_m) = \lim_m\sum_{k \leq m}\mu(A_k)$.

\begin{proposition}[Inclusion--exclusion]\label{measure-theory:proposition-inclusion-exclusion}
For measurable sets $A$, $B$, $\mu(A \cup B) + \mu(A \cap B) = \mu(A) + \mu(B)$. If $\mu(A_1 \cup \cdots \cup A_m) < \infty$, then
\[
\mu\Bigl(\bigcup_{i=1}^mA_i\Bigr) = \sum_{\emptyset \neq S \subseteq \{1, \ldots, m\}}(-1)^{|S| - 1}\mu\Bigl(\bigcap_{i \in S}A_i\Bigr).
\]
\end{proposition}

\begin{proof}
$A \cup B = (A \setminus B) \sqcup (A \cap B) \sqcup (B \setminus A)$, while $A = (A \setminus B) \sqcup (A \cap B)$ and $B = (B \setminus A) \sqcup (A \cap B)$; add. For the general formula, it suffices to
show that both sides assign the same weight to every atom of the partition generated by $A_1, \ldots, A_m$: for $T \subseteq \{1, \ldots, m\}$ nonempty, let $C_T$ be the set of points lying in $A_i$
exactly for $i \in T$. The $C_T$ are measurable, disjoint, with union $\bigcup_iA_i$, and $\bigcap_{i \in S}A_i = \bigsqcup_{T \supseteq S}C_T$. So the right side is
$\sum_T\mu(C_T)\sum_{\emptyset \neq S \subseteq T}(-1)^{|S| - 1}$, and the inner sum is $1 - \sum_{S \subseteq T}(-1)^{|S|} = 1 - (1 - 1)^{|T|} = 1$ by the binomial theorem. All terms are finite, so the
rearrangements are allowed.
\end{proof}

\begin{example}\label{measure-theory:example-inclusion-exclusion}
A permutation of $\{1, \ldots, m\}$ chosen uniformly at random has no fixed point with probability $\sum_{k=0}^m(-1)^k/k! \to 1/e$. Indeed, let $A_i$ be the event that $i$ is fixed; for $|S| = k$,
$\mu(\bigcap_{i \in S}A_i) = (m - k)!/m!$, and there are $\binom mk$ such $S$, so by the proposition $\mu(\bigcup_iA_i) = \sum_{k=1}^m(-1)^{k-1}/k!$. The finiteness hypothesis is needed: for counting
measure on $\NN$ and $A = B = \NN$ the first formula reads $\infty + \infty = \infty + \infty$, but one cannot subtract to get $\mu(A \cup B) = \mu(A) + \mu(B) - \mu(A \cap B)$.
\end{example}

\section{Measurable functions}\label{measure-theory:section-measurable-functions}

\begin{definition}\label{measure-theory:definition-measurable-function}
A map $f \colon (X, \mathcal{A}) \to (Y, \mathcal{B})$ between measurable spaces is \emph{measurable} if $f^{-1}(B) \in \mathcal{A}$ for every $B \in \mathcal{B}$. Functions with values in $[-\infty, \infty]$,
$\RR$ or $\CC$ are understood with the Borel $\sigma$-algebras; a measurable function between topological spaces with their Borel $\sigma$-algebras is a \emph{Borel function}. A \emph{simple function} is a
finite linear combination $\sum_ia_i\mathbf{1}_{A_i}$ of indicator functions of measurable sets, with real or complex coefficients; equivalently, a measurable function taking finitely many values.
\end{definition}

Measurability of functions is the condition that makes integrals of $f$ meaningful: one needs to measure the sets $\{a < f \leq b\}$, on which $f$ is approximately constant. It plays the role that
continuity plays in topology, with measurable sets in place of open sets. As there, it suffices to check generators.

\begin{lemma}\label{measure-theory:lemma-measurable-generators}
Let $\mathcal{B} = \sigma(\mathcal{E})$. A map $f \colon (X, \mathcal{A}) \to (Y, \mathcal{B})$ is measurable if and only if $f^{-1}(E) \in \mathcal{A}$ for every $E \in \mathcal{E}$. Consequently continuous
maps between topological spaces are Borel, and composites of measurable maps are measurable.
\end{lemma}

\begin{proof}
The sets $B \subseteq Y$ with $f^{-1}(B) \in \mathcal{A}$ form a $\sigma$-algebra (Example~\ref{measure-theory:example-sigma-algebras}(4)); if it contains $\mathcal{E}$, it contains $\sigma(\mathcal{E})$. For a
continuous map, preimages of open sets are open, hence Borel. If $f$ and $g$ are measurable, $(g \circ f)^{-1}(C) = f^{-1}(g^{-1}(C))$.
\end{proof}

\begin{example}\label{measure-theory:example-measurable-functions}
\begin{enumerate}
\item An indicator function $\mathbf{1}_A$ is measurable if and only if $A$ is measurable, since $\{\mathbf{1}_A > a\}$ is $\emptyset$, $A$ or $X$. The Dirichlet function $\mathbf{1}_\QQ$, which is nowhere
continuous, is Borel, $\QQ$ being countable.
\item A monotone function $f \colon \RR \to \RR$ is Borel, since $\{f > a\}$ is an interval (Proposition~\ref{measure-theory:proposition-measurable-functions}(1) below). So is a function continuous outside
a countable set, such as $x \mapsto \lfloor x\rfloor$.
\item The derivative of a differentiable $f \colon \RR \to \RR$ is Borel, being the pointwise limit of the continuous functions $n(f(x + 1/n) - f(x))$
(Proposition~\ref{measure-theory:proposition-measurable-functions}(3) below), although it need not be continuous: for $f(x) = x^2\sin(1/x)$, $f(0) = 0$, the derivative exists everywhere and is discontinuous at $0$.
\item If $\mathcal{A} = \{\emptyset, X\}$, only the constant functions are measurable; if $\mathcal{A} = \mathcal{P}(X)$, every function is.
\end{enumerate}
\end{example}

\begin{proposition}\label{measure-theory:proposition-measurable-functions}
Let $(X, \mathcal{A})$ be a measurable space.
\begin{enumerate}
\item $f \colon X \to [-\infty, \infty]$ is measurable if and only if $\{f > a\}$ is measurable for all $a \in \RR$; the same holds with $\{f \geq a\}$, $\{f < a\}$ or $\{f \leq a\}$ in place of $\{f > a\}$.
\item If $f_1, \ldots, f_k \colon X \to \RR$ are measurable and $\Phi \colon \RR^k \to \RR$ is continuous, $\Phi(f_1, \ldots, f_k)$ is measurable; so sums and products of measurable functions are
measurable, and $f \colon X \to \CC$ is measurable if and only if its real and imaginary parts are.
\item Pointwise $\sup$, $\inf$, $\limsup$ and $\liminf$ of sequences of measurable functions $X \to [-\infty, \infty]$ are measurable; so is the pointwise limit where it exists, and the set where a
sequence of real measurable functions converges is measurable.
\item Every measurable $f \colon X \to [0, \infty]$ is the pointwise limit of an increasing sequence of simple functions $0 \leq s_1 \leq s_2 \leq \cdots$; the convergence is uniform on every set where $f$ is
bounded. Every measurable $f \colon X \to \CC$ is the pointwise limit of simple functions $s_n$ with $|s_n| \leq |f|$.
\item For measurable $f, g \colon X \to [-\infty, \infty]$, the functions $\max(f, g)$, $\min(f, g)$, $f^+ = \max(f, 0)$, $f^- = \max(-f, 0)$ and $|f|$ are measurable, and the sets $\{f < g\}$ and
$\{f = g\}$ are measurable.
\end{enumerate}
\end{proposition}

\begin{proof}
(1) The rays $(a, \infty]$ generate $\mathcal{B}([-\infty, \infty])$ (Proposition~\ref{measure-theory:proposition-borel-generators}(3)), so Lemma~\ref{measure-theory:lemma-measurable-generators} applies. The other
families: $\{f \geq a\} = \bigcap_k\{f > a - 1/k\}$, $\{f > a\} = \bigcup_k\{f \geq a + 1/k\}$, and $\{f < a\}$, $\{f \leq a\}$ are the complements.

(2) The map $F = (f_1, \ldots, f_k) \colon X \to \RR^k$ is measurable: by Proposition~\ref{measure-theory:proposition-borel-generators}(2) it suffices to consider open boxes, and
$F^{-1}(\prod_i(a_i, b_i)) = \bigcap_if_i^{-1}((a_i, b_i))$. Then $\Phi \circ F$ is measurable by Lemma~\ref{measure-theory:lemma-measurable-generators}. For complex $f$, identify $\CC$ with $\RR^2$:
$f = (\operatorname{Re}f, \operatorname{Im}f)$, and $\operatorname{Re}$, $\operatorname{Im}$ are continuous.

(3) $\{\sup_nf_n > a\} = \bigcup_n\{f_n > a\}$; $\inf_nf_n = -\sup_n(-f_n)$; $\limsup_nf_n = \inf_n\sup_{k \geq n}f_k$ and $\liminf_nf_n = \sup_n\inf_{k \geq n}f_k$. The limit, where it exists, equals
$\limsup_nf_n$. For real $f_n$, the set of convergence is $\{\limsup_nf_n = \liminf_nf_n\} \cap \{\limsup_nf_n \in \RR\}$, measurable by (5).

(4) Let $s_n = \min(n, 2^{-n}\lfloor 2^nf\rfloor)$, with $s_n = n$ where $f = \infty$. Each $s_n$ takes finitely many values, on the measurable sets $\{k2^{-n} \leq f < (k + 1)2^{-n}\}$ and
$\{f \geq n\}$. Passing from $n$ to $n + 1$ refines the dyadic grid, so $s_n \leq s_{n+1}$; and where $f(x) < n$, $0 \leq f(x) - s_n(x) < 2^{-n}$. Hence $s_n \to f$ pointwise, uniformly on
$\{f \leq M\}$. For complex $f$ apply this to $(\operatorname{Re}f)^\pm$ and $(\operatorname{Im}f)^\pm$, obtaining simple $0 \leq s_n^{(j)} \uparrow$ the four parts, and put
$s_n = s_n^{(1)} - s_n^{(2)} + i(s_n^{(3)} - s_n^{(4)})$; at each point only one of each pair is nonzero, so $|\operatorname{Re}s_n| \leq |\operatorname{Re}f|$, $|\operatorname{Im}s_n| \leq |\operatorname{Im}f|$
and $|s_n| \leq |f|$.

(5) $\{\max(f, g) > a\} = \{f > a\} \cup \{g > a\}$ and $\min(f, g) = -\max(-f, -g)$; then $f^\pm$ and $|f| = f^+ + f^-$ (a sum without $\infty - \infty$). Finally
$\{f < g\} = \bigcup_{q \in \QQ}\{f < q\} \cap \{q < g\}$, and $\{f = g\}$ is the complement of $\{f < g\} \cup \{g < f\}$.
\end{proof}

\begin{counterexample}\label{measure-theory:counterexample-measurable-functions}
\begin{enumerate}
\item Measurability of $|f|$ does not imply that of $f$: on $X = \{1, 2\}$ with $\mathcal{A} = \{\emptyset, X\}$, $f(1) = 1$, $f(2) = -1$ is not measurable but $|f| = 1$ is.
\item The supremum of an uncountable family of measurable functions need not be measurable: for a non-Lebesgue-measurable set $N \subseteq \RR$ (Counterexample~\ref{measure-theory:counterexample-vitali}),
$\mathbf{1}_N = \sup_{x \in N}\mathbf{1}_{\{x\}}$ is a supremum of Borel functions, but it is not Lebesgue measurable. Countability in (3) is essential.
\item A composite $g \circ f$ of a Lebesgue measurable $g \colon \RR \to \RR$ (that is, $g^{-1}(B)$ Lebesgue measurable for Borel $B$) and a continuous $f$ need not be Lebesgue measurable. The
definition requires preimages of all Borel sets, not of all Lebesgue measurable sets, and a continuous map can pull a Lebesgue null set back to a non-measurable set; an example built from the Cantor
function is in Chapter~\ref{lebesgue-differentiation}. This is why Borel functions, and not Lebesgue measurable ones, are the natural class to compose.
\end{enumerate}
\end{counterexample}

A measurable map transports measures.

\begin{definition}\label{measure-theory:definition-image-measure}
Let $f \colon (X, \mathcal{A}) \to (Y, \mathcal{B})$ be measurable and $\mu$ a measure on $\mathcal{A}$. The \emph{image measure} (or \emph{push-forward}) of $\mu$ under $f$ is the measure
$f_*\mu$ on $\mathcal{B}$ given by $f_*\mu(B) = \mu(f^{-1}(B))$. If $\mu$ is a probability measure and $f$ is real-valued, $f_*\mu$ is the \emph{distribution} of $f$.
\end{definition}

It is a measure because preimages preserve disjointness and unions. For example, if $f \colon \{1, \ldots, 6\}^2 \to \{2, \ldots, 12\}$ is the sum of two dice and $\mu$ is the uniform measure, then
$f_*\mu(\{7\}) = \frac{6}{36}$; and for Lebesgue measure on $\RR$ and $f(x) = ax + b$ with $a \neq 0$, $f_*\lambda = |a|^{-1}\lambda$ (Theorem~\ref{measure-theory:theorem-lebesgue-measure} and
Example~\ref{measure-theory:example-distribution-function}: both measures give $(c, d]$ the measure $(d - c)/|a|$). The formula $\int g\,d(f_*\mu) = \int g \circ f\,d\mu$ is
Proposition~\ref{measure-theory:proposition-image-measure} below.

\section{Outer measures and Lebesgue measure}\label{measure-theory:section-lebesgue}

\begin{definition}\label{measure-theory:definition-outer-measure}
An \emph{outer measure} on $X$ is a map $\mu^* \colon \mathcal{P}(X) \to [0, \infty]$ with $\mu^*(\emptyset) = 0$, monotone and countably subadditive. A set $A$ is \emph{$\mu^*$-measurable} if
$\mu^*(E) = \mu^*(E \cap A) + \mu^*(E \setminus A)$ for all $E \subseteq X$.
\end{definition}

\begin{theorem}[Carath\'eodory]\label{measure-theory:theorem-caratheodory}
The $\mu^*$-measurable sets form a $\sigma$-algebra $\mathcal{M}$, and $\mu^*|_\mathcal{M}$ is a complete measure.
\end{theorem}

\begin{proof}
By subadditivity, $A$ is measurable if and only if $\mu^*(E \cap A) + \mu^*(E \setminus A) \leq \mu^*(E)$ for all $E$. The definition is symmetric in $A$ and $X \setminus A$. For measurable $A, B$ and any $E$,
\[\begin{aligned}
\mu^*(E \cap (A \cup B)) + \mu^*(E \setminus (A \cup B)) &\leq \mu^*(E \cap A) + \mu^*((E \setminus A) \cap B) \\
&\quad + \mu^*((E \setminus A) \setminus B) \\
&= \mu^*(E \cap A) + \mu^*(E \setminus A) = \mu^*(E),
\end{aligned}\]
so $\mathcal{M}$ is closed under finite unions. Let $A_1, A_2, \ldots \in \mathcal{M}$ be disjoint, $B_n = A_1 \cup \cdots \cup A_n$ and $B = \bigcup_nA_n$. Measurability of $A_n$ gives
$\mu^*(E \cap B_n) = \mu^*(E \cap A_n) + \mu^*(E \cap B_{n-1})$, so $\mu^*(E \cap B_n) = \sum_{j \leq n}\mu^*(E \cap A_j)$, and
$\mu^*(E) = \mu^*(E \cap B_n) + \mu^*(E \setminus B_n) \geq \sum_{j \leq n}\mu^*(E \cap A_j) + \mu^*(E \setminus B)$. Letting $n \to \infty$ and using subadditivity,
$\mu^*(E) \geq \sum_j\mu^*(E \cap A_j) + \mu^*(E \setminus B) \geq \mu^*(E \cap B) + \mu^*(E \setminus B)$. So $B \in \mathcal{M}$, and $E = B$ gives countable additivity. General countable unions reduce
to disjoint ones. If $\mu^*(A) = 0$, then $\mu^*(E \cap A) + \mu^*(E \setminus A) \leq \mu^*(E)$, so $A$ is measurable.
\end{proof}

\begin{definition}\label{measure-theory:definition-lebesgue}
A \emph{box} in $\RR^n$ is a product $Q = \prod_i[a_i, b_i]$ with volume $v(Q) = \prod_i(b_i - a_i)$. The \emph{Lebesgue outer measure} of $A \subseteq \RR^n$ is
$\lambda^*(A) = \inf\{\sum_kv(Q_k) : A \subseteq \bigcup_kQ_k\}$, the infimum over countable covers by boxes. The $\lambda^*$-measurable sets are \emph{Lebesgue measurable}, and $\lambda = \lambda^n$ is
\emph{Lebesgue measure}.
\end{definition}

\begin{theorem}\label{measure-theory:theorem-lebesgue-measure}
$\lambda^*$ is an outer measure, and $\lambda^*(Q) = v(Q)$ for boxes. Every Borel set is Lebesgue measurable, so $\lambda$ is a complete measure on a $\sigma$-algebra containing $\mathcal{B}(\RR^n)$. It is
translation invariant and \emph{regular}: for measurable $E$, $\lambda(E) = \inf\{\lambda(U) : U \supseteq E \text{ open}\} = \sup\{\lambda(K) : K \subseteq E \text{ compact}\}$.
\end{theorem}

\begin{proof}
Countable subadditivity: given covers of $A_k$ with total volume $\leq \lambda^*(A_k) + \varepsilon2^{-k}$, their union covers $\bigcup_kA_k$.

\emph{Volume of a box.} Clearly $\lambda^*(Q) \leq v(Q)$. Let $Q \subseteq \bigcup_kQ_k$ and $\varepsilon > 0$. Enlarge each $Q_k$ to a box $Q'_k$ whose interior contains $Q_k$, with
$v(Q'_k) \leq v(Q_k) + \varepsilon2^{-k}$. By compactness finitely many interiors cover $Q$, say those of $Q'_1, \ldots, Q'_N$. The endpoints of $Q$ and of the $Q'_k$ in each coordinate define a grid
dividing $Q$ into boxes $C$ with disjoint interiors, and $v(Q) = \sum_Cv(C)$ (volume is multiplicative and length is additive). The centre of each $C$ lies in some $Q'_k$, and then $C \subseteq Q'_k$ because
no grid line crosses the interior of $C$. The cells contained in $Q'_k$ have disjoint interiors, so their volumes add up to at most $v(Q'_k)$. Hence $v(Q) \leq \sum_kv(Q'_k) \leq \sum_kv(Q_k) + \varepsilon$.

\emph{Borel sets.} The half-spaces $H = \{x_i \leq c\}$ generate $\mathcal{B}(\RR^n)$, since open sets are countable unions of boxes with rational vertices. Given $E$ and a cover by boxes $Q_k$ with
$\sum v(Q_k) \leq \lambda^*(E) + \varepsilon$, the boxes $Q_k \cap H$ and $Q_k \cap \{x_i \geq c\}$ cover $E \cap H$ and $E \setminus H$ and have total volume $\sum_kv(Q_k)$. So
$\lambda^*(E \cap H) + \lambda^*(E \setminus H) \leq \lambda^*(E) + \varepsilon$, and $H$ is measurable by Theorem~\ref{measure-theory:theorem-caratheodory}.

\emph{Invariance and regularity.} Translation maps boxes to boxes of the same volume. Enlarging the covering boxes slightly into open boxes gives outer regularity. For inner regularity, first let $E$ be bounded, contained in a
closed box $Q$. Given $\varepsilon$, choose open $U \supseteq Q \setminus E$ with $\lambda(U) \leq \lambda(Q \setminus E) + \varepsilon$; then $K = Q \setminus U \subseteq E$ is compact with
$\lambda(E \setminus K) \leq \lambda(U \setminus (Q \setminus E)) \leq \varepsilon$. For general $E$, apply this to $E \cap [-m, m]^n$ and use continuity from below.
\end{proof}

\begin{definition}\label{measure-theory:definition-regular}
Let $\mu$ be a measure on the Borel sets of a topological space $X$ (or on a $\sigma$-algebra containing them). It is \emph{locally finite} if every point has a neighbourhood of finite measure;
\emph{outer regular} if $\mu(E) = \inf\{\mu(U) : U \supseteq E \text{ open}\}$, \emph{inner regular} if $\mu(E) = \sup\{\mu(K) : K \subseteq E \text{ compact}\}$, for all measurable $E$, and
\emph{regular} if it is both. A locally finite regular measure on a locally compact Hausdorff space is a \emph{Radon measure}.
\end{definition}

\begin{example}\label{measure-theory:example-regular}
Lebesgue measure is locally finite and regular (Theorem~\ref{measure-theory:theorem-lebesgue-measure}), so it is a Radon measure, and so is a Dirac measure $\delta_x$. Counting measure on $\ZZ$ is
Radon, but counting measure on $\RR$ is neither locally finite nor outer regular: $\{0\}$ has measure $1$, while every open set containing $0$ is infinite and has measure $\infty$. The
$1$-dimensional Hausdorff measure on $\RR^2$ (Chapter~\ref{geometric-measure-theory}) is not locally finite either, since every open set has infinite length.
\end{example}

\begin{counterexample}\label{measure-theory:counterexample-vitali}
There are subsets of $\RR$ that are not Lebesgue measurable. Choose (by the axiom of choice) a set $V \subseteq [0, 1]$ containing exactly one element of each class of $\RR/\QQ$. The translates $V + q$, $q \in \QQ \cap [-1, 1]$,
are disjoint and their union lies between $[0, 1]$ and $[-1, 2]$. If $V$ were measurable, translation invariance and countable additivity would give $1 \leq \sum_q\lambda(V) \leq 3$, which is impossible
whether $\lambda(V) = 0$ or $\lambda(V) > 0$.
\end{counterexample}

\begin{proposition}\label{measure-theory:proposition-free-group-paradox}
Let $F$ be the free group on two generators $a$, $b$ (Chapter~\ref{groups}), and for $x \in \{a, a^{-1}, b, b^{-1}\}$ let $W(x)$ be the set of reduced words beginning with $x$. Then
\[\begin{aligned}
F &= \{e\} \sqcup W(a) \sqcup W(a^{-1}) \sqcup W(b) \sqcup W(b^{-1}), \\
F &= W(a) \sqcup aW(a^{-1}) = W(b) \sqcup bW(b^{-1}).
\end{aligned}\]
So $F$ is \emph{paradoxical}: it is the disjoint union of four pieces, two of which can be translated to give $F$ again, and so can the other two.
\end{proposition}

\begin{proof}
The first decomposition is the classification of reduced words by their first letter. If $w \notin W(a)$, then $a^{-1}w$ is reduced and begins with $a^{-1}$, so $w = a(a^{-1}w) \in aW(a^{-1})$. If
$u \in W(a^{-1})$, $u = a^{-1}u^{\prime}$ with $u^{\prime}$ reduced and not beginning with $a$, so $au = u^{\prime} \notin W(a)$. Hence $aW(a^{-1})$ is exactly the complement of $W(a)$; the same for $b$.
\end{proof}

\begin{remark}\label{measure-theory:remark-banach-tarski}
\emph{The Banach--Tarski paradox} \cite{banach-tarski-1924}: the closed unit ball of $\RR^3$ is the disjoint union of finitely many pieces which can be moved by rigid motions to form two
disjoint copies of the ball. The proof transfers Proposition~\ref{measure-theory:proposition-free-group-paradox} to geometry: the rotation group $\mathrm{SO}(3)$ contains a free group on two generators (for
instance the rotations by $\arccos\frac13$ about two perpendicular axes); it acts freely on the sphere minus a countable set, and choosing, by the axiom of choice, one point in each orbit transports the
paradoxical decomposition of $F$ to the sphere and then, along radii, to the ball, the countable exceptional sets being absorbed by a further trick. The pieces cannot be Lebesgue measurable, since
rigid motions preserve measure and $\lambda(B) \neq 2\lambda(B)$; so the theorem says that, as for the Vitali set (Counterexample~\ref{measure-theory:counterexample-vitali}), no rotation-invariant notion
of volume, even finitely additive, extends to all subsets of $\RR^3$. On the line, Vitali's argument uses countable additivity in an essential way.
\end{remark}

The Vitali set is non-measurable in a mild way: its translates by rationals have the same outer measure, and nothing more is known. A sharper example is a set that is ``half of every interval''
in the strongest sense: it meets every set of positive measure, and so does its complement. It is built by transfinite recursion, and its construction needs one fact about closed sets.

\begin{lemma}[Cantor--Bendixson]\label{measure-theory:lemma-cantor-bendixson}
Every uncountable closed set $F \subseteq \RR$ has cardinality $2^{\aleph_0}$. More precisely, there is an injective map $\{0, 1\}^\NN \to F$.
\end{lemma}

\begin{proof}
\emph{Step 1: condensation points.} Call $x \in F$ a \emph{condensation point} if every open interval containing $x$ meets $F$ in uncountably many points, and let $P$ be the set of condensation points.
A point of $F \setminus P$ lies in an open interval with rational endpoints meeting $F$ in a countable set; there are countably many such intervals, so $F \setminus P$ is countable, and $P$ is
uncountable, in particular nonempty. Moreover, if an open interval $I$ meets $P$, then $I \cap F$ is uncountable, so $I \cap P = (I \cap F) \setminus (F \setminus P)$ is uncountable too.

\emph{Step 2: splitting.} If $J$ is a closed bounded interval whose interior meets $P$, its interior contains two points of $P$, and so there are two disjoint closed subintervals $J_0, J_1 \subseteq J$
of length at most half that of $J$ whose interiors meet $P$.

\emph{Step 3: the tree.} Start with a closed bounded interval $J_\emptyset$ whose interior meets $P$, and apply Step 2 repeatedly: for every finite string $s$ of $0$s and $1$s we obtain $J_s$ with
$J_{s0}, J_{s1} \subseteq J_s$ disjoint, of length at most $2^{-|s|}$ times that of $J_\emptyset$, with interiors meeting $P$. For $\omega \in \{0, 1\}^\NN$ the nested closed intervals
$J_{\omega_1 \cdots \omega_k}$ have lengths tending to $0$; their left endpoints form a nondecreasing bounded sequence, whose limit lies in all of them, so their intersection is a single point $\phi(\omega)$. Each $J_{\omega_1 \cdots \omega_k}$
meets $P \subseteq F$, so it contains points of $F$ at distance at most its length from $\phi(\omega)$, and $\phi(\omega) \in F$ as $F$ is closed. If $\omega \neq \omega^{\prime}$ first differ at place
$k$, $\phi(\omega)$ and $\phi(\omega^{\prime})$ lie in the disjoint intervals $J_{s0}$, $J_{s1}$ with $|s| = k - 1$. So $\phi$ is injective, and
$2^{\aleph_0} = |\{0, 1\}^\NN| \leq |F| \leq |\RR| = 2^{\aleph_0}$ (Theorem~\ref{sets:theorem-schroder-bernstein}).
\end{proof}

\begin{theorem}[Bernstein sets]\label{measure-theory:theorem-bernstein}
Assume the axiom of choice. There is a set $B \subseteq \RR$ such that both $B$ and $\RR \setminus B$ meet every uncountable closed subset of $\RR$. Such a set satisfies
$\lambda^*(B \cap E) = \lambda^*(E \setminus B) = \lambda(E)$ for every Lebesgue measurable $E$; in particular, $B \cap E$ is not measurable whenever $\lambda(E) > 0$.
\end{theorem}

\begin{proof}
\emph{Step 1: counting closed sets.} A closed set is the complement of an open set, and an open set is the union of the intervals with rational endpoints it contains; so there are at most
$|\mathcal{P}(\NN)| = 2^{\aleph_0}$ closed sets, and exactly that many uncountable ones (the intervals $[0, t]$ are distinct). Let $\mathfrak{c} = 2^{\aleph_0}$, regarded as a cardinal, that is, as an
ordinal no smaller ordinal is equinumerous with (Definition~\ref{sets:definition-cardinal}; $\RR$ can be well-ordered by the axiom of choice, Theorem~\ref{sets:theorem-choice-equivalents}). Enumerate the
uncountable closed sets as $(F_\alpha)_{\alpha < \mathfrak{c}}$.

\emph{Step 2: recursion.} By transfinite recursion (Theorem~\ref{sets:theorem-transfinite-recursion}) choose, for $\alpha < \mathfrak{c}$, points $a_\alpha \neq b_\alpha$ of $F_\alpha$ different from all
$a_\beta$, $b_\beta$ with $\beta < \alpha$. This is possible: the set of points already chosen has cardinality at most $|\alpha \times \{0, 1\}|$, which is finite or equal to $|\alpha|$
(Corollary~\ref{sets:corollary-cardinal-arithmetic}), hence less than $\mathfrak{c}$ since $\alpha < \mathfrak{c}$ and $\mathfrak{c}$ is a cardinal; while $|F_\alpha| = \mathfrak{c}$ by
Lemma~\ref{measure-theory:lemma-cantor-bendixson}. (The choices are made with a fixed well-ordering of $\RR$: take the two least admissible points.)

\emph{Step 3: the set.} Let $B = \{a_\alpha : \alpha < \mathfrak{c}\}$. The $a_\alpha$ and $b_\beta$ are all distinct, so every $b_\alpha$ lies in $\RR \setminus B$; thus $B$ meets each $F_\alpha$ in
$a_\alpha$, and $\RR \setminus B$ meets it in $b_\alpha$.

\emph{Step 4: outer measures.} Let $E$ be measurable. If $M \subseteq E \setminus B$ is measurable with $\lambda(M) > 0$, then by inner regularity (Theorem~\ref{measure-theory:theorem-lebesgue-measure}) $M$
contains a compact $K$ with $\lambda(K) > 0$; $K$ is closed and uncountable (countable sets are null), so it meets $B$, contradicting $K \subseteq M \subseteq \RR \setminus B$. So every measurable
subset of $E \setminus B$ is null. Now choose a measurable $G \supseteq B \cap E$ with $\lambda^*(B \cap E) = \lambda(G)$ (outer regularity gives open sets $U_k \supseteq B \cap E$ with
$\lambda(U_k) \leq \lambda^*(B \cap E) + 1/k$; take $G = \bigcap_kU_k$); we may replace $G$ by $G \cap E$. Then $E \setminus G \subseteq E \setminus B$ is measurable, hence null, and
$\lambda^*(B \cap E) = \lambda(G) = \lambda(E)$. The same argument with $B$ and $\RR \setminus B$ exchanged gives $\lambda^*(E \setminus B) = \lambda(E)$. Finally, if $B \cap E$ were measurable with
$\lambda(E) > 0$, take $E$ of finite positive measure (intersect with a bounded interval); then $\lambda(E) = \lambda(B \cap E) + \lambda(E \setminus B) = 2\lambda(E)$, a contradiction.
\end{proof}

\begin{remark}\label{measure-theory:remark-solovay}
The axiom of choice cannot be avoided in Counterexample~\ref{measure-theory:counterexample-vitali} and Theorem~\ref{measure-theory:theorem-bernstein}: Solovay showed that, assuming the consistency of
an inaccessible cardinal, there is a model of set theory with the axiom of dependent choice (enough for all of analysis on separable spaces, including countable additivity arguments) in which every set
of real numbers is Lebesgue measurable \cite{solovay-1970}. On the other hand, the non-measurable sets of the theorem are as wild as possible: by Theorem~\ref{measure-theory:theorem-bernstein}, the
inner measure $\sup\{\lambda(K) : K \subseteq B \text{ compact}\}$ of a Bernstein set is $0$ while its outer measure in every interval is the length of the interval.
\end{remark}

\section{Extension and completion of measures}\label{measure-theory:section-extension}

Lebesgue measure was obtained in two steps: the volume of boxes was given, and Carath\'eodory's theorem extended it to a $\sigma$-algebra through an outer measure. The same pattern occurs whenever a
measure is known on a small family of sets: a probability distribution on $\RR$ is described by the probabilities of the half-lines $(-\infty, x]$, and the distribution of infinitely many coin tosses by
the probabilities of the events depending on finitely many tosses. This section answers the two questions such descriptions raise. When does the data determine the measure (uniqueness)? And when
does a measure with the prescribed values exist (extension)? We also show how to make any measure complete, and construct all locally finite Borel measures on the real line from their distribution
functions.

\subsection{Uniqueness of measures}\label{measure-theory:subsection-uniqueness}

To show that two measures agreeing on a family $\mathcal{P}$ agree on $\sigma(\mathcal{P})$, one considers the sets on which they agree. This family is closed under complements and countable
disjoint unions (when the measures are finite), but not obviously under intersections. The following definitions isolate the two halves of the notion of $\sigma$-algebra.

\begin{definition}\label{measure-theory:definition-dynkin-system}
A family $\mathcal{P}$ of subsets of $X$ is a \emph{$\pi$-system} if it is closed under finite intersections. A family $\mathcal{L}$ of subsets of $X$ is a \emph{Dynkin system} (or
\emph{$\lambda$-system}) if it contains $X$ and is closed under complements and countable disjoint unions.
\end{definition}

\begin{example}\label{measure-theory:example-dynkin-systems}
\begin{enumerate}
\item A $\sigma$-algebra is both a $\pi$-system and a Dynkin system. Conversely, a Dynkin system $\mathcal{L}$ that is a $\pi$-system is a $\sigma$-algebra: it is closed under differences,
$A \setminus B = A \cap (X \setminus B)$, so a countable union $\bigcup_kA_k$ can be written as the disjoint union of the sets $A_k \setminus (A_1 \cup \cdots \cup A_{k-1})
= A_k \cap (X \setminus A_1) \cap \cdots \cap (X \setminus A_{k-1}) \in \mathcal{L}$.
\item The intervals $(a, b]$, $-\infty \leq a < b \leq \infty$ (with $(a, \infty]$ read as $(a, \infty)$), together with $\emptyset$, form a $\pi$-system on $\RR$ generating the Borel $\sigma$-algebra; so do
the half-lines $(-\infty, x]$, $x \in \RR$.
\item On $X = \{1, 2, 3, 4\}$ the subsets with an even number of elements form a Dynkin system which is not a $\sigma$-algebra: complements and disjoint unions of even sets are even, but
$\{1, 2\} \cap \{2, 3\} = \{2\}$.
\item If $\mu$ and $\nu$ are finite measures on $(X, \mathcal{A})$ with $\mu(X) = \nu(X)$, the sets $A \in \mathcal{A}$ with $\mu(A) = \nu(A)$ form a Dynkin system: if $\mu(A) = \nu(A)$, then
$\mu(X \setminus A) = \mu(X) - \mu(A) = \nu(X \setminus A)$, and countable additivity handles disjoint unions. In general it is not a $\pi$-system: for $\mu = \frac12(\delta_1 + \delta_2)$ and
$\nu = \frac12(\delta_3 + \delta_4)$ on $\{1, 2, 3, 4\}$, both $\{1, 3\}$ and $\{2, 3\}$ have measure $\frac12$ under $\mu$ and $\nu$, but their intersection $\{3\}$ does not.
\end{enumerate}
\end{example}

\begin{lemma}[Dynkin]\label{measure-theory:lemma-dynkin}
Let $\mathcal{P}$ be a family of subsets of $X$ closed under finite intersections, and $\mathcal{L} \supseteq \mathcal{P}$ a family containing $X$, closed under complements and countable disjoint unions.
Then $\mathcal{L}$ contains the $\sigma$-algebra generated by $\mathcal{P}$. Consequently two measures that agree on $\mathcal{P}$ and are finite on an increasing sequence of sets of $\mathcal{P}$ exhausting $X$ agree on
$\sigma(\mathcal{P})$.
\end{lemma}

\begin{proof}
Let $\mathcal{L}_0$ be the smallest such family containing $\mathcal{P}$. For $A \in \mathcal{P}$, the sets $B$ with $A \cap B \in \mathcal{L}_0$ form such a family (for complements use
$A \setminus B = X \setminus ((X \setminus A) \sqcup (A \cap B))$) containing $\mathcal{P}$, hence $\mathcal{L}_0$. Repeating with $A \in \mathcal{L}_0$ shows that $\mathcal{L}_0$ is closed under finite
intersections, hence under all countable unions (disjointify), so it is a $\sigma$-algebra. For the consequence, if $\mu = \nu$ on $\mathcal{P}$ and $\mu(X_m) = \nu(X_m) < \infty$, the sets $B$ with
$\mu(B \cap X_m) = \nu(B \cap X_m)$ form such a family; let $m \to \infty$.
\end{proof}

\begin{example}\label{measure-theory:example-distribution-function}
A finite Borel measure $\mu$ on $\RR$ is determined by its \emph{distribution function} $F_\mu(x) = \mu((-\infty, x])$. Indeed, two finite measures with the same distribution function agree on the
$\pi$-system of half-lines, and on $\RR = \bigcup_m(-\infty, m]$ (finite measure), so they agree on the $\sigma$-algebra generated, which is $\mathcal{B}(\RR)$ by
Example~\ref{measure-theory:example-dynkin-systems}(2). Likewise, a Borel measure on $\RR^n$ that is finite on bounded sets is determined by its values on boxes $\prod_i(a_i, b_i]$; for instance it
is Lebesgue measure if it gives every such box its volume.
\end{example}

\begin{counterexample}\label{measure-theory:counterexample-uniqueness}
The finiteness hypothesis in Lemma~\ref{measure-theory:lemma-dynkin} cannot be dropped. On the Borel sets of $\RR$ let $\mu(E)$ be the number of rational points in $E$ (possibly $\infty$), and
$\nu = 2\mu$. Both are measures, and they agree on the $\pi$-system of intervals $(a, b]$, all of which have measure $\infty$, and on $\emptyset$. But $\mu(\{0\}) = 1 \neq 2 = \nu(\{0\})$.
\end{counterexample}

\subsection{Premeasures and the extension theorem}\label{measure-theory:subsection-extension}

\begin{definition}\label{measure-theory:definition-premeasure}
An \emph{algebra} of subsets of $X$ is a family $\mathcal{A}_0$ containing $X$ and closed under complements and finite unions (hence under finite intersections and differences). A
\emph{premeasure} on $\mathcal{A}_0$ is a map $\mu_0 \colon \mathcal{A}_0 \to [0, \infty]$ with $\mu_0(\emptyset) = 0$ such that $\mu_0(\bigsqcup_kA_k) = \sum_k\mu_0(A_k)$ for every sequence of
pairwise disjoint $A_k \in \mathcal{A}_0$ \emph{whose union lies in $\mathcal{A}_0$}. It is \emph{$\sigma$-finite} if $X$ is a countable union of sets $A \in \mathcal{A}_0$ with $\mu_0(A) < \infty$.
The \emph{outer measure induced} by $\mu_0$ is
\[
\mu^*(E) = \inf\Bigl\{\sum_k\mu_0(A_k) : A_k \in \mathcal{A}_0, \ E \subseteq \bigcup_kA_k\Bigr\} \qquad (E \subseteq X).
\]
\end{definition}

A premeasure is finitely additive and monotone on $\mathcal{A}_0$ ($\mu_0(B) = \mu_0(A) + \mu_0(B \setminus A)$ for $A \subseteq B$). The condition that the union lies in $\mathcal{A}_0$ is
essential, since an algebra need not contain countable unions: the finite and cofinite subsets of $\NN$ form an algebra containing all singletons but not the set of even numbers.

\begin{theorem}[Carath\'eodory--Hahn extension theorem]\label{measure-theory:theorem-extension}
Let $\mu_0$ be a premeasure on an algebra $\mathcal{A}_0$ and $\mu^*$ the induced outer measure. Then:
\begin{enumerate}
\item $\mu^*$ is an outer measure, and $\mu^*(A) = \mu_0(A)$ for $A \in \mathcal{A}_0$;
\item every $A \in \mathcal{A}_0$ is $\mu^*$-measurable; hence $\mu = \mu^*|_{\sigma(\mathcal{A}_0)}$ is a measure on $\sigma(\mathcal{A}_0)$ extending $\mu_0$;
\item if $\mu_0$ is $\sigma$-finite, $\mu$ is the only measure on $\sigma(\mathcal{A}_0)$ extending $\mu_0$.
\end{enumerate}
\end{theorem}

\begin{proof}
(1) $\mu^*(\emptyset) = 0$ and monotonicity are clear. Countable subadditivity: given $E_1, E_2, \ldots$ and $\varepsilon > 0$, cover each $E_j$ by sets of $\mathcal{A}_0$ with
$\sum_k\mu_0(A_{jk}) \leq \mu^*(E_j) + \varepsilon2^{-j}$; all the $A_{jk}$ together cover $\bigcup_jE_j$, so $\mu^*(\bigcup_jE_j) \leq \sum_j\mu^*(E_j) + \varepsilon$. For $A \in \mathcal{A}_0$,
$\mu^*(A) \leq \mu_0(A)$ by the cover $\{A\}$. Conversely, if $A \subseteq \bigcup_kA_k$ with $A_k \in \mathcal{A}_0$, the sets $B_k = A \cap A_k \setminus (A_1 \cup \cdots \cup A_{k-1})$ lie in
$\mathcal{A}_0$, are disjoint and have union $A \in \mathcal{A}_0$; so $\mu_0(A) = \sum_k\mu_0(B_k) \leq \sum_k\mu_0(A_k)$ by the premeasure property and monotonicity. Hence
$\mu_0(A) \leq \mu^*(A)$.

(2) Let $A \in \mathcal{A}_0$, $E \subseteq X$ and $\varepsilon > 0$, and choose $A_k \in \mathcal{A}_0$ covering $E$ with $\sum_k\mu_0(A_k) \leq \mu^*(E) + \varepsilon$. The sets $A_k \cap A$ cover
$E \cap A$ and the sets $A_k \setminus A$ cover $E \setminus A$, so by additivity of $\mu_0$,
\[\begin{aligned}
\mu^*(E \cap A) + \mu^*(E \setminus A) &\leq \sum_k\bigl(\mu_0(A_k \cap A) + \mu_0(A_k \setminus A)\bigr) \\
&= \sum_k\mu_0(A_k) \leq \mu^*(E) + \varepsilon .
\end{aligned}\]
So $A$ is $\mu^*$-measurable. By Carath\'eodory's theorem (Theorem~\ref{measure-theory:theorem-caratheodory}) the $\mu^*$-measurable sets form a $\sigma$-algebra on which $\mu^*$ is a measure; it
contains $\mathcal{A}_0$, hence $\sigma(\mathcal{A}_0)$, and $\mu$ extends $\mu_0$ by (1).

(3) An algebra is a $\pi$-system. If $X = \bigcup_kC_k$ with $C_k \in \mathcal{A}_0$ and $\mu_0(C_k) < \infty$, the sets $X_m = C_1 \cup \cdots \cup C_m$ lie in $\mathcal{A}_0$, increase to $X$
and have finite measure. Two measures extending $\mu_0$ agree on $\mathcal{A}_0$, so they agree on $\sigma(\mathcal{A}_0)$ by Lemma~\ref{measure-theory:lemma-dynkin}.
\end{proof}

\begin{remark}\label{measure-theory:remark-extension}
\begin{enumerate}
\item The construction of Lebesgue measure in Theorem~\ref{measure-theory:theorem-lebesgue-measure} is an instance of this pattern: its step ``volume of a box'' is the statement that volume is
countably subadditive on boxes, which is what makes $\mu^*$ agree with the given values; here this is the premeasure property.
\item The measure $\mu^*$ on \emph{all} $\mu^*$-measurable sets is complete (Theorem~\ref{measure-theory:theorem-caratheodory}), and in general larger than $\mu$ on $\sigma(\mathcal{A}_0)$. For
Lebesgue measure it is exactly the completion of $\mu$ in the sense of the next subsection (Proposition~\ref{measure-theory:proposition-lebesgue-completion}), and the same holds for
every $\sigma$-finite $\mu_0$ \cite[Section 1.4]{folland-real}.
\item Without $\sigma$-finiteness the extension need not be unique. On $X = \QQ$ let $\mathcal{A}_0$ be the algebra of finite unions of sets $(a, b] \cap \QQ$ with $a < b$ real (and their complements, which
are again such unions with $a = -\infty$ or $b = \infty$ allowed), and $\mu_0(A) = \infty$ for $A \neq \emptyset$. This is a premeasure, $\sigma(\mathcal{A}_0)$ is the power set (every singleton is a
countable intersection of such sets), and both counting measure and twice counting measure extend $\mu_0$.
\item A central application is the construction of measures on infinite products, such as the probability measure describing infinitely many independent coin tosses on $\{0, 1\}^\NN$: the premeasure is
defined on the algebra of sets depending on finitely many coordinates, and its countable additivity follows from compactness of $\{0, 1\}^\NN$, or in general from Fubini's theorem; this is done in
Subsection~\ref{measure-theory:subsection-infinite-products}.
\end{enumerate}
\end{remark}

\subsection{Completion}\label{measure-theory:subsection-completion}

Lebesgue measure is complete, but its restriction to the Borel sets is not (Remark~\ref{measure-theory:remark-lebesgue-not-borel} below). Completeness is convenient: properties holding almost
everywhere then concern measurable sets, and modifying a measurable function on a null set keeps it measurable. Every measure can be completed, in a unique way.

\begin{theorem}\label{measure-theory:theorem-completion}
Let $(X, \mathcal{A}, \mu)$ be a measure space, and let $\overline{\mathcal{A}}$ be the family of sets $E \cup F$ with $E \in \mathcal{A}$ and $F \subseteq N$ for some $N \in \mathcal{A}$ with
$\mu(N) = 0$. Then $\overline{\mathcal{A}}$ is a $\sigma$-algebra containing $\mathcal{A}$, and $\bar\mu(E \cup F) = \mu(E)$ defines the unique measure on $\overline{\mathcal{A}}$ extending $\mu$. The
measure $\bar\mu$ is complete; it is the \emph{completion} of $\mu$.
\end{theorem}

\begin{proof}
\emph{$\sigma$-algebra.} Countable unions: $\bigcup_k(E_k \cup F_k) = (\bigcup_kE_k) \cup (\bigcup_kF_k)$ and $\bigcup_kF_k \subseteq \bigcup_kN_k$, a null set. Complements: given $E \cup F$ with
$F \subseteq N$, we may assume $E \cap N = \emptyset$, replacing $F$ by $F \setminus E$ and $N$ by $N \setminus E$. Then
$X \setminus (E \cup F) = (X \setminus (E \cup N)) \cup (N \setminus F)$, with $X \setminus (E \cup N) \in \mathcal{A}$ and $N \setminus F \subseteq N$.

\emph{$\bar\mu$ is well defined.} If $E_1 \cup F_1 = E_2 \cup F_2$ with $F_j \subseteq N_j$, then $E_1 \subseteq E_2 \cup N_2$, so $\mu(E_1) \leq \mu(E_2) + \mu(N_2) = \mu(E_2)$, and symmetrically.

\emph{Measure, completeness, uniqueness.} If the sets $E_k \cup F_k$ are disjoint, so are the $E_k$, and $\bar\mu(\bigsqcup_k(E_k \cup F_k)) = \mu(\bigsqcup_kE_k) = \sum_k\mu(E_k)$. If
$\bar\mu(E \cup F) = 0$, then $E \cup F \subseteq E \cup N$, a $\mu$-null set of $\mathcal{A}$, so every subset of $E \cup F$ lies in $\overline{\mathcal{A}}$. If $\nu$ is any measure on
$\overline{\mathcal{A}}$ extending $\mu$, then $\mu(E) = \nu(E) \leq \nu(E \cup F) \leq \nu(E \cup N) = \mu(E \cup N) = \mu(E)$.
\end{proof}

\begin{example}\label{measure-theory:example-completion}
\begin{enumerate}
\item The completion of the Dirac measure $\delta_0$ on $\mathcal{B}(\RR)$ is defined on all subsets of $\RR$: every $A \subseteq \RR$ is $(A \cap \{0\}) \cup (A \setminus \{0\})$, and
$A \setminus \{0\}$ lies in the null set $\RR \setminus \{0\}$.
\item A complete measure stays the same: $\overline{\mathcal{A}} = \mathcal{A}$.
\item For the completion, a function $g$ that agrees $\mu$-almost everywhere with an $\mathcal{A}$-measurable $f$ is $\overline{\mathcal{A}}$-measurable: $\{g > a\}$ differs from $\{f > a\}$ by a subset of
a null set.
\end{enumerate}
\end{example}

\begin{proposition}\label{measure-theory:proposition-lebesgue-completion}
A set $E \subseteq \RR^n$ is Lebesgue measurable if and only if there are a countable union $A$ of compact sets and a countable intersection $B$ of open sets with $A \subseteq E \subseteq B$ and
$\lambda(B \setminus A) = 0$. Consequently the Lebesgue $\sigma$-algebra is the completion of $\mathcal{B}(\RR^n)$ with respect to $\lambda$, and $\lambda$ is the completion of its restriction to
$\mathcal{B}(\RR^n)$.
\end{proposition}

\begin{proof}
Let $E$ be measurable and $E_m = E \cap [-m, m]^n$, of finite measure. By the regularity in Theorem~\ref{measure-theory:theorem-lebesgue-measure}, for all $m, j \geq 1$ there are a compact
$K_{m,j} \subseteq E_m$ with $\lambda(E_m \setminus K_{m,j}) < 1/j$ and an open $U_{m,j} \supseteq E_m$ with $\lambda(U_{m,j} \setminus E_m) < 2^{-m}/j$. Let $A = \bigcup_{m,j}K_{m,j}$ and
$B = \bigcap_j\bigcup_mU_{m,j}$. Then $A \subseteq E \subseteq B$. For each $m$, $\lambda(E_m \setminus A) < 1/j$ for all $j$, so $\lambda(E \setminus A) \leq \sum_m\lambda(E_m \setminus A) = 0$; and
$\bigcup_mU_{m,j} \setminus E \subseteq \bigcup_m(U_{m,j} \setminus E_m)$ has measure $< 1/j$, so $\lambda(B \setminus E) = 0$.

Conversely, if such $A$ and $B$ exist, then $E = A \cup (E \setminus A)$ with $A$ Borel and $E \setminus A$ contained in the Borel null set $B \setminus A$; so $E$ lies in the completion of
$\mathcal{B}(\RR^n)$, and it is Lebesgue measurable because $\lambda$ is complete (Theorem~\ref{measure-theory:theorem-caratheodory}). This proves both statements, the measures agreeing on Borel sets.
\end{proof}

\begin{remark}\label{measure-theory:remark-lebesgue-not-borel}
The completion is strictly larger: there are Lebesgue measurable sets that are not Borel. Such a set can be found inside the Cantor set $C$, which has measure $0$
(Exercise~\ref{measure-theory:exercise-cantor}), so that all its subsets are Lebesgue measurable by completeness; an explicit non-Borel subset is constructed in Chapter~\ref{lebesgue-differentiation}
with the help of the Cantor function.
\end{remark}

\subsection{Lebesgue--Stieltjes measures on the line}\label{measure-theory:subsection-stieltjes}

By Example~\ref{measure-theory:example-distribution-function}, a finite Borel measure on $\RR$ is determined by its distribution function, which is nondecreasing. We now show that every nondecreasing
right-continuous function arises, and allow infinite measures as long as bounded intervals have finite measure. The natural building blocks are the half-open intervals: they fit together without
overlaps or gaps, $(a, b] \sqcup (b, c] = (a, c]$.

\begin{lemma}\label{measure-theory:lemma-half-open-algebra}
Call \emph{h-intervals} the sets $(a, b]$ with $-\infty \leq a < b < \infty$, the sets $(a, \infty)$ with $-\infty \leq a < \infty$, and $\emptyset$. The finite disjoint unions of h-intervals form an
algebra $\mathcal{A}_h$ on $\RR$, and $\sigma(\mathcal{A}_h) = \mathcal{B}(\RR)$.
\end{lemma}

\begin{proof}
The intersection of two h-intervals is an h-interval, and the complement of an h-interval is a disjoint union of at most two h-intervals (for instance $\RR \setminus (a, b] = (-\infty, a] \sqcup (b, \infty)$).
If $A = \bigsqcup_iI_i$ and $B = \bigsqcup_jJ_j$ are in $\mathcal{A}_h$, then $A \cap B = \bigsqcup_{i,j}I_i \cap J_j \in \mathcal{A}_h$, and $\RR \setminus A = \bigcap_i(\RR \setminus I_i)$ is a finite
intersection of members of $\mathcal{A}_h$, hence in $\mathcal{A}_h$. Closure under intersections and complements gives closure under unions. Every h-interval is Borel, and the open intervals are
countable unions of h-intervals, $(a, b) = \bigcup_k(a, b - 1/k]$, so $\sigma(\mathcal{A}_h)$ contains the open sets.
\end{proof}

Let $F \colon \RR \to \RR$ be nondecreasing and right-continuous, and write $F(\infty) = \lim_{x \to \infty}F(x)$ and $F(-\infty) = \lim_{x \to -\infty}F(x)$, in $[-\infty, \infty]$. Define
$\mu_0((a, b]) = F(b) - F(a)$ and $\mu_0((a, \infty)) = F(\infty) - F(a)$ for h-intervals (with $F(a) = F(-\infty)$ if $a = -\infty$, and the value $\infty$ if a difference involves an infinite
limit), $\mu_0(\emptyset) = 0$, and $\mu_0(\bigsqcup_iI_i) = \sum_i\mu_0(I_i)$ for finite disjoint unions. This is well defined: if an h-interval $I$ is the disjoint union of h-intervals
$I_1, \ldots, I_m$, ordering them from left to right shows that the right endpoint of each is the left endpoint of the next, so $\sum_i\mu_0(I_i)$ telescopes to $\mu_0(I)$; and two
representations of the same set have a common refinement by the intersections $I_i \cap J_j$.

\begin{lemma}\label{measure-theory:lemma-stieltjes-premeasure}
$\mu_0$ is a $\sigma$-finite premeasure on $\mathcal{A}_h$.
\end{lemma}

\begin{proof}
$\sigma$-finiteness: $\RR = \bigcup_m(-m, m]$ with $\mu_0((-m, m]) = F(m) - F(-m) < \infty$. Since every element of $\mathcal{A}_h$ is a finite disjoint union of h-intervals, and each piece of a
countable disjoint decomposition can be intersected with these finitely many h-intervals, it suffices to show: if an h-interval $I$ is the disjoint union of countably many h-intervals $I_k$, then
$\mu_0(I) = \sum_k\mu_0(I_k)$.

\emph{Step 1: $\geq$.} For each $n$, $I \setminus (I_1 \cup \cdots \cup I_n) \in \mathcal{A}_h$, and finite additivity gives $\mu_0(I) \geq \sum_{k \leq n}\mu_0(I_k)$.

\emph{Step 2: a finite covering inequality.} If a compact interval $[c, d]$ is covered by finitely many open intervals $(\alpha_1, \beta_1), \ldots, (\alpha_N, \beta_N)$, then
$F(d) - F(c) \leq \sum_i(F(\beta_i) - F(\alpha_i))$. Indeed, choose an interval of the cover containing $c$, call it $(\alpha_{i_1}, \beta_{i_1})$; if $\beta_{i_1} \leq d$, choose one containing
$\beta_{i_1}$, call it $(\alpha_{i_2}, \beta_{i_2})$; and so on. The right endpoints increase strictly, so no interval is chosen twice and the process stops after $r \leq N$ steps with $\beta_{i_r} > d$.
Since $\alpha_{i_1} < c$ and $\alpha_{i_{s+1}} < \beta_{i_s}$, and $F$ is nondecreasing,
\[
F(d) - F(c) \leq F(\beta_{i_r}) - F(\alpha_{i_1}) = \sum_{s=1}^r\bigl(F(\beta_{i_s}) - F(\beta_{i_{s-1}})\bigr) \leq \sum_{s=1}^r\bigl(F(\beta_{i_s}) - F(\alpha_{i_s})\bigr),
\]
where $\beta_{i_0} = \alpha_{i_1}$.

\emph{Step 3: $\leq$ for bounded $I = (a, b]$.} Let $I_k = (a_k, b_k]$ and $\varepsilon > 0$. By right continuity choose $0 < \delta < b - a$ with $F(a + \delta) < F(a) + \varepsilon$, and $\delta_k > 0$ with
$F(b_k + \delta_k) < F(b_k) + \varepsilon2^{-k}$. The open intervals $(a_k, b_k + \delta_k)$ cover the compact interval $[a + \delta, b]$, so finitely many of them do, and Step 2 gives
\[
F(b) - F(a + \delta) \leq \sum_k\bigl(F(b_k + \delta_k) - F(a_k)\bigr) \leq \sum_k\mu_0(I_k) + \varepsilon .
\]
Hence $\mu_0(I) = F(b) - F(a) \leq \sum_k\mu_0(I_k) + 2\varepsilon$.

\emph{Step 4: unbounded $I$.} For $M > 0$, $I \cap (-M, M]$ is a bounded h-interval (or empty), the disjoint union of the h-intervals $I_k \cap (-M, M]$, so by Step 3
$\mu_0(I \cap (-M, M]) \leq \sum_k\mu_0(I_k \cap (-M, M]) \leq \sum_k\mu_0(I_k)$. As $M \to \infty$ the left side tends to $\mu_0(I)$, by the definition of $F(\pm\infty)$.
\end{proof}

\begin{theorem}[Lebesgue--Stieltjes measures]\label{measure-theory:theorem-lebesgue-stieltjes}
\begin{enumerate}
\item For every nondecreasing right-continuous $F \colon \RR \to \RR$ there is a unique Borel measure $\mu_F$ on $\RR$ with $\mu_F((a, b]) = F(b) - F(a)$ for all $a < b$.
\item $\mu_F = \mu_G$ if and only if $F - G$ is constant.
\item Conversely, let $\mu$ be a Borel measure on $\RR$ that is finite on bounded intervals, and put $F(x) = \mu((0, x])$ for $x \geq 0$ and $F(x) = -\mu((x, 0])$ for $x < 0$. Then $F$ is nondecreasing and
right-continuous, and $\mu = \mu_F$.
\item For all $x$ and $a < b$, $\mu_F(\{x\}) = F(x) - F(x-)$, $\mu_F([a, b]) = F(b) - F(a-)$ and $\mu_F((a, b)) = F(b-) - F(a)$, where $F(x-) = \lim_{t \uparrow x}F(t)$. So $\mu_F$ has an atom
exactly at each point where $F$ jumps, of mass equal to the jump.
\end{enumerate}
\end{theorem}

\begin{proof}
(1) Existence follows from Theorem~\ref{measure-theory:theorem-extension}, Lemma~\ref{measure-theory:lemma-stieltjes-premeasure} and $\sigma(\mathcal{A}_h) = \mathcal{B}(\RR)$. A Borel measure with the
given values on bounded h-intervals has the values of $\mu_0$ on all h-intervals (by continuity from below, $(a, \infty) = \bigcup_m(a, m]$, and similarly for $a = -\infty$), hence on
$\mathcal{A}_h$; uniqueness follows from Theorem~\ref{measure-theory:theorem-extension}(3).

(2) If $\mu_F = \mu_G$, then $F(b) - F(a) = G(b) - G(a)$ for all $a < b$. Conversely, $F$ and $F + c$ define the same $\mu_0$.

(3) In every case $\mu((x, y]) = F(y) - F(x)$ for $x < y$ (for $x < 0 \leq y$, $(x, y] = (x, 0] \sqcup (0, y]$), so $F$ is nondecreasing. If $x_k \downarrow x$, the sets $(x, x_k]$ decrease to
$\emptyset$ and have finite measure, so $F(x_k) - F(x) = \mu((x, x_k]) \to 0$ (Proposition~\ref{measure-theory:proposition-measure-properties}): $F$ is right-continuous. Then $\mu$ and $\mu_F$ agree on
bounded h-intervals, and $\mu = \mu_F$ by (1).

(4) $\{x\} = \bigcap_k(x - 1/k, x]$, a decreasing sequence of sets of finite measure, so $\mu_F(\{x\}) = \lim_k(F(x) - F(x - 1/k)) = F(x) - F(x-)$. Then
$\mu_F([a, b]) = \mu_F(\{a\}) + \mu_F((a, b])$ and $\mu_F((a, b)) = \mu_F((a, b]) - \mu_F(\{b\})$.
\end{proof}

\begin{definition}\label{measure-theory:definition-lebesgue-stieltjes}
The measure $\mu_F$ is the \emph{Lebesgue--Stieltjes measure} of $F$; its completion is also called by this name. For a finite Borel measure $\mu$ on $\RR$ one usually normalises $F$ by
$F(-\infty) = 0$, so that $F(x) = \mu((-\infty, x])$ is the distribution function of Example~\ref{measure-theory:example-distribution-function}.
\end{definition}

\begin{example}\label{measure-theory:example-lebesgue-stieltjes}
\begin{enumerate}
\item $F(x) = x$ gives Lebesgue measure: $\lambda((a, b]) = b - a$ by Theorem~\ref{measure-theory:theorem-lebesgue-measure} (the point $a$ has measure $\lambda^*([a, a]) = 0$), so $\lambda = \mu_F$ on
$\mathcal{B}(\RR)$ by uniqueness. This is a second construction of one-dimensional Lebesgue measure.
\item $F = \mathbf{1}_{[0, \infty)}$ gives $\delta_0$, and $F(x) = \lfloor x\rfloor$ gives counting measure on $\ZZ$, whose atoms are the jumps of the floor function.
\item \emph{Probability distributions.} The distribution functions of Borel probability measures on $\RR$ are exactly the nondecreasing right-continuous $F$ with $F(-\infty) = 0$ and $F(\infty) = 1$.
For instance $F(x) = (1 - e^{-x})^+$ is the exponential distribution, $\mu_F = e^{-x}\mathbf{1}_{[0, \infty)}\lambda$ (the two measures agree on h-intervals by the fundamental theorem of calculus and
Proposition~\ref{measure-theory:proposition-riemann-lebesgue}), and a fair die has the step function $F(x) = \frac16\#\{k \in \{1, \ldots, 6\} : k \leq x\}$.
\item \emph{A dense set of jumps.} Let $q_1, q_2, \ldots$ enumerate $\QQ$ and $F(x) = \sum_{q_n \leq x}2^{-n}$. Then $F$ is nondecreasing, and right-continuous: $F(x + h) - F(x)$ is the sum of $2^{-n}$
over the $n$ with $x < q_n \leq x + h$; given $N$, for small $h$ no $q_n$ with $n \leq N$ lies in $(x, x + h]$, so the sum is at most $2^{-N}$. By (4) of the theorem, $F$ jumps by $2^{-n}$ at $q_n$
and is continuous at every irrational point, and $\mu_F = \sum_n2^{-n}\delta_{q_n}$, a probability measure carried by the countable set $\QQ$.
\item \emph{A continuous singular distribution.} The Cantor function of Chapter~\ref{lebesgue-differentiation}, extended by $0$ on $(-\infty, 0)$ and $1$ on $(1, \infty)$, is continuous, so $\mu_F$ has no
atoms; yet $\mu_F$ is carried by the Cantor set, which has Lebesgue measure $0$.
\end{enumerate}
\end{example}

\begin{counterexample}\label{measure-theory:counterexample-right-continuity}
Right continuity cannot be dropped. For $F = \mathbf{1}_{(0, \infty)}$, which is nondecreasing but not right-continuous at $0$, the formula $\mu_0((a, b]) = F(b) - F(a)$ is still finitely additive on
$\mathcal{A}_h$, but $(0, 1] = \bigsqcup_k(\frac1{k+1}, \frac1k]$ has $\mu_0((0, 1]) = 1$ while every piece has $\mu_0 = 0$. So $\mu_0$ is not a premeasure, and no measure $\mu$ satisfies
$\mu((a, b]) = F(b) - F(a)$: it would have $\mu((0, 1/k]) = 1$ for all $k$, contradicting continuity from above. The convention that distribution functions are right-continuous corresponds to the choice
of half-open intervals $(a, b]$; left-continuous functions go with $[a, b)$.
\end{counterexample}

\begin{proposition}\label{measure-theory:proposition-stieltjes-regular}
Every Lebesgue--Stieltjes measure $\mu_F$ is regular: for every Borel set $E$,
\[
\mu_F(E) = \inf\{\mu_F(U) : U \supseteq E \text{ open}\} = \sup\{\mu_F(K) : K \subseteq E \text{ compact}\}.
\]
\end{proposition}

\begin{proof}
\emph{Outer regularity.} By Theorem~\ref{measure-theory:theorem-extension}, $\mu_F(E) = \mu^*(E)$ is the infimum of $\sum_k\mu_0(I_k)$ over covers of $E$ by h-intervals (members of
$\mathcal{A}_h$ being finite unions of h-intervals); splitting $(a, \infty)$ into $(a, a + 1], (a + 1, a + 2], \ldots$ we may take all $I_k = (a_k, b_k]$ with $b_k < \infty$. Given $\varepsilon > 0$ and
such a cover with $\sum_k\mu_0(I_k) \leq \mu_F(E) + \varepsilon$, choose $\delta_k > 0$ with $F(b_k + \delta_k) < F(b_k) + \varepsilon2^{-k}$. Then $U = \bigcup_k(a_k, b_k + \delta_k)$ is open, contains $E$, and
$\mu_F(U) \leq \sum_k(F(b_k + \delta_k) - F(a_k)) \leq \mu_F(E) + 2\varepsilon$.

\emph{Inner regularity.} If $E$ is bounded, contained in a compact interval $J$, choose by outer regularity an open $U \supseteq J \setminus E$ with $\mu_F(U) \leq \mu_F(J \setminus E) + \varepsilon$; then
$K = J \setminus U$ is compact, contained in $E$, and $E \setminus K \subseteq U \setminus (J \setminus E)$ has measure $\leq \varepsilon$ (all measures here are finite). For general $E$ apply this to
$E \cap [-m, m]$ and use continuity from below.
\end{proof}

The same argument, with boxes in place of h-intervals, would construct Lebesgue measure on $\RR^n$ once more; distribution functions in several variables exist as well, but nondecreasing in
each variable is no longer enough: one needs that every box $\prod_i(a_i, b_i]$ receives a nonnegative value from the alternating sum of $F$ over its $2^n$ vertices.

\section{The integral}\label{measure-theory:section-integral}

\begin{definition}\label{measure-theory:definition-integral}
For a simple function $s = \sum_ia_i\mathbf{1}_{A_i} \geq 0$ with disjoint $A_i$, $\int s\,d\mu = \sum_ia_i\mu(A_i)$ (independent of the representation, by passing to a common refinement). For measurable
$f \colon X \to [0, \infty]$, $\int f\,d\mu = \sup\{\int s\,d\mu : 0 \leq s \leq f \text{ simple}\}$, and $\int_Ef = \int f\mathbf{1}_E$. A measurable $f \colon X \to \CC$ is \emph{integrable} if $\int|f| < \infty$; then
$\int f = \int(\operatorname{Re}f)^+ - \int(\operatorname{Re}f)^- + i(\int(\operatorname{Im}f)^+ - \int(\operatorname{Im}f)^-)$, where $g^\pm = \max(\pm g, 0)$.
\end{definition}

\emph{The idea of the definition.} The Riemann integral cuts the \emph{domain} into intervals; the Lebesgue integral cuts the \emph{range}: for $f \geq 0$ and levels $0 = y_0 < y_1 < \cdots$, the simple
function $\sum_iy_i\mathbf{1}_{\{y_i \leq f < y_{i+1}\}}$ lies below $f$, its integral is $\sum_iy_i\,\mu(\{y_i \leq f < y_{i+1}\})$, and refining the levels approximates $\int f$ from below. This only needs the
measures of the level sets, which is why it works on any measure space and for very irregular functions.

\begin{example}\label{measure-theory:example-integrals}
\begin{enumerate}
\item \emph{Counting measure on $\NN$.} For $f \colon \NN \to [0, \infty]$, $\int f\,d\mu = \sum_nf(n)$. If some $f(n) = \infty$ both sides are $\infty$. Otherwise the functions $s_N = f\mathbf{1}_{\{0, \ldots, N\}}$ are simple and satisfy $\int s_N = \sum_{n \leq N}f(n)$ and increase to $f$, so the monotone convergence theorem below gives the claim; a function $f \colon \NN \to \CC$ is integrable
exactly when $\sum_n|f(n)| < \infty$. So series are integrals, and theorems about integrals specialise to theorems about series.
\item \emph{Dirac measure.} $\int f\,d\delta_x = f(x)$: a simple function $s \leq f$ has $\int s\,d\delta_x = s(x) \leq f(x)$, and the simple function $f(x)\mathbf{1}_{\{x\}}$ attains it.
\item \emph{The indicator of the rationals.} $\int_{[0, 1]}\mathbf{1}_\QQ\,d\lambda = \lambda(\QQ \cap [0, 1]) = 0$, as countable sets have Lebesgue measure $0$. The Riemann integral of $\mathbf{1}_\QQ$ does
not exist: every upper sum is $1$ and every lower sum is $0$.
\end{enumerate}
\end{example}

\begin{theorem}[Monotone convergence]\label{measure-theory:theorem-monotone-convergence}
If $0 \leq f_1 \leq f_2 \leq \cdots$ are measurable and $f = \lim f_n$, then $\int f = \lim\int f_n$.
\end{theorem}

\begin{proof}
The integrals increase and are at most $\int f$. Let $0 \leq s \leq f$ be simple and $0 < c < 1$, and $E_n = \{f_n \geq cs\}$, which increase to $X$. Then $\int f_n \geq c\int_{E_n}s$, and
$A \mapsto \int_As$ is a measure, so $\int_{E_n}s \to \int s$ (Proposition~\ref{measure-theory:proposition-measure-properties}). Hence $\lim\int f_n \geq c\int s$ for all such $c$ and $s$.
\end{proof}

\begin{corollary}\label{measure-theory:corollary-integral-properties}
\begin{enumerate}
\item For measurable $f, g \geq 0$ and $a \geq 0$, $\int(af + g) = a\int f + \int g$, and $\int\sum_nf_n = \sum_n\int f_n$.
\item (Fatou) For measurable $f_n \geq 0$, $\int\liminf f_n \leq \liminf\int f_n$.
\item The integral is linear on integrable functions and $|\int f| \leq \int|f|$.
\item For measurable $f \geq 0$, $\int f = 0$ if and only if $f = 0$ a.e.; and $\mu(\{f \geq t\}) \leq \frac1t\int f$ for $t > 0$ (Markov's inequality).
\end{enumerate}
\end{corollary}

\begin{proof}
(1) For simple functions additivity is clear; approximate by increasing simple functions (Proposition~\ref{measure-theory:proposition-measurable-functions}(4)) and use monotone convergence. (2) Apply
monotone convergence to $g_n = \inf_{k \geq n}f_k$, which satisfy $\int g_n \leq \inf_{k \geq n}\int f_k$. (3) Linearity follows from (1) by decomposing into positive and negative parts. For complex $f$, write
$\int f = |\int f|e^{i\theta}$; then $|\int f| = \int\operatorname{Re}(e^{-i\theta}f) \leq \int|f|$. (4) $t\mathbf{1}_{\{f \geq t\}} \leq f$; and $\{f > 0\} = \bigcup_n\{f \geq 1/n\}$.
\end{proof}

\begin{remark}\label{measure-theory:remark-markov-chebyshev}
\emph{Names.} The inequality $\mu(\{f \geq t\}) \leq \frac1t\int f$ is usually called \emph{Markov's inequality}, and some books call it Chebyshev's. \emph{Chebyshev's inequality} proper is the
special case applied to $(f - m)^2$: for a probability measure, $f \in L^2$ with mean $m = \int f$ and variance $\sigma^2 = \int(f - m)^2$,
\[
\mu(\{|f - m| \geq t\}) = \mu(\{(f - m)^2 \geq t^2\}) \leq \frac{\sigma^2}{t^2}.
\]
More generally $\mu(\{|f| \geq t\}) \leq t^{-p}\int|f|^p$ for $p > 0$. The proof is one line, $t\mathbf{1}_{\{f \geq t\}} \leq f$, and it is often all one needs: for instance it shows that
convergence in $L^p$ implies convergence in measure (Proposition~\ref{measure-theory:proposition-modes}).
\end{remark}

\begin{proposition}[Jensen's inequality]\label{measure-theory:proposition-jensen}
Let $\mu$ be a probability measure, $I \subseteq \RR$ an open interval, $\varphi \colon I \to \RR$ convex, and $f$ an integrable function with values in $I$. Then $\int f \in I$, and
\[
\varphi\Bigl(\int f\,d\mu\Bigr) \leq \int\varphi \circ f\,d\mu,
\]
where the right side is defined in $(-\infty, \infty]$.
\end{proposition}

\begin{proof}
\emph{Convex functions.} For $s < t < u$ in $I$, writing $t = \lambda s + (1 - \lambda)u$, convexity $\varphi(t) \leq \lambda\varphi(s) + (1 - \lambda)\varphi(u)$ is equivalent to each of the inequalities
between slopes
\[
\frac{\varphi(t) - \varphi(s)}{t - s} \leq \frac{\varphi(u) - \varphi(s)}{u - s} \leq \frac{\varphi(u) - \varphi(t)}{u - t}.
\]
So for $c \in I$ the slope $(\varphi(t) - \varphi(c))/(t - c)$ is nondecreasing in $t \neq c$; with $m = \sup_{t < c}(\varphi(t) - \varphi(c))/(t - c) \leq \inf_{t > c}(\varphi(t) - \varphi(c))/(t - c)$,
a finite number, $\varphi(t) \geq \varphi(c) + m(t - c)$ for all $t \in I$: the graph lies above a line through $(c, \varphi(c))$. The same monotonicity shows that $\varphi$ has finite one-sided
derivatives, so it is continuous, and $\varphi \circ f$ is measurable.

\emph{Jensen.} Let $c = \int f$. If $I$ is bounded above by $b$, then $c \leq b$, and $c = b$ would force $\int(b - f) = 0$ with $b - f > 0$, which is impossible; similarly at the lower end, so
$c \in I$. With $m$ as above, $\varphi(f(x)) \geq \varphi(c) + m(f(x) - c)$; the right side is integrable, so $\int\varphi \circ f$ is defined in $(-\infty, \infty]$, and integrating, using $\mu(X) = 1$,
gives $\int\varphi \circ f \geq \varphi(c) + m(c - c) = \varphi(c)$.
\end{proof}

\begin{example}\label{measure-theory:example-jensen}
\begin{enumerate}
\item With the convex function $\varphi(t) = |t|^p$ on $\RR$, $p \geq 1$, applied to $|f|$: $(\int|f|)^p \leq \int|f|^p$ for a probability measure.
\item With $\varphi = \exp$: $\exp(\int f) \leq \int e^f$.
\item \emph{Arithmetic and geometric means.} Let $\mu$ give mass $\frac1n$ to each of $n$ points, where $f$ takes the values $a_1, \ldots, a_n > 0$, and $\varphi = -\log$. Then
$-\log(\frac1n\sum a_i) \leq -\frac1n\sum\log a_i$, that is, $(a_1\cdots a_n)^{1/n} \leq \frac1n(a_1 + \cdots + a_n)$.
\end{enumerate}
The hypothesis $\mu(X) = 1$ matters: for Lebesgue measure on $[0, 2]$ and $f = 1$, $\varphi(t) = t^2$ gives $\varphi(\int f) = 4 > 2 = \int\varphi(f)$.
\end{example}

\begin{theorem}[Dominated convergence]\label{measure-theory:theorem-dominated-convergence}
Let $f_n$ be measurable with $f_n \to f$ a.e. and $|f_n| \leq g$ a.e. for an integrable $g$. Then $f$ is integrable and $\int|f_n - f| \to 0$; in particular $\int f_n \to \int f$.
\end{theorem}

\begin{proof}
Modifying on a null set, we may assume convergence and domination everywhere. Then $|f| \leq g$, and Fatou's lemma applied to $2g - |f_n - f| \geq 0$ gives
$\int 2g \leq \int 2g - \limsup\int|f_n - f|$.
\end{proof}

\begin{proposition}\label{measure-theory:proposition-riemann-lebesgue}
For continuous $f \colon [a, b] \to \CC$, the Lebesgue integral $\int_{[a,b]}f\,d\lambda$ equals the Riemann integral.
\end{proposition}

\begin{proof}
For a partition $P$ with mesh tending to $0$, the step functions $\sum_i\min_{[t_{i-1}, t_i]}\operatorname{Re}f \cdot \mathbf{1}_{[t_{i-1}, t_i)}$ converge to $\operatorname{Re}f$ uniformly, and their Lebesgue integrals are
the lower sums (Theorem~\ref{measure-theory:theorem-lebesgue-measure}). Apply dominated convergence, and treat the imaginary part in the same way.
\end{proof}

\begin{proposition}[Abstract change of variables]\label{measure-theory:proposition-image-measure}
Let $f \colon (X, \mathcal{A}) \to (Y, \mathcal{B})$ be measurable, $\mu$ a measure on $\mathcal{A}$, and $g \colon Y \to [0, \infty]$ measurable. Then
\[
\int_Yg\,d(f_*\mu) = \int_Xg \circ f\,d\mu.
\]
A measurable $g \colon Y \to \CC$ is integrable for $f_*\mu$ if and only if $g \circ f$ is integrable for $\mu$, and then the same formula holds.
\end{proposition}

\begin{proof}
\emph{Step 1: indicators.} For $g = \mathbf{1}_B$, $g \circ f = \mathbf{1}_{f^{-1}(B)}$, and both sides equal $\mu(f^{-1}(B)) = f_*\mu(B)$.

\emph{Step 2: simple functions.} Both sides are additive and positively homogeneous in $g$ (Corollary~\ref{measure-theory:corollary-integral-properties}(1)), so the formula holds for nonnegative simple $g$.

\emph{Step 3: nonnegative functions.} Choose simple $0 \leq s_n \uparrow g$ (Proposition~\ref{measure-theory:proposition-measurable-functions}(4)). Then $s_n \circ f$ are simple and increase to $g \circ f$, and
monotone convergence on both sides passes the formula of Step 2 to the limit.

\emph{Step 4: complex functions.} Applying Step 3 to $|g|$ gives $\int|g|\,d(f_*\mu) = \int|g \circ f|\,d\mu$, which is the statement on integrability. Since $(\operatorname{Re}g)^\pm \circ f = (\operatorname{Re}(g \circ f))^\pm$,
and similarly for the imaginary part, Step 3 applied to the four parts and Definition~\ref{measure-theory:definition-integral} give the formula.
\end{proof}

\begin{example}\label{measure-theory:example-image-measure}
\begin{enumerate}
\item \emph{Expectation through the distribution.} If $\mu$ is a probability measure and $f$ is a real measurable function with distribution $\nu = f_*\mu$, then $\int f\,d\mu = \int_\RR t\,d\nu(t)$ whenever $f \geq 0$ or $f$ is integrable
(apply the proposition to $g(t) = t$ on $[0, \infty)$, or to $g(t) = t^\pm$); all quantities depending only on the values of $f$ (mean, variance, $\int\varphi \circ f\,d\mu$) are integrals against $\nu$. For the sum $f$ of two fair dice, $\int f\,d\mu = \sum_{k=2}^{12}k\,\nu(\{k\}) = 7$.
\item \emph{Affine maps of $\RR$.} With $f(x) = ax + b$, $a \neq 0$, and $f_*\lambda = |a|^{-1}\lambda$ (Definition~\ref{measure-theory:definition-image-measure}), the proposition gives
$\int g(ax + b)\,d\lambda(x) = |a|^{-1}\int g\,d\lambda$, the one-dimensional case of the change of variables formula of Section~\ref{measure-theory:section-change-of-variables}.
\item \emph{Integrals on a circle.} Let $f \colon [0, 2\pi) \to S^1$, $f(\theta) = e^{i\theta}$, and let $\sigma = f_*\lambda$. Then $\int_{S^1}z^k\,d\sigma(z) = \int_0^{2\pi}e^{ik\theta}\,d\theta$, which is $2\pi$ for $k = 0$ and $0$
for other integers $k$. The measure $\sigma$ is invariant under rotations: the rotation by $e^{i\alpha}$ corresponds to the translation $\theta \mapsto \theta + \alpha$ modulo $2\pi$.
\end{enumerate}
\end{example}

\begin{proposition}\label{measure-theory:proposition-riemann-integrable}
Every bounded Riemann integrable function $f \colon [a, b] \to \RR$ (Definition~\ref{real-analysis:definition-riemann}) is Lebesgue measurable and Lebesgue integrable, with the same integral.
\end{proposition}

\begin{proof}
Choose partitions $P_1, P_2, \ldots$, each refining the previous one, with $U(f, P_k) - L(f, P_k) < 1/k$ (Proposition~\ref{real-analysis:proposition-riemann}(1)). Let $\phi_k$ and $\psi_k$ be the step functions equal, on the interior of each interval of $P_k$,
to the infimum and supremum of $f$ there, and to $f$ at the finitely many partition points. Then $\phi_k \leq f \leq \psi_k$, $\phi_k$ increases and $\psi_k$ decreases (refinement), and their Lebesgue integrals are $L(f, P_k)$ and $U(f, P_k)$,
since points have measure zero. The limits $\phi = \lim\phi_k$ and $\psi = \lim\psi_k$ are measurable, $\phi \leq f \leq \psi$, and by monotone convergence $\int(\psi - \phi) = \lim(U(f, P_k) - L(f, P_k)) = 0$. So $\psi = \phi$ a.e.
(Corollary~\ref{measure-theory:corollary-integral-properties}(4)), hence $f = \phi$ a.e., and $f$ is measurable because Lebesgue measure is complete. Finally $\int f\,d\lambda = \lim\int\phi_k\,d\lambda = \lim L(f, P_k)$, which is the Riemann integral.
\end{proof}

\begin{counterexample}\label{measure-theory:counterexample-dirichlet}
The converse fails. The Dirichlet function $\mathbf{1}_\QQ$ on $[0, 1]$ is Lebesgue integrable with integral $0$, since $\QQ$ is countable and hence of measure zero. It is not Riemann integrable: every nondegenerate interval contains rational
and irrational points, so $L(\mathbf{1}_\QQ, P) = 0$ and $U(\mathbf{1}_\QQ, P) = 1$ for every partition $P$. It is the increasing pointwise limit of the Riemann integrable functions $\mathbf{1}_{\{q_1, \ldots, q_n\}}$, each of integral $0$, for an
enumeration of $\QQ \cap [0, 1]$: the Riemann integral does not behave well under monotone limits, which is what the Lebesgue integral repairs (Theorem~\ref{measure-theory:theorem-monotone-convergence}).
\end{counterexample}

\begin{proposition}[Improper integrals]\label{measure-theory:proposition-improper}
Let $f \colon [a, \infty) \to \RR$ be continuous, and call $\lim_{b \to \infty}\int_a^bf$, when it exists, the \emph{improper (Riemann) integral} $\int_a^\infty f$.
\begin{enumerate}
\item If $f \geq 0$, then $\int_{[a, \infty)}f\,d\lambda = \lim_{b \to \infty}\int_a^bf$ in $[0, \infty]$.
\item $f \in L^1([a, \infty))$ if and only if $\int_a^\infty|f| < \infty$, and then the improper integral exists and equals the Lebesgue integral.
\item The improper integral can exist without $f$ being Lebesgue integrable: $\int_\pi^\infty\frac{\sin x}{x}\,dx$ converges, but $\int_\pi^\infty\frac{|\sin x|}{x}\,dx = \infty$.
\end{enumerate}
The same holds for continuous functions on $(a, b]$ with a singularity at $a$; for instance $\int_0^1x^{-p}\,dx$ is finite if and only if $p < 1$.
\end{proposition}

\begin{proof}
(1) $f\mathbf{1}_{[a, b]}$ increases to $f$ as $b \to \infty$ (along any sequence $b_n \uparrow \infty$), and $\int_{[a, b]}f\,d\lambda = \int_a^bf$ (Proposition~\ref{measure-theory:proposition-riemann-lebesgue});
apply monotone convergence. (2) The first statement is (1) for $|f|$. Then $|f\mathbf{1}_{[a, b]}| \leq |f|$, and dominated convergence gives $\int_a^bf \to \int f\,d\lambda$. (3) On
$[k\pi, (k + 1)\pi]$ the integral of $|\sin x|/x$ is at least $\frac{2}{(k + 1)\pi}$, and $\sum_k\frac1{k + 1} = \infty$. The integrals $c_k = \int_{k\pi}^{(k+1)\pi}\frac{\sin x}{x}\,dx$ alternate in sign,
decrease in absolute value (substitute $x \mapsto x + \pi$) and tend to $0$, so $\sum_kc_k$ converges (Lemma~\ref{real-analysis:lemma-alternating}); and for $b$ between $k\pi$ and $(k + 1)\pi$,
$\int_\pi^b$ differs from $\sum_{j < k}c_j$ by at most $|c_k| \to 0$. For $x^{-p}$, $\int_\varepsilon^1x^{-p}\,dx = \frac{1 - \varepsilon^{1-p}}{1 - p}$ for $p \neq 1$ and $-\log\varepsilon$ for $p = 1$.
\end{proof}

So for nonnegative functions the improper integral is the Lebesgue integral, and in general Lebesgue integrability is absolute convergence of the improper integral; the convergent improper integrals
of non-integrable functions are the analogue of conditionally convergent series.

\begin{proposition}[Differentiation under the integral sign]\label{measure-theory:proposition-differentiation-integral}
Let $(X, \mu)$ be a measure space, $I \subseteq \RR$ an open interval, and $f \colon X \times I \to \RR$ such that $f(\cdot, t)$ is integrable for every $t$, $f(x, \cdot)$ is differentiable for every $x$,
and $|\partial_tf(x, t)| \leq g(x)$ for an integrable $g$ independent of $t$. Then $F(t) = \int_Xf(x, t)\,d\mu(x)$ is differentiable, with $F^{\prime}(t) = \int_X\partial_tf(x, t)\,d\mu(x)$.
\end{proposition}

\begin{proof}
For $h_n \to 0$, $(f(x, t + h_n) - f(x, t))/h_n \to \partial_tf(x, t)$, and by the mean value theorem these quotients are bounded by $g(x)$. So $\partial_tf(\cdot, t)$ is measurable, as a limit of
measurable functions, and dominated convergence gives $(F(t + h_n) - F(t))/h_n \to \int\partial_tf(x, t)\,d\mu$.
\end{proof}

\begin{example}\label{measure-theory:example-feynman}
\emph{Feynman's trick.} Introducing a parameter and differentiating under the integral sign computes integrals that have no elementary antiderivative.
\begin{enumerate}
\item $F(t) = \int_0^1\frac{x^t - 1}{\log x}\,dx$ for $t > -1$: here $\partial_t\frac{x^t - 1}{\log x} = x^t$, dominated by $x^{s}$ for $t \geq s > -1$, so $F^{\prime}(t) = \int_0^1x^t\,dx = \frac{1}{t + 1}$, and
$F(0) = 0$ gives $F(t) = \log(1 + t)$.
\item $F(t) = \int_\RR e^{-x^2}\cos(tx)\,dx$. Differentiating, with domination by $|x|e^{-x^2}$, and integrating by parts,
$F^{\prime}(t) = -\int_\RR xe^{-x^2}\sin(tx)\,dx = -\frac t2F(t)$; so $F(t) = F(0)e^{-t^2/4} = \sqrt\pi\,e^{-t^2/4}$ by Example~\ref{measure-theory:example-gaussian}. This is the Fourier transform of the
Gaussian.
\item \emph{The Dirichlet integral} $\int_0^\infty\frac{\sin x}{x}\,dx = \frac\pi2$ (an improper integral, Proposition~\ref{measure-theory:proposition-improper}(3)). For $t > 0$ let
$F(t) = \int_0^\infty e^{-tx}\frac{\sin x}{x}\,dx$, a Lebesgue integral. On $t \geq \delta > 0$ the derivative $-e^{-tx}\sin x$ is dominated by $e^{-\delta x}$, so
$F^{\prime}(t) = -\int_0^\infty e^{-tx}\sin x\,dx = -\operatorname{Im}\int_0^\infty e^{-(t - i)x}\,dx = -\operatorname{Im}\frac{1}{t - i} = -\frac{1}{1 + t^2}$. Since $|\sin x/x| \leq 1$,
$|F(t)| \leq 1/t \to 0$ as $t \to \infty$; hence $F(t) = \frac\pi2 - \arctan t$, with $\arctan$ the inverse of $\tan \colon (-\frac\pi2, \frac\pi2) \to \RR$. Finally $F(t) \to \int_0^\infty\frac{\sin x}x$ as
$t \to 0^+$: for a nonincreasing $C^1$ function $u \geq 0$, integrating by parts, $|\int_R^Su(x)\sin x\,dx| \leq 2u(R)$; with $u(x) = e^{-tx}/x$ this bounds the tails beyond $R$ by $2/R$ for all
$t \geq 0$, while $|\int_0^R(e^{-tx} - 1)\frac{\sin x}{x}\,dx| \leq \int_0^Rtx\,dx = \frac{tR^2}2$. Choosing first $R$ large and then $t$ small gives the limit.
\end{enumerate}
\end{example}

\section{Lebesgue spaces}\label{measure-theory:section-lp}

\begin{definition}\label{measure-theory:definition-lp}
For $1 \leq p < \infty$, $L^p(\mu)$ is the space of measurable $f \colon X \to \CC$ with $\|f\|_p = (\int|f|^p)^{1/p} < \infty$, modulo equality a.e. $L^\infty(\mu)$ consists of the essentially bounded
functions with $\|f\|_\infty = \inf\{c : |f| \leq c \text{ a.e.}\}$, modulo equality a.e. On $L^2(\mu)$ we have the inner product $\langle f, g \rangle = \int f\bar g$.
\end{definition}

\begin{notation}\label{measure-theory:notation-lp}
As usual, an element of $L^p(\mu)$ is written as a function $f$, although it is a class of functions equal almost everywhere; integrals, norms and a.e.\ statements do not depend on the representative.
For $1 \leq p \leq \infty$ the \emph{conjugate exponent} $p^{\prime} = q$ is defined by $1/p + 1/q = 1$, with $1/\infty = 0$; so $1$ and $\infty$ are conjugate, and $2$ is its own conjugate. When $\mu$ is counting
measure on a set $S$ we write $\ell^p(S)$, and $\ell^p = \ell^p(\NN)$: the sequences $(a_n)$ with $\sum_n|a_n|^p < \infty$, respectively the bounded sequences.
\end{notation}

The infimum defining $\|f\|_\infty$ is attained: $|f| \leq \|f\|_\infty$ a.e., since $\{|f| > \|f\|_\infty\} = \bigcup_k\{|f| > \|f\|_\infty + 1/k\}$ is a countable union of null sets. For counting
measure only the empty set is null, and $\|f\|_\infty = \sup|f|$; for Lebesgue measure, values on null sets are ignored: $\|\mathbf{1}_\QQ\|_\infty = 0$, and the function on $[0, 1]$ equal to $x$ at irrational $x$ and to $5$ at
rational $x$ has $L^\infty$ norm $1$, although its supremum is $5$. Identifying functions that agree a.e.\ is needed for $\|\cdot\|_p$ to be a norm: $\|f\|_p = 0$
means only that $f = 0$ a.e.

The triangle inequality for $\|\cdot\|_p$ is not obvious for $p \neq 1, 2, \infty$; it follows from H\"older's inequality, which in turn rests on an elementary inequality for numbers.

\begin{lemma}[Young's inequality for products]\label{measure-theory:lemma-young-product}
Let $1 < p < \infty$ with conjugate $q$, and $a, b \geq 0$. Then $ab \leq \dfrac{a^p}{p} + \dfrac{b^q}{q}$, with equality if and only if $a^p = b^q$.
\end{lemma}

\begin{proof}
If $a = 0$ or $b = 0$ the inequality is clear, with equality exactly when both vanish. Let $a, b > 0$. The function $\log$ is strictly concave on $(0, \infty)$, since its derivative $1/t$ is strictly decreasing;
with the weights $1/p + 1/q = 1$,
\[
\log\Bigl(\frac{a^p}{p} + \frac{b^q}{q}\Bigr) \geq \frac1p\log a^p + \frac1q\log b^q = \log(ab),
\]
with equality if and only if the two points $a^p$ and $b^q$ coincide. Apply $\exp$.
\end{proof}

For $p = q = 2$ the lemma is $ab \leq \frac12(a^2 + b^2)$, that is, $(a - b)^2 \geq 0$. Geometrically: the area $ab$ of a rectangle is at most the area under the curve $y = x^{p-1}$ from $0$ to $a$ plus
the area to the left of the same curve, $x = y^{q-1}$, from $0$ to $b$.

\begin{proposition}\label{measure-theory:proposition-holder}
Let $1 \leq p, q \leq \infty$ with $1/p + 1/q = 1$. Then $\int|fg| \leq \|f\|_p\|g\|_q$ (H\"older) and $\|f + g\|_p \leq \|f\|_p + \|g\|_p$ (Minkowski). So $L^p(\mu)$ is a normed space.
\end{proposition}

\begin{proof}
The inequalities are statements about the measurable functions $|f|$, $|g|$, with values in $[0, \infty]$ and norms in $[0, \infty]$ defined by the same formulas; we prove them in that generality, with
the convention $0 \cdot \infty = 0$.

\emph{Step 1: H\"older for $p = 1$, $q = \infty$.} Since $|g| \leq \|g\|_\infty$ a.e., $|fg| \leq |f|\,\|g\|_\infty$ a.e., and integrating gives $\int|fg| \leq \|f\|_1\|g\|_\infty$. The case $p = \infty$ is symmetric.

\emph{Step 2: degenerate cases for $1 < p < \infty$.} If $\|f\|_p = 0$, then $f = 0$ a.e., so $fg = 0$ a.e.\ and both sides vanish (the right side is $0$ even if $\|g\|_q = \infty$). The same holds if
$\|g\|_q = 0$. If neither norm is $0$ and one of them is $\infty$, the right side is $\infty$.

\emph{Step 3: the main case.} Let $0 < \|f\|_p, \|g\|_q < \infty$. Both sides of H\"older's inequality are multiplied by $st$ when $f$, $g$ are replaced by $sf$, $tg$ ($s, t > 0$); so we may assume
$\|f\|_p = \|g\|_q = 1$. By Lemma~\ref{measure-theory:lemma-young-product} applied pointwise, $|fg| \leq |f|^p/p + |g|^q/q$, and integrating,
\[
\int|fg| \leq \frac1p\int|f|^p + \frac1q\int|g|^q = \frac1p + \frac1q = 1 = \|f\|_p\|g\|_q.
\]

\emph{Step 4: Minkowski for $p = 1$ and $p = \infty$.} For $p = 1$ integrate $|f + g| \leq |f| + |g|$. For $p = \infty$, $|f| \leq \|f\|_\infty$ and $|g| \leq \|g\|_\infty$ outside a null set, so
$|f + g| \leq \|f\|_\infty + \|g\|_\infty$ a.e.

\emph{Step 5: Minkowski for $1 < p < \infty$.} First, $\|f + g\|_p$ is finite when $\|f\|_p$ and $\|g\|_p$ are: $|f + g| \leq 2\max(|f|, |g|)$, so $|f + g|^p \leq 2^p(|f|^p + |g|^p)$. (If $\|f\|_p$ or
$\|g\|_p$ is infinite there is nothing to prove.) Write
\[
|f + g|^p \leq |f|\,|f + g|^{p-1} + |g|\,|f + g|^{p-1}.
\]
The conjugate exponent satisfies $(p - 1)q = p$, so $\||f + g|^{p-1}\|_q = (\int|f + g|^p)^{1/q} = \|f + g\|_p^{p/q}$, and H\"older applied to both terms gives
$\|f + g\|_p^p \leq (\|f\|_p + \|g\|_p)\|f + g\|_p^{p/q}$. If $\|f + g\|_p = 0$ there is nothing to prove; otherwise divide by the finite number $\|f + g\|_p^{p/q}$, using $p - p/q = 1$.

\emph{Step 6: norm.} By Step 5 and Step 4, $L^p(\mu)$ is closed under addition, and clearly under scalar multiplication, with $\|cf\|_p = |c|\,\|f\|_p$. If $\|f\|_p = 0$ with $p < \infty$, then $\int|f|^p = 0$
and $f = 0$ a.e.\ by Corollary~\ref{measure-theory:corollary-integral-properties}(4); for $p = \infty$ this is the remark that $|f| \leq \|f\|_\infty$ a.e. So $\|\cdot\|_p$ is a norm on classes modulo
equality a.e.
\end{proof}

\begin{remark}\label{measure-theory:remark-holder-equality}
\emph{Equality in H\"older.} Let $1 < p < \infty$ and $0 < \|f\|_p, \|g\|_q < \infty$. Equality holds in Step 3 after normalisation if and only if $|f|^p = |g|^q$ a.e., by the equality case of
Lemma~\ref{measure-theory:lemma-young-product} (the integral of the nonnegative difference vanishes exactly when the difference is $0$ a.e.). Undoing the normalisation: $\int|fg| = \|f\|_p\|g\|_q$ if
and only if $\alpha|f|^p = \beta|g|^q$ a.e.\ for some constants $\alpha, \beta \geq 0$ not both zero. For instance, given $f \in L^p$ with $\|f\|_p = 1$, the function $g = |f|^{p-1}$ satisfies
$\|g\|_q = 1$ and $\int|f|g = 1$; this choice of $g$ is the one used in Step 5 and in the duality of $L^p$ spaces (Chapter~\ref{signed-measures}).
\end{remark}

\begin{example}\label{measure-theory:example-holder}
\begin{enumerate}
\item \emph{Cauchy--Schwarz.} For $p = q = 2$, $|\langle f, g \rangle| \leq \int|f\bar g| \leq \|f\|_2\|g\|_2$; for counting measure on $\{1, \ldots, n\}$ this is
$|\sum_ia_ib_i| \leq (\sum_i|a_i|^2)^{1/2}(\sum_i|b_i|^2)^{1/2}$.
\item \emph{Sums.} For counting measure, H\"older reads $\sum_n|a_nb_n| \leq (\sum_n|a_n|^p)^{1/p}(\sum_n|b_n|^q)^{1/q}$. With $b_n = 1$ for $n \leq N$ and $b_n = 0$ otherwise: $\sum_{n \leq N}|a_n| \leq N^{1/q}(\sum_{n \leq N}|a_n|^p)^{1/p}$, a
power mean inequality.
\item \emph{Averages.} For a probability measure and $1 \leq p \leq r \leq \infty$, $\|f\|_p \leq \|f\|_r$: apply H\"older to $|f|^p \cdot 1$ with exponents $r/p$ and its conjugate. So
$L^r(\mu) \subseteq L^p(\mu)$, a special case of Exercise~\ref{measure-theory:exercise-lp-inclusions}. For counting measure the inclusions go the other way: $\ell^p \subseteq \ell^r$, since a sequence in $\ell^p$
has $|a_n| \leq 1$ for large $n$, and then $|a_n|^r \leq |a_n|^p$.
\item \emph{Three functions.} If $1/p + 1/q + 1/r = 1$, then $\int|fgh| \leq \|f\|_p\|g\|_q\|h\|_r$: apply H\"older with exponents $p$ and $p^{\prime}$, where $1/p^{\prime} = 1/q + 1/r$, and then to
$\|gh\|_{p^{\prime}}^{p^{\prime}} = \int|g|^{p^{\prime}}|h|^{p^{\prime}}$ with exponents $q/p^{\prime}$ and $r/p^{\prime}$.
\end{enumerate}
\end{example}

\begin{proposition}[Interpolation of $L^p$ norms]\label{measure-theory:proposition-lp-interpolation}
Let $1 \leq p < r < s \leq \infty$ and $0 < \theta < 1$ with $\frac1r = \frac\theta p + \frac{1 - \theta}s$. Then $\|f\|_r \leq \|f\|_p^\theta\|f\|_s^{1 - \theta}$; in particular
$L^p(\mu) \cap L^s(\mu) \subseteq L^r(\mu)$.
\end{proposition}

\begin{proof}
If $s = \infty$, then $\theta = p/r$ and $\int|f|^r \leq \|f\|_\infty^{r - p}\int|f|^p$, which is the claim raised to the power $r$. If $s < \infty$, write $|f|^r = |f|^{\theta r}|f|^{(1 - \theta)r}$ and apply
H\"older with the conjugate exponents $p/(\theta r)$ and $s/((1 - \theta)r)$, whose reciprocals add up to $\theta r/p + (1 - \theta)r/s = 1$:
\[
\int|f|^r \leq \Bigl(\int|f|^p\Bigr)^{\theta r/p}\Bigl(\int|f|^s\Bigr)^{(1 - \theta)r/s} = \bigl(\|f\|_p^\theta\|f\|_s^{1 - \theta}\bigr)^r.
\]
\end{proof}

In other words, $1/r \mapsto \log\|f\|_r$ is convex. Neither inclusion $L^p \subseteq L^s$ nor $L^s \subseteq L^p$ holds in general (Exercise~\ref{measure-theory:exercise-lp-inclusions}), but the
intermediate spaces always contain the intersection.

\begin{counterexample}\label{measure-theory:counterexample-p-less-than-one}
For $0 < p < 1$ the expression $\|f\|_p = (\int|f|^p)^{1/p}$ is not a norm. On $[0, 1]$ with Lebesgue measure let $f = \mathbf{1}_{[0, 1/2]}$ and $g = \mathbf{1}_{(1/2, 1]}$. Then
$\|f\|_p = \|g\|_p = 2^{-1/p}$ and $\|f + g\|_p = 1$, while $\|f\|_p + \|g\|_p = 2^{1 - 1/p} < 1$: the triangle inequality fails. The reason is that $t \mapsto t^p$ is concave for $p < 1$, so the unit
ball $\{\|f\|_p \leq 1\}$ is not convex. H\"older's inequality also needs $p \geq 1$: for $p = \frac12$ the formula $1/p + 1/q = 1$ forces $q = -1$, which has no meaning as an exponent of an $L^q$ space.
\end{counterexample}

\begin{theorem}[Riesz--Fischer]\label{measure-theory:theorem-riesz-fischer}
$L^p(\mu)$ is a Banach space for $1 \leq p \leq \infty$, and every convergent sequence in $L^p$, $p < \infty$, has a subsequence converging a.e.
\end{theorem}

\begin{proof}
Let $p < \infty$ and $(f_n)$ Cauchy. Choose $n_1 < n_2 < \cdots$ with $\|f_{n_{k+1}} - f_{n_k}\|_p < 2^{-k}$, and let $g = \sum_k|f_{n_{k+1}} - f_{n_k}|$. By monotone convergence and Minkowski,
$\|g\|_p \leq 1$, so $g < \infty$ a.e., and $f = f_{n_1} + \sum_k(f_{n_{k+1}} - f_{n_k})$ converges absolutely a.e. Since $|f - f_{n_k}| \leq \sum_{j \geq k}|f_{n_{j+1}} - f_{n_j}|$, Minkowski gives
$\|f - f_{n_k}\|_p \leq 2^{1-k}$, and a Cauchy sequence with a convergent subsequence converges. For $p = \infty$, outside a null set $|f_n - f_m| \leq \|f_n - f_m\|_\infty$ for all $n, m$, so $(f_n)$ converges
uniformly there.
\end{proof}

\begin{proposition}\label{measure-theory:proposition-density}
For $1 \leq p < \infty$, the continuous functions with compact support are dense in $L^p(\RR^n, \lambda)$.
\end{proposition}

\begin{proof}
By Proposition~\ref{measure-theory:proposition-measurable-functions}(4) and dominated convergence, simple functions $\sum a_i\mathbf{1}_{E_i}$ with $\lambda(E_i) < \infty$ are dense. Given $E$ with $\lambda(E) < \infty$ and
$\varepsilon > 0$, choose compact $K \subseteq E \subseteq U$ open with $\lambda(U \setminus K) < \varepsilon$ (Theorem~\ref{measure-theory:theorem-lebesgue-measure}) and, by
Lemma~\ref{general-topology:lemma-lch-shrinking}, a continuous $\phi \colon \RR^n \to [0, 1]$ with compact support in $U$ and $\phi = 1$ on $K$. Then $\|\mathbf{1}_E - \phi\|_p^p \leq \lambda(U \setminus K) < \varepsilon$.
\end{proof}

\begin{counterexample}\label{measure-theory:counterexample-density-infinity}
The proposition fails for $p = \infty$. Let $g$ be continuous. If $|g| \leq c$ almost everywhere on an open set $W$, then $|g| \leq c$ everywhere on $W$, because $\{x \in W : |g(x)| > c\}$ is open and null,
and a nonempty open set has positive measure. Hence $\|g\|_\infty = \sup|g|$, a sequence in $C_c(\RR^n)$ that is Cauchy in $L^\infty$ is uniformly Cauchy, and the closure of $C_c(\RR^n)$ in $L^\infty$
consists of (classes of) continuous functions vanishing at infinity. For instance, let $H = \{x : x_1 \geq 0\}$ and suppose $\|g - \mathbf{1}_H\|_\infty = c < \frac12$ with $g$ continuous. Applying the
first remark to $g - 1$ on the interior of $H$ and to $g$ on $\RR^n \setminus H$, we get $|g - 1| \leq c$ on $\operatorname{int}H$ and $|g| \leq c$ on $\RR^n \setminus H$; at the origin, which is a limit of points
of both sets, continuity gives $g(0) \geq 1 - c > \frac12$ and $g(0) \leq c < \frac12$, a contradiction. So $\mathbf{1}_H$ has distance at least $\frac12$ from all continuous functions. In fact
$L^\infty([0, 1])$ is not separable at all: the functions $\mathbf{1}_{[0, t]}$, $0 < t \leq 1$, are uncountably many and pairwise at distance $1$.
\end{counterexample}

\section{Modes of convergence}\label{measure-theory:section-modes}

A sequence of functions can converge in many ways, and on a measure space there are several natural ones in addition to pointwise and uniform convergence: one may neglect null sets, neglect sets of
small measure, or measure the size of $f_n - f$ by an integral. Each of the convergence theorems of Section~\ref{measure-theory:section-integral} passes from one mode to another (dominated
convergence: from almost everywhere convergence to $L^1$ convergence), so it is worth knowing exactly which implications hold.

\begin{definition}\label{measure-theory:definition-modes}
Let $f_n$, $f$ be measurable functions on a measure space $(X, \mu)$ with values in $\CC$. Then $f_n \to f$ \emph{almost everywhere} if $f_n(x) \to f(x)$ for $x$ outside a null set; \emph{in measure} if
$\mu(\{|f_n - f| > \varepsilon\}) \to 0$ for every $\varepsilon > 0$; and \emph{in $L^p$} if $\|f_n - f\|_p \to 0$.
\end{definition}

\begin{definition}\label{measure-theory:definition-almost-uniform}
In the situation of Definition~\ref{measure-theory:definition-modes}, $f_n \to f$ \emph{almost uniformly} if for every $\delta > 0$ there is a measurable set $B$ with $\mu(B) < \delta$ such that
$f_n \to f$ uniformly on $X \setminus B$. The sequence is \emph{Cauchy in measure} if for every $\varepsilon > 0$, $\mu(\{|f_n - f_m| > \varepsilon\}) \to 0$ as $n, m \to \infty$.
\end{definition}

\begin{notation}\label{measure-theory:notation-modes}
We write $f_n \to f$ a.e., $f_n \to f$ in measure, and $f_n \to f$ in $L^p$; in all three cases the limit is determined only up to equality almost everywhere. Uniform convergence on $X$ is convergence in $L^\infty$ for counting measure; for a general $\mu$,
convergence in $L^\infty(\mu)$ is uniform convergence outside a null set.
\end{notation}

In probability, where $\mu(X) = 1$ and measurable functions are random variables, almost everywhere convergence is called convergence \emph{almost surely}, and convergence in measure is
convergence \emph{in probability}. The difference between the two is the difference between the strong and the weak law of large numbers.

\begin{example}\label{measure-theory:example-modes}
\begin{enumerate}
\item On $[0, 1]$ with Lebesgue measure, $f_n(x) = x^n$ converges to $0$ at every $x < 1$, hence almost everywhere; not uniformly, since $\sup_{[0, 1)}x^n = 1$; but uniformly on $[0, 1 - \delta]$ for every $\delta > 0$,
hence almost uniformly; in measure, since $\{x^n > \varepsilon\} = (\varepsilon^{1/n}, 1]$ has length $1 - \varepsilon^{1/n} \to 0$; and in every $L^p$, $p < \infty$, since $\|x^n\|_p^p = 1/(np + 1)$.
It does not converge to $0$ in $L^\infty$: $\|x^n\|_\infty = 1$.
\item For counting measure on $\NN$, a set of measure less than $1$ is empty. So convergence in measure is uniform convergence, almost uniform convergence is uniform convergence, and almost
everywhere convergence is pointwise convergence.
\item On $[0, 1]$, $f_n = \mathbf{1}_{[0, 1/n]}$ converges to $0$ in measure and in every $L^p$, $p < \infty$, and almost everywhere (at every $x > 0$), but not at $x = 0$: modes of convergence that
neglect null sets cannot see a single point.
\end{enumerate}
\end{example}

\begin{theorem}[Egorov]\label{measure-theory:theorem-egorov}
Let $\mu(X) < \infty$ and $f_n \to f$ almost everywhere. For every $\delta > 0$ there is a measurable $E$ with $\mu(X \setminus E) < \delta$ such that $f_n \to f$ uniformly on $E$.
\end{theorem}

\begin{proof}
\emph{Step 1: the bad sets.} For $k, n \geq 1$ let $E_{k,n} = \bigcup_{m \geq n}\{|f_m - f| \geq 1/k\}$, the set of points where some $f_m$ with $m \geq n$ is $1/k$-far from $f$. For fixed $k$ these sets
decrease in $n$, and their intersection consists of points where $|f_m - f| \geq 1/k$ for infinitely many $m$, a null set.

\emph{Step 2: finiteness.} As $\mu(X) < \infty$, continuity from above (Proposition~\ref{measure-theory:proposition-measure-properties}) gives $\mu(E_{k,n}) \to 0$ as $n \to \infty$.

\emph{Step 3: diagonal choice.} Choose $n_k$ with $\mu(E_{k,n_k}) < \delta2^{-k}$ and let $E = X \setminus \bigcup_kE_{k,n_k}$. Then $\mu(X \setminus E) < \sum_k\delta2^{-k} = \delta$, and on $E$,
$|f_m - f| < 1/k$ for all $m \geq n_k$, for every $k$; that is, $f_n \to f$ uniformly on $E$.
\end{proof}

In other words, on a finite measure space almost everywhere convergence implies almost uniform convergence. The converse holds on every measure space: if $f_n \to f$ uniformly outside sets $B_k$
with $\mu(B_k) < 1/k$, then $f_n \to f$ outside $\bigcap_kB_k$, a null set. Egorov's theorem is the second of \emph{Littlewood's three principles} (Section~\ref{measure-theory:section-lusin}): a
convergent sequence of measurable functions is nearly uniformly convergent.

\begin{proposition}\label{measure-theory:proposition-modes}
\begin{enumerate}
\item If $f_n \to f$ in $L^p$ for some $1 \leq p < \infty$, then $f_n \to f$ in measure.
\item If $f_n \to f$ in measure, a subsequence converges to $f$ almost everywhere.
\item If $\mu(X) < \infty$ and $f_n \to f$ almost everywhere, then $f_n \to f$ in measure.
\item If $f_n \to f$ almost uniformly, then $f_n \to f$ almost everywhere and in measure.
\item If $f_n \to f$ in $L^\infty$, then $f_n \to f$ almost uniformly, and in $L^p$ for all $p < \infty$ if $\mu(X) < \infty$.
\end{enumerate}
\end{proposition}

\begin{proof}
(1) By Markov's inequality (Corollary~\ref{measure-theory:corollary-integral-properties}(4)) applied to $|f_n - f|^p$, $\mu(\{|f_n - f| > \varepsilon\}) \leq \varepsilon^{-p}\|f_n - f\|_p^p$.

(2) Choose $n_1 < n_2 < \cdots$ with $\mu(A_k) < 2^{-k}$, where $A_k = \{|f_{n_k} - f| > 2^{-k}\}$. The set $A$ of points lying in infinitely many $A_k$ is contained in $\bigcup_{k \geq K}A_k$ for every $K$,
so $\mu(A) \leq \sum_{k \geq K}2^{-k} = 2^{1-K}$ for all $K$, and $\mu(A) = 0$. For $x \notin A$, $|f_{n_k}(x) - f(x)| \leq 2^{-k}$ for large $k$.

(3) Given $\varepsilon, \delta > 0$, Egorov's theorem gives $E$ with $\mu(X \setminus E) < \delta$ and $|f_n - f| < \varepsilon$ on $E$ for large $n$, so $\mu(\{|f_n - f| > \varepsilon\}) < \delta$ for large $n$.

(4) Almost everywhere convergence was shown after Egorov's theorem. Given $\varepsilon, \delta > 0$, let $\mu(B) < \delta$ with uniform convergence off $B$; for large $n$, $\{|f_n - f| > \varepsilon\} \subseteq B$.

(5) Outside the null set $\bigcup_n\{|f_n - f| > \|f_n - f\|_\infty\}$, $|f_n - f| \leq \|f_n - f\|_\infty \to 0$ uniformly. And $\|f_n - f\|_p \leq \mu(X)^{1/p}\|f_n - f\|_\infty$.
\end{proof}

The implications are summarised in the following diagram, where a label on an arrow is the hypothesis it needs, and the dashed arrow holds only for a subsequence (the
implication from a.e.\ convergence to convergence in measure for finite $\mu$, part (3), is the composite of two arrows); Counterexample~\ref{measure-theory:counterexample-modes} shows that no arrow can be added
without hypotheses.
\begin{center}
\begin{tikzcd}[column sep=large, row sep=large]
L^\infty \arrow[r] \arrow[d, "\mu(X) < \infty"'] & \text{almost uniform} \arrow[r, shift right] \arrow[d] & \text{a.e.} \arrow[l, shift right, "\mu(X) < \infty"'] \\
L^p \arrow[r] & \text{in measure} \arrow[ur, dashed, "\text{subsequence}"'] &
\end{tikzcd}
\end{center}

\begin{counterexample}\label{measure-theory:counterexample-modes}
No other implications hold in general. Take $X = [0, 1]$ or $\RR$ with Lebesgue measure.
\begin{enumerate}
\item $f_n = n\mathbf{1}_{(0, 1/n)}$ on $[0, 1]$ converges to $0$ everywhere, but $\int f_n = 1$: almost everywhere convergence and convergence in measure do not imply $L^1$ convergence, even on a finite measure space.
\item $f_n = \mathbf{1}_{[n, n+1]}$ on $\RR$ converges to $0$ everywhere but not in measure, as $\lambda(\{f_n > \frac12\}) = 1$: part (3) of Proposition~\ref{measure-theory:proposition-modes} and Egorov's theorem need $\mu(X) < \infty$.
\item (The typewriter sequence.) For $n = 2^k + j$ with $0 \leq j < 2^k$ let $f_n = \mathbf{1}_{[j2^{-k}, (j+1)2^{-k}]}$ on $[0, 1]$. Then $\|f_n\|_1 = 2^{-k} \to 0$, so $f_n \to 0$ in $L^1$ and in measure; but every $x \in [0, 1]$ lies in one of the
intervals for each $k$ and outside most of them, so $f_n(x) = 1$ and $f_n(x) = 0$ both happen infinitely often, and $(f_n(x))$ converges for no $x$. The subsequence in part (2) of
Proposition~\ref{measure-theory:proposition-modes} is necessary.
\item $f_n = n^{-1}\mathbf{1}_{[0, n]}$ on $\RR$ converges to $0$ uniformly, hence in $L^\infty$, almost uniformly and in measure, but $\|f_n\|_1 = 1$ and $\|f_n\|_p = n^{1/p - 1} \to 0$ only for
$p > 1$: on an infinite measure space $L^\infty$ convergence does not imply $L^1$ convergence.
\item Egorov's theorem cannot be improved to $\mu(X \setminus E) = 0$: for $f_n(x) = x^n$ on $[0, 1]$, if $f_n \to 0$ uniformly on $E$, then $E \cap (1 - \eta, 1) = \emptyset$ for some $\eta > 0$, since
$x^n \geq \frac12$ on $(2^{-1/n}, 1)$; so $\lambda([0, 1] \setminus E) \geq \eta > 0$.
\end{enumerate}
\end{counterexample}

\begin{theorem}[Riesz]\label{measure-theory:theorem-cauchy-in-measure}
A sequence of measurable functions that is Cauchy in measure converges in measure to a measurable function $f$, unique up to equality almost everywhere, and a subsequence converges to $f$
almost everywhere.
\end{theorem}

\begin{proof}
\emph{Step 1: a fast subsequence.} Choose $n_1 < n_2 < \cdots$ such that $A_k = \{|f_{n_{k+1}} - f_{n_k}| > 2^{-k}\}$ has $\mu(A_k) < 2^{-k}$. Let $B_K = \bigcup_{k \geq K}A_k$, so
$\mu(B_K) < 2^{1-K}$, and $N = \bigcap_KB_K$, a null set.

\emph{Step 2: almost uniform convergence.} For $x \notin B_K$ and $K \leq k < l$, $|f_{n_l}(x) - f_{n_k}(x)| \leq \sum_{j \geq k}2^{-j} = 2^{1-k}$. So $(f_{n_k})$ is uniformly Cauchy off $B_K$, and
converges off $N$; let $f$ be its limit there and $f = 0$ on $N$, a measurable function (Proposition~\ref{measure-theory:proposition-measurable-functions}(3)). Off $B_K$, $|f - f_{n_k}| \leq 2^{1-k}$
for $k \geq K$.

\emph{Step 3: the whole sequence.} Given $\varepsilon > 0$, $\{|f_n - f| > \varepsilon\} \subseteq \{|f_n - f_{n_k}| > \varepsilon/2\} \cup \{|f_{n_k} - f| > \varepsilon/2\}$. For $k$ large, with
$2^{1-k} < \varepsilon/2$, the second set lies in $B_k$, of measure $< 2^{1-k}$; and the first has small measure for $n$ and $n_k$ large, by the Cauchy condition. So $f_n \to f$ in measure.

\emph{Step 4: uniqueness.} If $f_n \to f$ and $f_n \to g$ in measure, then $\{|f - g| > 2\varepsilon\} \subseteq \{|f - f_n| > \varepsilon\} \cup \{|f_n - g| > \varepsilon\}$, so
$\mu(\{|f - g| > 2\varepsilon\}) = 0$ for all $\varepsilon > 0$, and $f = g$ a.e.
\end{proof}

\begin{proposition}\label{measure-theory:proposition-measure-metric}
Let $\mu$ be a finite measure, and for measurable $f, g \colon X \to \CC$ put $\rho(f, g) = \int\min(1, |f - g|)\,d\mu$. Then $\rho$ is a metric on the space $L^0(\mu)$ of measurable functions modulo
equality a.e., $f_n \to f$ in measure if and only if $\rho(f_n, f) \to 0$, and $(L^0(\mu), \rho)$ is complete.
\end{proposition}

\begin{proof}
\emph{Metric.} Symmetry is clear; $\rho(f, g) = 0$ means $\min(1, |f - g|) = 0$ a.e., that is, $f = g$ a.e. The triangle inequality follows from the pointwise inequality
$\min(1, s + t) \leq \min(1, s) + \min(1, t)$ for $s, t \geq 0$, applied to $s = |f - h|$, $t = |h - g|$ together with $|f - g| \leq s + t$.

\emph{Comparison.} For $0 < \varepsilon \leq 1$, Markov's inequality gives $\varepsilon\,\mu(\{|f - g| > \varepsilon\}) \leq \rho(f, g)$; and splitting $X$ into $\{|f - g| > \varepsilon\}$ and its
complement, $\rho(f, g) \leq \mu(\{|f - g| > \varepsilon\}) + \varepsilon\mu(X)$. The first inequality shows that $\rho$-convergence implies convergence in measure, the second the converse.

\emph{Completeness.} By the comparison, a $\rho$-Cauchy sequence is Cauchy in measure, and Theorem~\ref{measure-theory:theorem-cauchy-in-measure} gives a limit in measure, which is a $\rho$-limit.
\end{proof}

\begin{remark}\label{measure-theory:remark-ae-not-topological}
\emph{Almost everywhere convergence is not topological.} On $[0, 1]$ with Lebesgue measure there is no topology on $L^0$ whose convergent sequences are exactly the almost everywhere convergent
ones. In a topological space, if every subsequence of $(f_n)$ has a further subsequence converging to $f$, then $f_n \to f$: otherwise some neighbourhood of $f$ misses infinitely many $f_n$,
and that subsequence has no subsequence converging to $f$. The typewriter sequence (Counterexample~\ref{measure-theory:counterexample-modes}(3)) converges to $0$ in measure, so every subsequence
does, and by Proposition~\ref{measure-theory:proposition-modes}(2) has a further subsequence converging to $0$ almost everywhere; but the sequence itself converges nowhere. Convergence in measure, by
contrast, is the convergence of the metric $\rho$; it is the ``topological shadow'' of almost everywhere convergence: $f_n \to f$ in measure if and only if every subsequence has a further subsequence
converging almost everywhere to $f$ (for finite $\mu$; one direction is part (2), the other follows from part (3) and the subsequence principle just used).
\end{remark}

\begin{proposition}[Dominated convergence in measure]\label{measure-theory:proposition-dominated-in-measure}
Let $f_n \to f$ in measure and $|f_n| \leq g$ a.e.\ with $g$ integrable. Then $f$ is integrable and $\int|f_n - f| \to 0$.
\end{proposition}

\begin{proof}
Suppose not. Then there are $\eta > 0$ and a subsequence with $\int|f_{n_k} - f| \geq \eta$ for all $k$. This subsequence still converges to $f$ in measure, so by
Proposition~\ref{measure-theory:proposition-modes}(2) a further subsequence converges to $f$ almost everywhere, and by the dominated convergence theorem
(Theorem~\ref{measure-theory:theorem-dominated-convergence}) $f$ is integrable and the integrals $\int|f_{n_{k_j}} - f|$ tend to $0$, a contradiction. (That $f$ is integrable follows from the
same subsequence: $|f| \leq g$ a.e.)
\end{proof}

\section{Independence and the Borel--Cantelli lemmas}\label{measure-theory:section-borel-cantelli}

Measure theory is also the language of probability. A \emph{probability space} is a measure space $(\Omega, \mathcal{F}, P)$ with $P(\Omega) = 1$; its measurable sets are called \emph{events}, measurable
functions are \emph{random variables}, and ``almost everywhere'' becomes \emph{almost surely}. A typical question asks whether infinitely many of a sequence of events happen. For events $A_1, A_2, \ldots$
the event ``$A_k$ happens for infinitely many $k$'' is
\[
\limsup_kA_k = \bigcap_{n \geq 1}\bigcup_{k \geq n}A_k = \{\omega : \omega \in A_k \text{ for infinitely many } k\},
\]
and its complement is $\liminf_kA_k^c = \bigcup_n\bigcap_{k \geq n}A_k^c$, ``from some index on, no $A_k$ happens''. The two Borel--Cantelli lemmas decide, in many cases, whether $\limsup_kA_k$ has probability
$0$ or $1$. The first needs nothing; the second needs independence, which we now define.

\begin{definition}\label{measure-theory:definition-independent}
Events $A_1, A_2, \ldots$ (finitely or countably many) in a probability space are \emph{independent} if
\[
P(A_{k_1} \cap A_{k_2} \cap \cdots \cap A_{k_m}) = P(A_{k_1})P(A_{k_2}) \cdots P(A_{k_m})
\]
for all $m$ and all distinct indices $k_1, \ldots, k_m$.
\end{definition}

The condition must be imposed on all subfamilies, not only on pairs.

\begin{example}\label{measure-theory:example-pairwise-independent}
Toss two fair coins: $\Omega = \{HH, HT, TH, TT\}$, each point with probability $\frac14$. Let $A$ = ``the first coin shows heads'', $B$ = ``the second coin shows heads'', $C$ = ``the two coins agree''. Each
event has probability $\frac12$, and each intersection of two of them has probability $\frac14$; so they are pairwise independent. But $A \cap B \cap C = \{HH\}$ has probability $\frac14 \neq \frac18$: if
two of the events happen, so does the third.
\end{example}

\begin{lemma}\label{measure-theory:lemma-independent-complements}
If $A_1, A_2, \ldots$ are independent, so are the events obtained by replacing any of the $A_k$ by their complements.
\end{lemma}

\begin{proof}
It suffices to replace one event, say $A_1$, by $A_1^c$; replacing several is done one at a time. Let $k_2, \ldots, k_m$ be distinct indices different from $1$ and $B = A_{k_2} \cap \cdots \cap A_{k_m}$ (with
$B = \Omega$ if $m = 1$). Then $P(A_1^c \cap B) = P(B) - P(A_1 \cap B) = P(B) - P(A_1)P(B) = P(A_1^c)P(B)$, and $P(B)$ is the product of the $P(A_{k_j})$. Intersections not involving $A_1^c$ are unchanged.
\end{proof}

The basic example of independence lives on $[0, 1)$ with Lebesgue measure: the binary digits of a number chosen at random behave like independent tosses of a fair coin. For $x \in [0, 1)$ and $k \geq 1$ let
\[
d_k(x) = \lfloor 2^kx\rfloor - 2\lfloor 2^{k-1}x\rfloor \in \{0, 1\},
\]
the $k$-th digit of the binary expansion $x = \sum_kd_k(x)2^{-k}$ that does not end with infinitely many $1$s. For $K \geq 1$ the integer $\lfloor 2^Kx\rfloor$ has binary digits $d_1(x), \ldots, d_K(x)$, so
$d_1(x) = \delta_1, \ldots, d_K(x) = \delta_K$ holds exactly on the interval $[N2^{-K}, (N + 1)2^{-K})$ with $N = \sum_{k \leq K}\delta_k2^{K-k}$, of length $2^{-K}$. Call a set $E \subseteq [0, 1)$
\emph{determined by} a finite set $S$ of positions if whether $x \in E$ depends only on the digits $d_k(x)$, $k \in S$.

\begin{lemma}\label{measure-theory:lemma-binary-digits}
\begin{enumerate}
\item For a finite set $S$ of positions and prescribed digits $\delta_k \in \{0, 1\}$, $k \in S$, the set $\{x : d_k(x) = \delta_k \text{ for } k \in S\}$ has measure $2^{-|S|}$.
\item If $E_1, \ldots, E_m$ are determined by pairwise disjoint finite sets of positions $S_1, \ldots, S_m$, then $\lambda(E_1 \cap \cdots \cap E_m) = \lambda(E_1) \cdots \lambda(E_m)$. Consequently, events
$E_1, E_2, \ldots$ determined by pairwise disjoint finite sets of positions are independent; in particular the events $\{d_k = 1\}$, $k \geq 1$, are independent of probability $\frac12$.
\item If $g_i = G_i\bigl((d_k)_{k \in S_i}\bigr)$ for functions $G_i$ of finitely many digits and pairwise disjoint $S_1, \ldots, S_m$, then $\int_0^1g_1 \cdots g_m\,d\lambda = \prod_i\int_0^1g_i\,d\lambda$.
\end{enumerate}
\end{lemma}

\begin{proof}
(1) Let $K = \max S$ (the case $S = \emptyset$ is trivial). The set is the disjoint union, over the $2^{K - |S|}$ choices of the digits at the positions in $\{1, \ldots, K\} \setminus S$, of intervals of length
$2^{-K}$, so its measure is $2^{K - |S|}2^{-K}$.

(2) A set determined by $S_i$ is a disjoint union of $n_i$ sets as in (1), one for each admissible pattern of digits on $S_i$, so $\lambda(E_i) = n_i2^{-|S_i|}$. With $S = S_1 \cup \cdots \cup S_m$, the
intersection $E_1 \cap \cdots \cap E_m$ is the disjoint union of the sets of (1) for the patterns on $S$ whose restriction to each $S_i$ is admissible for $E_i$; there are $n_1 \cdots n_m$ of them, so its
measure is $n_1 \cdots n_m2^{-|S|} = \prod_in_i2^{-|S_i|}$. A subfamily of the $E_k$ is again determined by disjoint sets of positions, which gives independence.

(3) Each $g_i$ is a finite linear combination of indicator functions of sets determined by $S_i$; expand the product and apply (2) to each term.
\end{proof}

\begin{theorem}[Borel--Cantelli lemmas]\label{measure-theory:theorem-borel-cantelli}
Let $A_1, A_2, \ldots$ be events in a probability space.
\begin{enumerate}
\item If $\sum_kP(A_k) < \infty$, then $P(\limsup_kA_k) = 0$: almost surely only finitely many $A_k$ happen.
\item If the $A_k$ are independent and $\sum_kP(A_k) = \infty$, then $P(\limsup_kA_k) = 1$: almost surely infinitely many $A_k$ happen.
\end{enumerate}
\end{theorem}

\begin{proof}
(1) For every $n$, $\limsup_kA_k \subseteq \bigcup_{k \geq n}A_k$, so $P(\limsup_kA_k) \leq \sum_{k \geq n}P(A_k)$, the tail of a convergent series, which tends to $0$ as $n \to \infty$.

(2) We show that the complement $\bigcup_n\bigcap_{k \geq n}A_k^c$ is null; it suffices to show $P(\bigcap_{k \geq n}A_k^c) = 0$ for each $n$. For $N \geq n$, the events $A_k^c$ are independent
(Lemma~\ref{measure-theory:lemma-independent-complements}), and $1 - t \leq e^{-t}$ for real $t$ (the exponential lies above its tangent line at $0$), so
\[
P\Bigl(\bigcap_{k = n}^NA_k^c\Bigr) = \prod_{k = n}^N\bigl(1 - P(A_k)\bigr) \leq \exp\Bigl(-\sum_{k = n}^NP(A_k)\Bigr) \longrightarrow 0 \qquad (N \to \infty).
\]
These sets decrease in $N$ to $\bigcap_{k \geq n}A_k^c$, so by continuity of measures (Proposition~\ref{measure-theory:proposition-measure-properties}) its probability is $0$.
\end{proof}

For independent events $\limsup_kA_k$ therefore has probability $0$ or $1$, never anything in between: a first instance of a \emph{zero-one law}. The first lemma was already used, without the name, in the
proof of Proposition~\ref{measure-theory:proposition-modes}(2).

\begin{counterexample}\label{measure-theory:counterexample-borel-cantelli}
Independence cannot be dropped in (2). On $[0, 1)$ let $A_k = [0, 1/k)$. Then $\sum_k\lambda(A_k) = \sum_k1/k = \infty$, but the sets decrease, so
$\limsup_kA_k = \bigcap_kA_k = \{0\}$, a null set. The events are far from independent: $\lambda(A_j \cap A_k) = 1/\max(j, k)$, not $1/(jk)$. Even more simply, $A_k = A$ for all $k$ with $\lambda(A) = \frac12$ gives
$\limsup_kA_k = A$, of probability $\frac12$.
\end{counterexample}

\begin{example}\label{measure-theory:example-infinite-monkey}
\emph{Every pattern appears infinitely often.} Fix a block $\beta = (\beta_1, \ldots, \beta_m)$ of binary digits. For $j \geq 0$ let $B_j$ be the event that the digits at positions
$jm + 1, \ldots, jm + m$ form the block $\beta$. These events are determined by the disjoint sets of positions $\{jm + 1, \ldots, jm + m\}$, so they are independent, each of probability $2^{-m}$
(Lemma~\ref{measure-theory:lemma-binary-digits}). Since $\sum_j2^{-m} = \infty$, almost every $x \in [0, 1)$ lies in infinitely many $B_j$. There are countably many blocks, so for almost every $x$ every
finite block of digits occurs infinitely often in the binary expansion of $x$. (Encoding letters by blocks of digits, this is the ``infinite monkey theorem'': a random infinite text contains every
book infinitely often.)
\end{example}

The probability of a long block is small, and this limits how long the runs of $1$s in a random expansion can be. Let $R_n(x)$ be the length of the run of $1$s starting at position $n$: the largest $r \geq 0$ with
$d_n(x) = \cdots = d_{n+r-1}(x) = 1$, which is finite because the expansion does not end with infinitely many $1$s.

\begin{proposition}[Long runs]\label{measure-theory:proposition-long-runs}
For almost every $x \in [0, 1)$, $\displaystyle\limsup_{n \to \infty}\frac{R_n(x)}{\log_2n} = 1$.
\end{proposition}

\begin{proof}
\emph{Upper bound.} For $r > 0$, $R_n \geq r$ means that the $\lceil r\rceil$ digits from position $n$ on are $1$, so $\lambda(R_n \geq r) = 2^{-\lceil r\rceil} \leq 2^{-r}$ by
Lemma~\ref{measure-theory:lemma-binary-digits}(1). For $\varepsilon > 0$ and $r = (1 + \varepsilon)\log_2n$ this is at most $n^{-1-\varepsilon}$, which is summable. By
Theorem~\ref{measure-theory:theorem-borel-cantelli}(1), for almost every $x$ we have $R_n(x) < (1 + \varepsilon)\log_2n$ for all large $n$, so the $\limsup$ is at most $1 + \varepsilon$. Let
$\varepsilon = 1/m$, $m = 1, 2, \ldots$, and discard countably many null sets.

\emph{Lower bound.} We need independent events, so we use disjoint windows. For $j \geq 2$ let $L_j = \lfloor\log_2j\rfloor \geq 1$, $s_2 = 1$ and $s_{j+1} = s_j + L_j$, so that the windows
$W_j = \{s_j, \ldots, s_j + L_j - 1\}$ are disjoint. Let $E_j$ be the event that all digits in $W_j$ are $1$; then $E_j \subseteq \{R_{s_j} \geq L_j\}$, and $\lambda(E_j) = 2^{-L_j} \geq 1/j$. The $E_j$ are
independent (Lemma~\ref{measure-theory:lemma-binary-digits}(2)) and $\sum_j1/j = \infty$, so by Theorem~\ref{measure-theory:theorem-borel-cantelli}(2) almost every $x$ lies in infinitely many $E_j$. Now
$s_j = 1 + \sum_{2 \leq i < j}L_i \leq j\log_2j$, so $\log_2s_j \leq \log_2j + \log_2\log_2j$, while $L_j \geq \log_2j - 1$. Hence for such $x$,
\[
\frac{R_{s_j}(x)}{\log_2s_j} \geq \frac{\log_2j - 1}{\log_2j + \log_2\log_2j}
\]
for infinitely many $j$, and the right side tends to $1$. So the $\limsup$ is at least $1$.
\end{proof}

So among the first $n$ tosses of a fair coin one should expect a run of heads of length about $\log_2n$, and no essentially longer one: about $10$ in a thousand tosses, $20$ in a million. This is a form
of the \emph{Erd\H{o}s--R\'enyi law of long runs}; people asked to write down ``random'' sequences typically avoid such long runs, which is a simple way to detect fake randomness.

The most famous application of the first lemma is Borel's theorem that almost every number is \emph{normal}: its digits appear with the frequencies one expects. We prove the simplest case.

\begin{theorem}[Borel's normal number theorem]\label{measure-theory:theorem-borel-normal}
For almost every $x \in [0, 1)$,
\[
\lim_{n \to \infty}\frac{d_1(x) + d_2(x) + \cdots + d_n(x)}{n} = \frac12 .
\]
\end{theorem}

\begin{proof}
Let $\rho_k = 2d_k - 1 \in \{-1, 1\}$ and $S_n = \rho_1 + \cdots + \rho_n$; the claim is $S_n/n \to 0$ almost everywhere.

\emph{Step 1: the fourth moment.} Expanding,
$\int_0^1S_n^4 = \sum_{i, j, k, l \leq n}\int_0^1\rho_i\rho_j\rho_k\rho_l$. Since $\rho^2 = 1$, a product $\rho_i\rho_j\rho_k\rho_l$ in which some index occurs an odd number of times reduces to a product
$\rho_a$ or $\rho_a\rho_b$ ($a \neq b$), whose integral is $\int\rho_a \cdot \int\rho_b = 0$ by Lemma~\ref{measure-theory:lemma-binary-digits}(3), because $\int_0^1\rho_a = \frac12 - \frac12 = 0$. The
remaining terms equal $1$: either all four indices are equal ($n$ terms), or they take two different values, each twice ($3n(n - 1)$ terms: choose which of the three pairings of the four places carry equal
indices, then the two values in order). So $\int_0^1S_n^4 = n + 3n(n - 1) \leq 3n^2$.

\emph{Step 2: Borel--Cantelli.} Let $\varepsilon > 0$. By Markov's inequality (Corollary~\ref{measure-theory:corollary-integral-properties}(4)) applied to $S_n^4$,
\[
\lambda\bigl(|S_n| \geq \varepsilon n\bigr) \leq \frac{1}{\varepsilon^4n^4}\int_0^1S_n^4 \leq \frac{3}{\varepsilon^4n^2},
\]
which is summable in $n$. By Theorem~\ref{measure-theory:theorem-borel-cantelli}(1), for almost every $x$ we have $|S_n(x)| < \varepsilon n$ for all large $n$. Taking $\varepsilon = 1/m$ for all $m$ and
discarding countably many null sets, $S_n/n \to 0$ almost everywhere. Since $d_k = (1 + \rho_k)/2$, the average of the first $n$ digits is $\frac12 + S_n/(2n) \to \frac12$.
\end{proof}

\begin{remark}\label{measure-theory:remark-normal-numbers}
The same argument, applied to the blocks of Example~\ref{measure-theory:example-infinite-monkey} and to expansions in any base $b \geq 2$, shows that almost every real number is \emph{normal in every base}:
in base $b$, each block of $m$ digits occurs with limiting frequency $b^{-m}$ \cite{borel-1909}. The theorem is the strong law of large numbers for fair coin tosses. It is a striking example of an
existence proof by measure: almost all numbers are normal, yet it is not known whether $\sqrt2$, $e$ or $\pi$ is normal in any base. Explicit normal numbers can be constructed, such as Champernowne's
number $0.123456789101112\ldots$, obtained by writing the positive integers one after another, which is normal in base $10$ \cite{champernowne-1933}.
\end{remark}

\begin{example}\label{measure-theory:example-diophantine}
\emph{Almost no number is very well approximable.} By Dirichlet's approximation theorem (Exercise~\ref{measure-theory:exercise-dirichlet-approximation}), every irrational $x$ has infinitely many
approximations $p/q$ with $|x - p/q| < 1/q^2$. The exponent $2$ cannot be improved for almost all $x$. Fix $\varepsilon > 0$ and for $q \geq 1$ let $A_q$ be the set of $x \in [0, 1]$ for which some integer
$p$ satisfies $|x - p/q| < q^{-2-\varepsilon}$. Only $p = -1, 0, \ldots, q + 1$ can occur, so $A_q$ is covered by $q + 3$ intervals of length $2q^{-2-\varepsilon}$, and
$\lambda(A_q) \leq 2(q + 3)q^{-2-\varepsilon} \leq 8q^{-1-\varepsilon}$, which is summable. By Theorem~\ref{measure-theory:theorem-borel-cantelli}(1), almost every $x$ lies in only finitely many $A_q$, and
hence has only finitely many approximations with $|x - p/q| < q^{-2-\varepsilon}$. In particular the \emph{Liouville numbers}, the irrational $x$ that for every $n$ have an approximation $p/q$ with $q \geq 2$ and
$|x - p/q| < q^{-n}$, form a null set. An example is $\sum_{k \geq 1}10^{-k!}$: cutting the series after the term $k = m$ gives $p/q$ with $q = 10^{m!}$ and error less than
$2 \cdot 10^{-(m+1)!} = 2q^{-(m+1)}$.
\end{example}

\section{Littlewood's principles and Lusin's theorem}\label{measure-theory:section-lusin}

Littlewood summarised the theory of Lebesgue measure on $\RR^n$ in three principles: every measurable set of finite measure is nearly a finite union of boxes;
every measurable function is nearly continuous; every convergent sequence of measurable functions is nearly uniformly convergent. The third is Egorov's
theorem (Theorem~\ref{measure-theory:theorem-egorov}); this section makes the first two precise. Throughout, $\lambda$ is Lebesgue measure on $\RR^n$ and
$E \mathbin{\triangle} A = (E \setminus A) \cup (A \setminus E)$ is the symmetric difference.

\begin{lemma}\label{measure-theory:lemma-dyadic-cubes}
Every open set $W \subseteq \RR^n$ is a countable union of closed \emph{dyadic cubes} $\prod_i[k_i2^{-j}, (k_i + 1)2^{-j}]$ ($j \in \NN$, $k_i \in \ZZ$) with pairwise
disjoint interiors, and $\lambda(W)$ is the sum of their volumes.
\end{lemma}

\begin{proof}
Two dyadic cubes of side at most $1$ are either nested or have disjoint interiors, and a dyadic cube of side at most $1$ is contained in only finitely many
dyadic cubes of side at most $1$. Every $x \in W$ lies in a dyadic cube of side at most $1$ contained in $W$, since $W$ contains a ball around $x$. Let
$(Q_k)$ be the dyadic cubes of side at most $1$ contained in $W$ that are maximal for inclusion among these; every dyadic cube of side at most $1$ in $W$ lies in
one of them, so they cover $W$, and distinct maximal cubes have disjoint interiors. There are countably many dyadic cubes. The boundary of a cube is a
finite union of boxes of volume $0$, a null set (Theorem~\ref{measure-theory:theorem-lebesgue-measure}), so the interiors, which are disjoint, have the same
measures as the cubes, and $\lambda(W) = \sum_k\lambda(Q_k)$ by countable additivity.
\end{proof}

\begin{proposition}[First principle]\label{measure-theory:proposition-approximation-boxes}
Let $E \subseteq \RR^n$ be Lebesgue measurable with $\lambda(E) < \infty$, and $\varepsilon > 0$. There is a finite union $A$ of boxes with
$\lambda(E \mathbin{\triangle} A) < \varepsilon$.
\end{proposition}

\begin{proof}
By Theorem~\ref{measure-theory:theorem-lebesgue-measure} there is an open $W \supseteq E$ with $\lambda(W \setminus E) < \varepsilon/2$. Write $W = \bigcup_kQ_k$ as in
Lemma~\ref{measure-theory:lemma-dyadic-cubes}. Since $\sum_k\lambda(Q_k) = \lambda(W) < \infty$, there is $N$ with $\sum_{k > N}\lambda(Q_k) < \varepsilon/2$. Let
$A = Q_1 \cup \cdots \cup Q_N$. Then $\lambda(A \setminus E) \leq \lambda(W \setminus E) < \varepsilon/2$ and
$\lambda(E \setminus A) \leq \lambda(W \setminus A) \leq \sum_{k > N}\lambda(Q_k) < \varepsilon/2$.
\end{proof}

\begin{theorem}[Lusin]\label{measure-theory:theorem-lusin}
Let $E \subseteq \RR^n$ be Lebesgue measurable with $\lambda(E) < \infty$, let $f \colon E \to \CC$ be measurable, and $\varepsilon > 0$. There is a compact set
$K \subseteq E$ with $\lambda(E \setminus K) < \varepsilon$ such that the restriction $f|_K$ is continuous. Moreover there is $g \in C_c(\RR^n)$ with $g = f$ on $K$,
and if $|f| \leq c$ on $E$ we can take $|g| \leq c$.
\end{theorem}

\begin{proof}
\emph{Step 1: indicator functions.} Let $f = \mathbf{1}_A$ with $A \subseteq E$ measurable. By inner regularity
(Theorem~\ref{measure-theory:theorem-lebesgue-measure}) there are compact sets $K_1 \subseteq A$ and $K_2 \subseteq E \setminus A$ with
$\lambda(A \setminus K_1) < \varepsilon/2$ and $\lambda((E \setminus A) \setminus K_2) < \varepsilon/2$. Then $K = K_1 \cup K_2$ is compact with $\lambda(E \setminus K) < \varepsilon$.
As $K_1$ and $K_2$ are disjoint and closed, each is open in $K$, and $f|_K$ is constant on each, hence continuous.

\emph{Step 2: simple functions.} If $f = \sum_{i=1}^ma_i\mathbf{1}_{A_i}$, apply Step 1 to each $A_i$ with $\varepsilon/m$ and intersect the compact sets.

\emph{Step 3: real functions.} By Proposition~\ref{measure-theory:proposition-measurable-functions}(4), applied to $f^+$ and $f^-$, there are simple functions
$s_k$ with $s_k \to f$ pointwise on $E$. By Step 2 there are compact $K_k \subseteq E$ with $\lambda(E \setminus K_k) < \varepsilon2^{-k-2}$ and $s_k|_{K_k}$
continuous. By Egorov's theorem (Theorem~\ref{measure-theory:theorem-egorov}, on the finite measure space $E$) there is a measurable $F \subseteq E$ with
$\lambda(E \setminus F) < \varepsilon/4$ on which $s_k \to f$ uniformly, and by inner regularity a compact $K_0 \subseteq F$ with $\lambda(F \setminus K_0) < \varepsilon/4$.
The compact set $K = K_0 \cap \bigcap_kK_k$ satisfies $\lambda(E \setminus K) < \varepsilon/4 + \varepsilon/4 + \sum_k\varepsilon2^{-k-2} \leq \varepsilon$, and on it the
continuous functions $s_k|_K$ converge uniformly to $f|_K$, which is therefore continuous (Proposition~\ref{general-topology:proposition-uniform-limit}). For
complex $f$, treat the real and imaginary parts with $\varepsilon/2$ each and intersect.

\emph{Extension.} The space $\RR^n$ is metric, hence normal (Proposition~\ref{general-topology:proposition-normal-examples}), and $K$ is closed, so by the Tietze
extension theorem (Theorem~\ref{general-topology:theorem-tietze}) the real and imaginary parts of $f|_K$ extend to continuous functions on $\RR^n$; let $h$ be the
resulting extension of $f|_K$. If $|f| \leq c$, replace $h$ by $r \circ h$, where $r(z) = z$ for $|z| \leq c$ and $r(z) = cz/|z|$ otherwise is a continuous retraction
of $\CC$ onto the closed disc of radius $c$. Finally choose $\phi \in C_c(\RR^n)$ with values in $[0, 1]$ and $\phi = 1$ on $K$
(Lemma~\ref{general-topology:lemma-lch-shrinking}), and let $g = \phi h$.
\end{proof}

\begin{example}\label{measure-theory:example-lusin-dirichlet}
Lusin's theorem says that the \emph{restriction} of $f$ to $K$ is continuous, not that $f$ is continuous at the points of $K$. The Dirichlet function
$\mathbf{1}_\QQ$ on $[0, 1]$ is discontinuous at every point. Enumerate $\QQ \cap [0, 1] = \{q_1, q_2, \ldots\}$ and let $U$ be the union of the intervals
$(q_k - \varepsilon2^{-k-2}, q_k + \varepsilon2^{-k-2})$. Then $\lambda(U) \leq \varepsilon/2$, the set $K = [0, 1] \setminus U$ is compact with
$\lambda([0, 1] \setminus K) < \varepsilon$, and $\mathbf{1}_\QQ$ vanishes on $K$, so its restriction to $K$ is continuous.
\end{example}

\begin{corollary}\label{measure-theory:corollary-lusin-limit}
Every measurable function $f \colon \RR^n \to \CC$ is the almost everywhere limit of a sequence of continuous functions.
\end{corollary}

\begin{proof}
For $k \geq 1$ apply Theorem~\ref{measure-theory:theorem-lusin} to the restriction of $f$ to the ball $B(0, k)$ with $\varepsilon = 2^{-k}$, to get
$g_k \in C_c(\RR^n)$ equal to $f$ on $B(0, k)$ outside a set $N_k$ of measure less than $2^{-k}$. By monotone convergence,
$\int\sum_k\mathbf{1}_{N_k} = \sum_k\lambda(N_k) < \infty$, so $\sum_k\mathbf{1}_{N_k} < \infty$ almost everywhere: almost every $x$ lies in only finitely many $N_k$.
For such $x$ and all large $k > \|x\|$, $g_k(x) = f(x)$.
\end{proof}

\section{Product measures and Fubini's theorem}\label{measure-theory:section-fubini}

The area of a rectangle is the product of the lengths of its sides, and the volume of a solid is the integral of the areas of its slices. The first statement defines a measure on a product of
measure spaces; the second, Cavalieri's principle, is a theorem about it, and its generalisation to functions is Fubini's theorem: a double integral can be computed as an iterated integral, in
either order. Both statements need hypotheses, and the counterexamples below show which.

\subsection{The product $\sigma$-algebra}\label{measure-theory:subsection-product-sigma-algebra}

\begin{definition}\label{measure-theory:definition-product-measure}
For measurable spaces $(X, \mathcal{A})$ and $(Y, \mathcal{B})$, the \emph{product $\sigma$-algebra} $\mathcal{A} \otimes \mathcal{B}$ on $X \times Y$ is generated by the \emph{rectangles} $A \times B$. For
$E \subseteq X \times Y$, the sections are $E_x = \{y : (x, y) \in E\}$ and $E^y = \{x : (x, y) \in E\}$.
\end{definition}

The rectangles form a $\pi$-system, since $(A \times B) \cap (A^{\prime} \times B^{\prime}) = (A \cap A^{\prime}) \times (B \cap B^{\prime})$, and the finite disjoint unions of rectangles form an algebra:
the complement of $A \times B$ is $((X \setminus A) \times Y) \sqcup (A \times (Y \setminus B))$. The product $\sigma$-algebra is the smallest $\sigma$-algebra for which the two projections
$X \times Y \to X$ and $X \times Y \to Y$ are measurable, since $A \times B = \pi_X^{-1}(A) \cap \pi_Y^{-1}(B)$. Products of finitely many measurable spaces are defined in the same way, and are
associative: $(\mathcal{A}_1 \otimes \mathcal{A}_2) \otimes \mathcal{A}_3 = \mathcal{A}_1 \otimes (\mathcal{A}_2 \otimes \mathcal{A}_3) = \mathcal{A}_1 \otimes \mathcal{A}_2 \otimes \mathcal{A}_3$, as each is generated by
the boxes $A_1 \times A_2 \times A_3$ (a set $E \times A_3$ with $E \in \mathcal{A}_1 \otimes \mathcal{A}_2$ lies in the $\sigma$-algebra generated by boxes, because the sets $E$ with this property form a
$\sigma$-algebra containing the rectangles).

\begin{lemma}\label{measure-theory:lemma-sections}
Let $E \in \mathcal{A} \otimes \mathcal{B}$ and let $f \colon X \times Y \to [-\infty, \infty]$ or $\CC$ be $\mathcal{A} \otimes \mathcal{B}$-measurable. For every $x \in X$ and $y \in Y$ the sections
$E_x$ and $E^y$ are measurable, and so are the functions $f(x, -) \colon Y \to [-\infty, \infty]$ and $f(-, y) \colon X \to [-\infty, \infty]$.
\end{lemma}

\begin{proof}
Fix $x$. The sets $E \subseteq X \times Y$ with $E_x \in \mathcal{B}$ form a $\sigma$-algebra, since $(X \times Y \setminus E)_x = Y \setminus E_x$ and $(\bigcup_kE_k)_x = \bigcup_k(E_k)_x$; it contains the
rectangles, as $(A \times B)_x$ is $B$ or $\emptyset$. So it contains $\mathcal{A} \otimes \mathcal{B}$. For $f$, note $\{y : f(x, y) > a\} = \{f > a\}_x$. The other variable is symmetric.
\end{proof}

\begin{proposition}\label{measure-theory:proposition-borel-product}
For topological spaces $X$, $Y$, $\mathcal{B}(X) \otimes \mathcal{B}(Y) \subseteq \mathcal{B}(X \times Y)$, with equality if $X$ and $Y$ have countable bases of open sets. In particular
$\mathcal{B}(\RR^m) \otimes \mathcal{B}(\RR^n) = \mathcal{B}(\RR^{m+n})$.
\end{proposition}

\begin{proof}
The projections are continuous, hence Borel (Lemma~\ref{measure-theory:lemma-measurable-generators}); so $A \times B = \pi_X^{-1}(A) \cap \pi_Y^{-1}(B)$ is Borel for Borel $A$, $B$, and the rectangles,
hence their $\sigma$-algebra, lie in $\mathcal{B}(X \times Y)$. If $(U_i)$ and $(V_j)$ are countable bases, the products $U_i \times V_j$ form a countable base of $X \times Y$, so every open subset of
$X \times Y$ is a countable union of rectangles $U_i \times V_j$, and lies in $\mathcal{B}(X) \otimes \mathcal{B}(Y)$. For $\RR^n$ take the open boxes with rational vertices.
\end{proof}

\subsection{Product measures}\label{measure-theory:subsection-product-measures}

\begin{theorem}\label{measure-theory:theorem-product-measure}
Let $\mu$ and $\nu$ be $\sigma$-finite measures on $X$ and $Y$. For $E \in \mathcal{A} \otimes \mathcal{B}$, the sections are measurable, the functions $x \mapsto \nu(E_x)$ and $y \mapsto \mu(E^y)$ are measurable, and
\[
(\mu \times \nu)(E) = \int\nu(E_x)\,d\mu(x) = \int\mu(E^y)\,d\nu(y)
\]
defines the unique measure on $\mathcal{A} \otimes \mathcal{B}$ with $(\mu \times \nu)(A \times B) = \mu(A)\nu(B)$.
\end{theorem}

\begin{proof}
Measurability of the sections is Lemma~\ref{measure-theory:lemma-sections}.

\emph{Step 1: finite measures.} Assume $\mu(X), \nu(Y) < \infty$, and let $\mathcal{L}$ be the family of $E \in \mathcal{A} \otimes \mathcal{B}$ such that $x \mapsto \nu(E_x)$ and $y \mapsto \mu(E^y)$ are
measurable and $\int\nu(E_x)\,d\mu = \int\mu(E^y)\,d\nu$. A rectangle lies in $\mathcal{L}$: $\nu((A \times B)_x) = \nu(B)\mathbf{1}_A(x)$ and $\mu((A \times B)^y) = \mu(A)\mathbf{1}_B(y)$, and both
integrals equal $\mu(A)\nu(B)$. In particular $X \times Y \in \mathcal{L}$. If $E \in \mathcal{L}$, then $\nu((X \times Y \setminus E)_x) = \nu(Y) - \nu(E_x)$, a measurable function, and the two integrals
are $\mu(X)\nu(Y)$ minus the (finite) integrals for $E$; so the complement lies in $\mathcal{L}$. If $E = \bigsqcup_kE_k$ with $E_k \in \mathcal{L}$, then $\nu(E_x) = \sum_k\nu((E_k)_x)$ is measurable,
and by monotone convergence (Corollary~\ref{measure-theory:corollary-integral-properties}(1)) both integrals are the sums of the integrals for the $E_k$. So $\mathcal{L}$ is a Dynkin system
containing the $\pi$-system of rectangles, and by Lemma~\ref{measure-theory:lemma-dynkin} it is all of $\mathcal{A} \otimes \mathcal{B}$.

\emph{Step 2: the formula defines a measure.} With $(\mu \times \nu)(E) = \int\nu(E_x)\,d\mu$, the empty set has measure $0$, and countable additivity is the computation with $\sum_k$ just made.

\emph{Step 3: $\sigma$-finite measures.} Choose $X_m \uparrow X$ and $Y_m \uparrow Y$ measurable with $\mu(X_m), \nu(Y_m) < \infty$. Apply Step 1 to the finite measures $\mu_m = \mu_{X_m}$ and
$\nu_m = \nu_{Y_m}$ (restrictions, Example~\ref{measure-theory:example-measures}(3)): $x \mapsto \nu(E_x \cap Y_m)\mathbf{1}_{X_m}(x)$ and $y \mapsto \mu(E^y \cap X_m)\mathbf{1}_{Y_m}(y)$ are measurable, and their
integrals against $\mu$ and $\nu$ agree. As $m \to \infty$ these functions increase to $\nu(E_x)$ and $\mu(E^y)$ (continuity from below), which are therefore measurable, and monotone convergence gives
the equality of the integrals. Countable additivity follows as in Step 2.

\emph{Step 4: uniqueness.} Two measures with the stated values on rectangles agree on the $\pi$-system of rectangles, and are finite on the rectangles $X_m \times Y_m$ increasing to $X \times Y$; by
Lemma~\ref{measure-theory:lemma-dynkin} they agree on $\mathcal{A} \otimes \mathcal{B}$.
\end{proof}

Product measures of finitely many $\sigma$-finite measures are defined inductively, $\mu_1 \times \cdots \times \mu_n = (\mu_1 \times \cdots \times \mu_{n-1}) \times \mu_n$; by uniqueness, this is the only
measure with $(\mu_1 \times \cdots \times \mu_n)(A_1 \times \cdots \times A_n) = \prod_i\mu_i(A_i)$, so the bracketing does not matter.

\begin{example}\label{measure-theory:example-product-measures}
\begin{enumerate}
\item For counting measures on countable sets $X$ and $Y$, $\mu \times \nu$ is counting measure on $X \times Y$: both are measures on $\mathcal{P}(X \times Y)$ (every subset is a countable union of
points, which are rectangles) and agree on points.
\item For Dirac measures, $\delta_a \times \delta_b = \delta_{(a, b)}$. For a fair die, the product of two copies of the uniform measure on $\{1, \ldots, 6\}$ is the uniform measure on the $36$ pairs:
the model of two independent throws.
\item For Lebesgue measure, $\lambda^m \times \lambda^n$ and $\lambda^{m+n}$ agree on the Borel sets (Remark~\ref{measure-theory:remark-lebesgue-product} below). The area of a disc of radius $r$ is
$\int_{-r}^r2\sqrt{r^2 - x^2}\,dx = \pi r^2$, the integral of the lengths of its vertical sections.
\end{enumerate}
\end{example}

\begin{counterexample}\label{measure-theory:counterexample-product-sigma-finite}
\emph{$\sigma$-finiteness is needed.} Let $\mu = \lambda$ on $X = [0, 1]$ with the Borel sets, $\nu$ counting measure on $Y = [0, 1]$ with the Borel sets, and $E = \{(x, x)\}$ the diagonal, which is
closed, hence in $\mathcal{B}([0, 1]) \otimes \mathcal{B}([0, 1]) = \mathcal{B}([0, 1]^2)$ (Proposition~\ref{measure-theory:proposition-borel-product}). Every section is a point, so
\[
\int\nu(E_x)\,d\mu(x) = \int_0^11\,dx = 1, \qquad \int\mu(E^y)\,d\nu(y) = \int0\,d\nu = 0.
\]
The counting measure on $[0, 1]$ is not $\sigma$-finite. The theory of products of measures that are not $\sigma$-finite is delicate (the product property on rectangles no longer determines a
measure), and we avoid it.
\end{counterexample}

\subsection{The theorems of Tonelli and Fubini}\label{measure-theory:subsection-tonelli-fubini}

\begin{theorem}[Tonelli and Fubini]\label{measure-theory:theorem-fubini}
Let $\mu$, $\nu$ be $\sigma$-finite.
\begin{enumerate}
\item If $f \colon X \times Y \to [0, \infty]$ is measurable, then $x \mapsto \int f(x, y)\,d\nu(y)$ is measurable and
$\int_{X \times Y}f\,d(\mu \times \nu) = \int\left(\int f(x, y)\,d\nu(y)\right)d\mu(x) = \int\left(\int f(x, y)\,d\mu(x)\right)d\nu(y)$.
\item If $f \in L^1(\mu \times \nu)$, then $f(x, -) \in L^1(\nu)$ for a.e. $x$, the function $x \mapsto \int f(x, y)\,d\nu(y)$ (defined a.e.) is integrable, and the same equalities hold.
\end{enumerate}
\end{theorem}

\begin{proof}
\emph{Step 1: indicators.} For $f = \mathbf{1}_E$, $\int f(x, y)\,d\nu(y) = \nu(E_x)$, and (1) is Theorem~\ref{measure-theory:theorem-product-measure}.

\emph{Step 2: simple functions.} Each of the three expressions is additive and positively homogeneous in $f \geq 0$ (Corollary~\ref{measure-theory:corollary-integral-properties}(1)), and a sum of
measurable functions is measurable; so (1) holds for nonnegative simple functions.

\emph{Step 3: Tonelli.} For measurable $f \geq 0$ choose simple $0 \leq s_k \uparrow f$ (Proposition~\ref{measure-theory:proposition-measurable-functions}(4)). For each $x$, $s_k(x, -) \uparrow f(x, -)$, so by
monotone convergence $\int s_k(x, y)\,d\nu(y) \uparrow \int f(x, y)\,d\nu(y)$; the limit of measurable functions is measurable, and monotone convergence on $X$ and on $X \times Y$ passes the equalities of
Step 2 to the limit. The other order is symmetric.

\emph{Step 4: Fubini.} Let $f \in L^1(\mu \times \nu)$. By (1) applied to $|f|$, $\int(\int|f(x, y)|\,d\nu(y))\,d\mu(x) = \int|f| < \infty$, so the inner integral is finite for $x$ outside a null set $N$
(Corollary~\ref{measure-theory:corollary-integral-properties}(4)), and $f(x, -) \in L^1(\nu)$ there. For real $f$ write $f = f^+ - f^-$ with $f^\pm \leq |f|$ integrable; for $x \notin N$,
$\int f(x, y)\,d\nu = \int f^+(x, y)\,d\nu - \int f^-(x, y)\,d\nu$, a difference of two integrable functions of $x$ (by (1) for $f^\pm$), and integrating over $X$ gives
$\int f^+\,d(\mu \times \nu) - \int f^-\,d(\mu \times \nu) = \int f\,d(\mu \times \nu)$. For complex $f$ treat the real and imaginary parts.
\end{proof}

In practice one uses the two parts together: to justify exchanging the order of integration for a function $f$ of variable sign, first compute an iterated integral of $|f|$ with Tonelli's theorem,
in whichever order is convenient; if it is finite, $f$ is integrable and Fubini's theorem applies in both orders.

\begin{counterexample}\label{measure-theory:counterexample-fubini}
Integrability is needed in (2). On $\NN \times \NN$ with counting measures let $f(n, n) = 1$, $f(n, n+1) = -1$ and $f = 0$ otherwise. Then $\sum_m\sum_nf(n, m) = 1$, since only the first column $m = \min\NN$ has nonzero
sum, whereas $\sum_n\sum_mf(n, m) = 0$, since every row sums to $0$.
\end{counterexample}

\begin{counterexample}\label{measure-theory:counterexample-fubini-continuous}
The same happens for a function continuous outside one point of the unit square. Let $f(x, y) = \dfrac{x^2 - y^2}{(x^2 + y^2)^2}$ on $(0, 1]^2$. Since
$\dfrac{\partial}{\partial y}\dfrac{y}{x^2 + y^2} = \dfrac{x^2 - y^2}{(x^2 + y^2)^2}$, for fixed $x > 0$
\[
\int_0^1f(x, y)\,dy = \frac{1}{1 + x^2}, \qquad \int_0^1\int_0^1f(x, y)\,dy\,dx = \frac\pi4,
\]
and by the antisymmetry $f(y, x) = -f(x, y)$ the other order gives $-\pi/4$. Tonelli's theorem therefore says that $\int\!\!\int|f| = \infty$; directly, in polar coordinates $|f| = |\cos2\theta|/r^2$, whose
integral over a quarter disc diverges like $\int_0r^{-1}\,dr$.
\end{counterexample}

\begin{remark}\label{measure-theory:remark-fubini-measurability}
\emph{Measurability on the product is needed.} Tonelli's theorem assumes that $f$ is measurable for $\mathcal{A} \otimes \mathcal{B}$, not only that all sections are measurable. Assuming the
continuum hypothesis, Sierpi\'nski constructed a set $E \subseteq [0, 1]^2$ all of whose vertical sections have countable complement and all of whose horizontal sections are countable; then
$\int\lambda(E_x)\,dx = 1$ and $\int\lambda(E^y)\,dy = 0$, so $E$ is not measurable although all its sections are \cite[Chapter 8]{rudin-rca}. The construction uses a well-ordering of $[0, 1]$ in which
every initial segment is countable: $E = \{(x, y) : x \prec y\}$.
\end{remark}

\begin{remark}\label{measure-theory:remark-lebesgue-product}
$\lambda^{m+n}$ and $\lambda^m \times \lambda^n$ agree on boxes, hence on Borel sets by Lemma~\ref{measure-theory:lemma-dynkin}, and $\lambda^{m+n}$ is the completion of $\lambda^m \times \lambda^n$. So
Fubini's theorem applies to Lebesgue measure on $\RR^{m+n} = \RR^m \times \RR^n$ for Borel functions, and for Lebesgue measurable functions with ``for almost every $x$'' qualifications.
The product of the complete measures $\lambda^m$ and $\lambda^n$, on the $\sigma$-algebra generated by products of Lebesgue measurable sets, is not complete: if $N \subseteq \RR^m$ is not Lebesgue
measurable, $N \times \{0\}$ is contained in the null set $\RR^m \times \{0\}$, but it is not measurable for the product, since its section at $y = 0$ is $N$ (Lemma~\ref{measure-theory:lemma-sections}).
\end{remark}

\begin{proposition}[Layer cake formula]\label{measure-theory:proposition-layer-cake}
Let $\mu$ be a $\sigma$-finite measure on $X$ and $f \colon X \to [0, \infty]$ measurable. Then
\[
\int_Xf\,d\mu = \int_0^\infty\mu(\{f > t\})\,dt.
\]
More generally, if $\psi \colon [0, \infty) \to [0, \infty)$ is Borel measurable and $\phi(s) = \int_0^s\psi(t)\,dt$ (with $\phi(\infty) = \int_0^\infty\psi$), then
$\int_X\phi \circ f\,d\mu = \int_0^\infty\psi(t)\,\mu(\{f > t\})\,dt$; in particular $\int_Xf^p\,d\mu = p\int_0^\infty t^{p-1}\mu(\{f > t\})\,dt$ for $p > 0$.
\end{proposition}

\begin{proof}
The set $E = \{(x, t) \in X \times [0, \infty) : t < f(x)\}$ is measurable for the product $\sigma$-algebra, because $t < f(x)$ holds exactly when $t < q < f(x)$ for some rational
$q$, so $E = \bigcup_{q \in \QQ_{>0}}\{f > q\} \times [0, q)$. Apply Tonelli's theorem (Theorem~\ref{measure-theory:theorem-fubini}(1)) to the non-negative measurable function
$(x, t) \mapsto \psi(t)\mathbf{1}_E(x, t)$ on $X \times [0, \infty)$ with $\mu \times \lambda$. Integrating first in $t$ gives $\int_0^{f(x)}\psi(t)\,dt = \phi(f(x))$, integrating first in
$x$ gives $\psi(t)\mu(\{f > t\})$. The first formula is the case $\psi = 1$, and the last one the case $\psi(t) = pt^{p-1}$.
\end{proof}

The name describes the picture: the region under the graph of $f$ is cut into horizontal layers $\{f > t\}$ of thickness $dt$.

\begin{example}\label{measure-theory:example-layer-cake}
\begin{enumerate}
\item \emph{Expectation from the tail.} For a random variable $f \geq 0$ on a probability space, $\int f\,dP = \int_0^\infty P(f > t)\,dt$. If $f$ takes values in $\{0, 1, 2, \ldots\}$, then
$P(f > t) = P(f > n)$ for $n \leq t < n + 1$, and $\int f\,dP = \sum_{n \geq 0}P(f > n)$. For the number $f$ of throws of a fair die up to and including the first six,
$P(f > n) = (5/6)^n$, and $\int f\,dP = \sum_n(5/6)^n = 6$.
\item \emph{Integrability of powers.} On $\RR^n$ and for $a > 0$, $\int_{\{|x| \leq 1\}}|x|^{-a}\,dx < \infty$ if and only if $a < n$: here $\lambda(\{|x|^{-a} > t\} \cap \{|x| \leq 1\})$ equals $\omega_nt^{-n/a}$ for $t \geq 1$, with
$\omega_n$ the volume of the unit ball, and $\int_1^\infty t^{-n/a}\,dt < \infty$ exactly when $n/a > 1$.
\end{enumerate}
\end{example}

\subsection{Infinite products of probability measures}\label{measure-theory:subsection-infinite-products}

Infinitely many independent throws of a die, or tosses of a coin, are described by a probability measure on an infinite product, and Fubini's theorem is the tool to construct it. Let
$(X_k, \mathcal{A}_k, \mu_k)$, $k \geq 1$, be probability spaces, $X = \prod_kX_k$, and $\pi_k \colon X \to X_k$ the projections. For $n \geq 1$ write $X^{(n)} = X_1 \times \cdots \times X_n$, with the product
$\sigma$-algebra $\mathcal{A}^{(n)}$ and the product measure $\mu^{(n)} = \mu_1 \times \cdots \times \mu_n$, and $p_n \colon X \to X^{(n)}$ for the projection onto the first $n$ coordinates.

\begin{definition}\label{measure-theory:definition-cylinder}
A \emph{cylinder} in $X$ is a set $C = p_n^{-1}(B)$ with $n \geq 1$ and $B \in \mathcal{A}^{(n)}$: whether $x \in C$ depends only on $x_1, \ldots, x_n$. The \emph{product $\sigma$-algebra}
$\bigotimes_k\mathcal{A}_k$ is the $\sigma$-algebra generated by the cylinders, equivalently the smallest one for which all $\pi_k$ are measurable.
\end{definition}

A cylinder $p_n^{-1}(B)$ is also $p_{n+1}^{-1}(B \times X_{n+1})$, so two cylinders can always be written with the same $n$; hence the cylinders form an algebra $\mathcal{C}$. On a cylinder put
\[
\mu_0(p_n^{-1}(B)) = \mu^{(n)}(B).
\]
This does not depend on the representation: $\mu^{(n+1)}(B \times X_{n+1}) = \mu^{(n)}(B)\mu_{n+1}(X_{n+1}) = \mu^{(n)}(B)$, and two representations of the same cylinder become equal after enlarging the
smaller $n$ (if $p_n^{-1}(B) = p_n^{-1}(B^{\prime})$ then $B = B^{\prime}$, as $p_n$ is surjective). Writing two disjoint cylinders with the same $n$ shows that $\mu_0$ is finitely additive on $\mathcal{C}$,
with $\mu_0(X) = 1$.

\begin{theorem}[Infinite product measures]\label{measure-theory:theorem-infinite-product}
There is a unique probability measure $\mu = \bigotimes_k\mu_k$ on $\bigotimes_k\mathcal{A}_k$ such that $\mu(p_n^{-1}(B)) = \mu^{(n)}(B)$ for all $n$ and $B \in \mathcal{A}^{(n)}$; in particular
\[
\mu(\{x : x_1 \in A_1, \ldots, x_n \in A_n\}) = \mu_1(A_1) \cdots \mu_n(A_n) \qquad (A_k \in \mathcal{A}_k).
\]
\end{theorem}

\begin{proof}
By the Carath\'eodory--Hahn theorem (Theorem~\ref{measure-theory:theorem-extension}) it suffices to show that $\mu_0$ is a premeasure on $\mathcal{C}$; the extension is unique since $\mu_0$ is finite, and it is
determined by its values on boxes, since the boxes form a $\pi$-system generating $\bigotimes_k\mathcal{A}_k$ (Lemma~\ref{measure-theory:lemma-dynkin}).

\emph{Step 1: reduction to continuity at $\emptyset$.} Suppose that every decreasing sequence of cylinders $C_1 \supseteq C_2 \supseteq \cdots$ with $\mu_0(C_j) \geq \varepsilon > 0$ for all $j$ has a nonempty
intersection. Let $C = \bigsqcup_kC^{\prime}_k$ with $C, C^{\prime}_k \in \mathcal{C}$. Then $D_j = C \setminus (C^{\prime}_1 \cup \cdots \cup C^{\prime}_j)$ are decreasing cylinders with empty intersection, so
$\mu_0(D_j) \to 0$ (otherwise $\mu_0(D_j) \geq \varepsilon$ for all $j$, the sequence being nonincreasing), and by finite additivity $\mu_0(C) = \sum_{k \leq j}\mu_0(C^{\prime}_k) + \mu_0(D_j) \to \sum_k\mu_0(C^{\prime}_k)$.

\emph{Step 2: sections of cylinders.} For $a \in X_1$ and $C \subseteq X$ let $C_a = \{(x_2, x_3, \ldots) : (a, x_2, x_3, \ldots) \in C\}$, a subset of $X^{\prime} = \prod_{k \geq 2}X_k$. Let $\mu^{\prime}_0$ be the
function defined in the same way as $\mu_0$ on the cylinders of $X^{\prime}$, built from $\mu_2, \mu_3, \ldots$. If $C = p_n^{-1}(B)$, then $C_a$ is the cylinder of $X^{\prime}$ over the section
$B_a \subseteq X_2 \times \cdots \times X_n$, and $X^{(n)} = X_1 \times (X_2 \times \cdots \times X_n)$ with $\mu^{(n)} = \mu_1 \times (\mu_2 \times \cdots \times \mu_n)$. By
Theorem~\ref{measure-theory:theorem-product-measure}, $a \mapsto \mu^{\prime}_0(C_a)$ is measurable and
\[
\mu_0(C) = \int_{X_1}\mu^{\prime}_0(C_a)\,d\mu_1(a).
\]

\emph{Step 3: choosing the first coordinate.} Let $C_j$ be decreasing cylinders with $\mu_0(C_j) \geq \varepsilon$. Let $F_j = \{a \in X_1 : \mu^{\prime}_0((C_j)_a) \geq \varepsilon/2\}$, a decreasing sequence of
measurable sets. Splitting the integral of Step 2 over $F_j$ and its complement, and using $\mu^{\prime}_0 \leq 1$,
$\varepsilon \leq \mu_0(C_j) \leq \mu_1(F_j) + \varepsilon/2$, so $\mu_1(F_j) \geq \varepsilon/2$. By continuity from above, $\mu_1(\bigcap_jF_j) \geq \varepsilon/2 > 0$; choose $a_1 \in \bigcap_jF_j$.
Then the sections $(C_j)_{a_1}$ are decreasing cylinders of $X^{\prime}$ with $\mu^{\prime}_0((C_j)_{a_1}) \geq \varepsilon/2$ for all $j$.

\emph{Step 4: induction.} Applying Step 3 to $X^{\prime}$ and the sections $(C_j)_{a_1}$, with $\varepsilon/2$ in place of $\varepsilon$, gives $a_2 \in X_2$ such that the sets
$\{(x_3, \ldots) : (a_1, a_2, x_3, \ldots) \in C_j\}$ have measure $\geq \varepsilon/4$ for all $j$; and so on. This produces $a = (a_1, a_2, \ldots) \in X$ such that for every $n$ and $j$ the set
$\{(x_{n+1}, \ldots) : (a_1, \ldots, a_n, x_{n+1}, \ldots) \in C_j\}$ has positive measure, so is nonempty.

\emph{Step 5: $a$ lies in every $C_j$.} Write $C_j = p_n^{-1}(B)$. By Step 4 some point $(a_1, \ldots, a_n, x_{n+1}, \ldots)$ lies in $C_j$; since membership depends only on the first $n$ coordinates,
$(a_1, \ldots, a_n) \in B$ and $a \in C_j$. So $\bigcap_jC_j \neq \emptyset$, and Step 1 concludes.
\end{proof}

\begin{definition}\label{measure-theory:definition-independent-variables}
Measurable functions (random variables) $f_k \colon \Omega \to Y_k$, $k \geq 1$, on a probability space $(\Omega, P)$ are \emph{independent} if for all measurable $B_k \subseteq Y_k$ the events
$\{f_k \in B_k\}$ are independent (Definition~\ref{measure-theory:definition-independent}).
\end{definition}

\begin{corollary}\label{measure-theory:corollary-independent-sequences}
For probability measures $\nu_k$ on measurable spaces $Y_k$ there is a probability space carrying independent random variables $f_k$ with distributions $\nu_k$: on $\prod_kY_k$ with
$\bigotimes_k\nu_k$, take the coordinate maps $f_k = \pi_k$.
\end{corollary}

\begin{proof}
$\pi_k$ is measurable and $(\pi_k)_*(\bigotimes\nu_j)(B) = \nu_k(B)$, by the theorem applied to the box with $A_k = B$ and all other $A_j = Y_j$. For distinct $k_1 < \cdots < k_m$ the event
$\{\pi_{k_1} \in B_1\} \cap \cdots \cap \{\pi_{k_m} \in B_m\}$ is a box, of measure $\prod_i\nu_{k_i}(B_i)$.
\end{proof}

\begin{example}\label{measure-theory:example-coin-tossing}
\emph{Coin tossing and Lebesgue measure.} Let $X = \{0, 1\}^{\NN_{\geq 1}}$ with $\mu = \bigotimes_k(\frac12\delta_0 + \frac12\delta_1)$, the model of infinitely many tosses of a fair coin, and
$\beta(x) = \sum_kx_k2^{-k} \in [0, 1]$. The map $\beta$ is measurable (a pointwise limit of functions depending on finitely many coordinates), and $\beta_*\mu = \lambda$ on $[0, 1]$. Indeed, for a dyadic
interval $I = [N2^{-K}, (N + 1)2^{-K}]$, $\beta^{-1}(I)$ consists of the sequences beginning with the $K$ binary digits of $N$, together with at most two sequences ending in all $0$s or all $1$s from
the $(K+1)$-st place on, which form a null set; so $\beta_*\mu(I) = 2^{-K} = \lambda(I)$. An interval $[0, t]$ with dyadic $t$ is a finite union of such intervals overlapping in finitely many points, so
again $\beta_*\mu([0, t]) = t$ (single points are null for both measures). These intervals form a $\pi$-system generating $\mathcal{B}([0, 1])$, and
Lemma~\ref{measure-theory:lemma-dynkin} gives $\beta_*\mu = \lambda$. Conversely, the binary digit maps $d_k$ of Lemma~\ref{measure-theory:lemma-binary-digits} send $\lambda$ to $\mu$. Thus ``a point of
$[0, 1]$ chosen uniformly at random'' and ``an infinite sequence of fair coin tosses'' are the same probability space, up to null sets; this is why Borel's normal number theorem
(Theorem~\ref{measure-theory:theorem-borel-normal}) is a theorem about coin tossing.
\end{example}

\begin{remark}\label{measure-theory:remark-kolmogorov-extension}
The theorem constructs only \emph{independent} sequences. Dependent ones, such as Markov chains or Brownian motion, are built from a consistent family of finite-dimensional distributions by
\emph{Kolmogorov's extension theorem}, whose proof is the same as above once the compactness used in Steps 3--5 is supplied by inner regularity of measures on $\RR^n$ (Part~\ref{part-probability}). For
products of uncountably many probability spaces the same construction works, since every set of the product $\sigma$-algebra depends on countably many coordinates.
\end{remark}

\section{Change of variables}\label{measure-theory:section-change-of-variables}

\begin{lemma}\label{measure-theory:lemma-linear-change}
For $T \in \GL_n(\RR)$ and Borel $E \subseteq \RR^n$, $\lambda(T(E)) = |\det T|\,\lambda(E)$.
\end{lemma}

\begin{proof}
Both sides are multiplicative in $T$, and every invertible matrix is a product of \emph{elementary} matrices (by Gaussian elimination): permutations of coordinates, scalings of one coordinate by
$c \neq 0$, and shears $x_i \mapsto x_i + cx_j$ ($i \neq j$). Permutations and scalings map boxes to boxes, multiplying volumes by $|\det T|$, so they multiply outer measure by $|\det T|$. A shear preserves each
line parallel to the $i$-th axis and acts on it by a translation; by Theorem~\ref{measure-theory:theorem-fubini} and Remark~\ref{measure-theory:remark-lebesgue-product}, integrating the one-dimensional
measures of these sections shows $\lambda(T(E)) = \lambda(E)$, and $\det T = 1$.
\end{proof}

\begin{theorem}[Change of variables]\label{measure-theory:theorem-change-of-variables}
Let $U, V \subseteq \RR^n$ be open and $\phi \colon U \to V$ a $C^1$-diffeomorphism. For every measurable $f \colon V \to [0, \infty]$, and every $f \in L^1(V)$,
\[
\int_Vf\,d\lambda = \int_Uf(\phi(x))\,|\det D\phi(x)|\,d\lambda(x).
\]
\end{theorem}

\begin{proof}
We use the norm $\|x\|_\infty$ and the corresponding operator norm; the mean value inequality (Proposition~\ref{real-analysis:proposition-derivative-basic}(4)) holds for it with the same proof. Cubes are balls
$Q = \{x : \|x - a\|_\infty \leq h\}$.

\emph{Step 1.} Let $Q \subseteq U$ be a cube with centre $a$ and $T \in \GL_n(\RR)$, and let $M = \sup_{y \in Q}\|T^{-1}D\phi(y)\|$. By the mean value inequality applied to $T^{-1}\phi$,
$\phi(Q) \subseteq \phi(a) + T(\{\|z\|_\infty \leq Mh\})$, so $\lambda(\phi(Q)) \leq |\det T|M^n\lambda(Q)$ by Lemma~\ref{measure-theory:lemma-linear-change}.

\emph{Step 2: $\lambda(\phi(Q)) \leq \int_Q|\det D\phi|$ for cubes $Q \subseteq U$.} Let $\varepsilon > 0$. Since $D\phi$ is uniformly continuous on $Q$ and $\|D\phi^{-1}\|$ is bounded there, we can divide $Q$ into
subcubes $Q_j$ with centres $x_j$ such that $\|D\phi(x_j)^{-1}D\phi(y)\| \leq 1 + \varepsilon$ for $y \in Q_j$, and $|\det D\phi|$ varies by less than $\varepsilon$ on each $Q_j$. Step 1 with $T = D\phi(x_j)$ gives
$\lambda(\phi(Q)) \leq \sum_j\lambda(\phi(Q_j)) \leq (1 + \varepsilon)^n\sum_j|\det D\phi(x_j)|\lambda(Q_j) \leq (1 + \varepsilon)^n(\int_Q|\det D\phi| + \varepsilon\lambda(Q))$.

\emph{Step 3: $\lambda(\phi(E)) \leq \int_E|\det D\phi|$ for measurable $E \subseteq U$.} An open $W \subseteq U$ is a countable union of cubes with disjoint interiors (dyadic cubes), whose boundaries are null, so the
inequality holds for $W$. Let $E \subseteq K$ with $K \subseteq U$ compact, and $C = \sup$ of $|\det D\phi|$ on a compact neighbourhood $N \subseteq U$ of $K$. By regularity choose open $W$ with $E \subseteq W \subseteq N$ and
$\lambda(W \setminus E) < \delta$; then $\lambda(\phi(E)) \leq \lambda(\phi(W)) \leq \int_E|\det D\phi| + C\delta$. This shows that $\phi$ maps null sets to null sets, hence Lebesgue measurable sets to
Lebesgue measurable sets (a measurable set is a Borel set up to a null set), and gives the inequality for $E$ in a compact subset; exhaust $U$ by compact sets for general $E$.

\emph{Step 4.} By Step 3, $\int_V\mathbf{1}_{\phi(E)} \leq \int_U\mathbf{1}_{\phi(E)} \circ \phi\,|\det D\phi|$; by linearity and monotone convergence, $\int_Vf \leq \int_U(f \circ \phi)|\det D\phi|$ for all measurable
$f \geq 0$. Applying this to $\phi^{-1}$ and the function $(f \circ \phi)|\det D\phi|$ on $U$, and using $D\phi(\phi^{-1}(y))D(\phi^{-1})(y) = 1$, gives the reverse inequality. The $L^1$ case follows by linearity.
\end{proof}

\begin{example}\label{measure-theory:example-gaussian}
With polar coordinates $(r, \theta) \mapsto (r\cos\theta, r\sin\theta)$, a diffeomorphism $(0, \infty) \times (0, 2\pi) \to \RR^2$ minus a closed half-line (a null set), with Jacobian determinant $r$,
Tonelli's theorem gives $(\int_\RR e^{-x^2}dx)^2 = \int_{\RR^2}e^{-\|x\|^2} = \int_0^{2\pi}\int_0^\infty e^{-r^2}r\,dr\,d\theta = \pi$. So $\int_\RR e^{-x^2}\,dx = \sqrt{\pi}$.
\end{example}

Fubini's theorem reduces integration by parts in several variables to one variable. When one of the functions has compact support there is no boundary term.

\begin{proposition}[Integration by parts in $\RR^n$]\label{measure-theory:proposition-integration-by-parts}
Let $f, g \colon \RR^n \to \CC$ be continuously differentiable, and suppose that one of them has compact support. Then for each $i$
\[
\int_{\RR^n}(\partial_if)\,g\,d\lambda = -\int_{\RR^n}f\,\partial_ig\,d\lambda .
\]
Consequently, for a continuously differentiable vector field $F$ and a continuously differentiable function $g$, one of them compactly supported,
$\int\operatorname{div}F\,g\,d\lambda = -\int\langle F, \nabla g\rangle\,d\lambda$, where $\operatorname{div}F = \sum_i\partial_iF^i$ and $\nabla g = (\partial_1g, \ldots, \partial_ng)$; and for $f$ of class $C^2$,
$\int(\sum_i\partial_i^2f)\,g\,d\lambda = -\int\langle\nabla f, \nabla g\rangle\,d\lambda$.
\end{proposition}

\begin{proof}
The product $h = fg$ is continuously differentiable with compact support, contained in a cube $[-R, R]^n$, and $\partial_ih = (\partial_if)g + f\,\partial_ig$ by the product rule; all these functions are
continuous with compact support, hence integrable. By Fubini's theorem (Theorem~\ref{measure-theory:theorem-fubini}) we may integrate first in the variable $x_i$, the others being fixed. For fixed
values of the other variables, $t \mapsto h(\ldots, t, \ldots)$ is continuously differentiable and vanishes for $|t| \geq R$, so by the fundamental theorem of calculus
(Theorem~\ref{real-analysis:theorem-ftc}) $\int_{-R}^R\partial_ih\,dx_i = h(\ldots, R, \ldots) - h(\ldots, -R, \ldots) = 0$. Hence $\int\partial_ih\,d\lambda = 0$, which is the formula. For the vector field,
apply it to $F^i$ and $g$ and sum over $i$; for the last formula take $F = \nabla f$.
\end{proof}

The formula says that the operator $\partial_i$ is \emph{anti-self-adjoint} for the $L^2$ inner product on compactly supported functions: $\langle\partial_if, \bar g\rangle = -\langle f, \partial_i\bar g\rangle$. It is the
reason why the derivative of a distribution is defined by $(\partial_iT)(\varphi) = -T(\partial_i\varphi)$ (Chapter~\ref{distributions}), and, applied twice, it shows that $-\sum_i\partial_i^2$ is a
non-negative operator: $\int(-\sum_i\partial_i^2f)\bar f\,d\lambda = \int|\nabla f|^2\,d\lambda \geq 0$ for compactly supported $f$ of class $C^2$.

\subsection{Volumes: Cavalieri's principle and solids of revolution}\label{measure-theory:subsection-volumes}

Fubini's theorem and the change of variables formula contain the classical recipes for computing volumes. We write $\lambda^n$ for Lebesgue measure on $\RR^n$ and call $\lambda^3(E)$ the \emph{volume}
and $\lambda^2(E)$ the \emph{area} of a set. The first recipe slices a solid by parallel planes.

\begin{proposition}[Cavalieri's principle]\label{measure-theory:proposition-cavalieri}
Let $E \subseteq \RR^n = \RR^{n-1} \times \RR$ be a Borel set, and for $t \in \RR$ let $E_t = \{y \in \RR^{n-1} : (y, t) \in E\}$ be its slice at height $t$. Then $t \mapsto \lambda^{n-1}(E_t)$ is measurable and
\[
\lambda^n(E) = \int_\RR\lambda^{n-1}(E_t)\,dt .
\]
In particular, two Borel sets whose slices at every height have the same $(n-1)$-dimensional measure have the same volume.
\end{proposition}

\begin{proof}
Apply Tonelli's theorem (Theorem~\ref{measure-theory:theorem-fubini}(1)) to $\mathbf{1}_E$ on $\RR^{n-1} \times \RR$, with $\lambda^n = \lambda^{n-1} \times \lambda^1$ on Borel sets
(Remark~\ref{measure-theory:remark-lebesgue-product}), integrating first over $y$: $\int\mathbf{1}_E(y, t)\,d\lambda^{n-1}(y) = \lambda^{n-1}(E_t)$.
\end{proof}

The principle is named after Bonaventura Cavalieri (1635), who stated it for plane figures and solids, but Archimedes had already used it, and Zu Chongzhi and his son Zu Geng used it in China in the
fifth century to find the volume of the ball.

\begin{example}\label{measure-theory:example-cavalieri}
\begin{enumerate}
\item \emph{Cones.} Let $B \subseteq \RR^2$ be a Borel set of finite area $A$ and $h > 0$, and let $C$ be the cone over $B$ with apex $(0, 0, h)$, the union of the segments from the points of $B \times \{0\}$ to the apex.
Its slice at height $t \in [0, h)$ is the scaled copy $(1 - t/h)B$, of area $(1 - t/h)^2A$ by Lemma~\ref{measure-theory:lemma-linear-change}. So $\lambda^3(C) = A\int_0^h(1 - t/h)^2\,dt = \frac13Ah$, whatever the shape of the
base: a third of the prism with the same base and height.
\item \emph{Archimedes' ball.} The slice of the half-ball $\{|x| \leq R, x_3 \geq 0\}$ at height $t \in [0, R]$ is a disc of area $\pi(R^2 - t^2)$. The slice of the cylinder $\{x_1^2 + x_2^2 \leq R^2, 0 \leq x_3 \leq R\}$ minus
the cone $\{\sqrt{x_1^2 + x_2^2} < x_3\}$ is an annulus of area $\pi R^2 - \pi t^2$. The slices have equal areas, so the half-ball has the volume of the cylinder minus the cone, $\pi R^3 - \frac13\pi R^3 = \frac23\pi R^3$, and
the ball has volume $\frac43\pi R^3$. This is Archimedes' argument, which he regarded as his finest: the ball is two thirds of the circumscribed cylinder.
\end{enumerate}
\end{example}

A \emph{solid of revolution} is obtained by rotating a plane region around an axis. We take the $z$-axis in $\RR^3$ as axis, and describe the region in the half-plane $H = \{(r, z) : r > 0\}$ of distances
from the axis and heights. For a Borel set $R \subseteq H$ let
\[
\operatorname{Rev}(R) = \bigl\{(x, y, z) \in \RR^3 : \bigl(\sqrt{x^2 + y^2}, z\bigr) \in R\bigr\}
\]
be the solid swept out by $R$, a Borel set since it is the preimage of $R$ under a continuous map on $\RR^3$ minus the axis.

\begin{proposition}[Volumes of solids of revolution]\label{measure-theory:proposition-revolution-volume}
For every Borel set $R \subseteq H$,
\[
\lambda^3(\operatorname{Rev}(R)) = 2\pi\int_Rr\,d\lambda^2(r, z) .
\]
In particular:
\begin{enumerate}
\item (\emph{Discs}) if $R = \{(r, z) : a \leq z \leq b,\ 0 < r \leq f(z)\}$ for a continuous $f \geq 0$ on $[a, b]$, then $\lambda^3(\operatorname{Rev}(R)) = \pi\int_a^bf(z)^2\,dz$;
\item (\emph{Shells}) if $R = \{(r, z) : a \leq r \leq b,\ 0 \leq z \leq g(r)\}$ for $0 < a < b$ and a continuous $g \geq 0$ on $[a, b]$, then $\lambda^3(\operatorname{Rev}(R)) = 2\pi\int_a^brg(r)\,dr$.
\end{enumerate}
\end{proposition}

\begin{proof}
Use \emph{cylindrical coordinates} $\phi(r, \theta, z) = (r\cos\theta, r\sin\theta, z)$, a $C^1$-diffeomorphism from $U = (0, \infty) \times (0, 2\pi) \times \RR$ onto $\RR^3$ minus the closed half-plane $\{y = 0, x \geq 0\}$,
which is a null set. Its Jacobian determinant is
\[
\det\begin{pmatrix}\cos\theta & -r\sin\theta & 0\\ \sin\theta & r\cos\theta & 0\\ 0 & 0 & 1\end{pmatrix} = r .
\]
Since $\mathbf{1}_{\operatorname{Rev}(R)}(\phi(r, \theta, z)) = \mathbf{1}_R(r, z)$, the change of variables formula (Theorem~\ref{measure-theory:theorem-change-of-variables}) and Tonelli's theorem give
$\lambda^3(\operatorname{Rev}(R)) = \int_U\mathbf{1}_R(r, z)\,r\,d\lambda^3 = \int_0^{2\pi}d\theta\int_Rr\,d\lambda^2 = 2\pi\int_Rr\,d\lambda^2$. For (1), integrate first over $r$: $\int_0^{f(z)}r\,dr = \frac12f(z)^2$. For
(2), integrate first over $z$: $\int_0^{g(r)}dz = g(r)$.
\end{proof}

The two special cases are the ``disc method'' and the ``shell method'' of calculus courses: (1) cuts the solid by planes perpendicular to the axis, which meet it in discs of area $\pi f(z)^2$, and is
also a case of Cavalieri's principle; (2) cuts it into thin cylindrical shells of radius $r$, height $g(r)$ and thickness $dr$, of volume $2\pi rg(r)\,dr$. For the ball, (1) with $f(z) = \sqrt{R^2 - z^2}$
gives $\pi\int_{-R}^R(R^2 - z^2)\,dz = \frac43\pi R^3$ once more.

The integral $\int_Rr\,d\lambda^2$ has a mechanical meaning. The \emph{centroid} of a Borel set $R \subseteq \RR^2$ of finite positive area is the point $(\bar r, \bar z)$ with
$\bar r = \lambda^2(R)^{-1}\int_Rr\,d\lambda^2$ and $\bar z = \lambda^2(R)^{-1}\int_Rz\,d\lambda^2$, its centre of mass for a uniform density. So the proposition can be read as follows.

\begin{corollary}[Pappus--Guldin theorem for volumes]\label{measure-theory:corollary-pappus-volume}
If $R \subseteq H$ has finite positive area and centroid at distance $\bar r$ from the axis, then
\[
\lambda^3(\operatorname{Rev}(R)) = 2\pi\bar r \cdot \lambda^2(R):
\]
the volume of a solid of revolution is the area of the region times the length of the circle described by its centroid.
\end{corollary}

\begin{proof}
$2\pi\int_Rr\,d\lambda^2 = 2\pi\bar r\lambda^2(R)$ by the definition of $\bar r$.
\end{proof}

The theorem was stated by Pappus of Alexandria around 320 and rediscovered by Paul Guldin in 1640; in Italian it is the \emph{teorema di Pappo--Guldino}. It turns volume computations into centroid
computations, and conversely.

\begin{example}\label{measure-theory:example-pappus-volume}
\begin{enumerate}
\item \emph{The torus.} Rotating the disc of radius $a$ centred at $(c, 0)$, with $0 < a < c$, gives a solid torus. The centroid is the centre, by symmetry, so the volume is $2\pi c \cdot \pi a^2 = 2\pi^2ca^2$: the
volume of a cylinder of radius $a$ whose length is the central circle of the torus. Computing it with discs requires the integral $\int_{-a}^a\pi\bigl((c + \sqrt{a^2 - z^2})^2 - (c - \sqrt{a^2 - z^2})^2\bigr)dz =
4\pi c\int_{-a}^a\sqrt{a^2 - z^2}\,dz$, which is the same thing, $4\pi c \cdot \frac12\pi a^2$.
\item \emph{The centroid of a half-disc.} Conversely, rotating the half-disc $\{r^2 + z^2 \leq a^2, r > 0\}$ around the $z$-axis gives the ball of radius $a$, so $\frac43\pi a^3 = 2\pi\bar r \cdot \frac12\pi a^2$, and the
centroid of a half-disc lies at distance $\bar r = \frac{4a}{3\pi}$ from its diameter.
\end{enumerate}
\end{example}

\section{Exercises}\label{measure-theory:section-exercises}

\begin{exercise}\label{measure-theory:exercise-cantor}
Show that the Cantor set, obtained from $[0, 1]$ by repeatedly removing open middle thirds, is compact, uncountable and of Lebesgue measure zero.
\end{exercise}

\begin{solution}
It is an intersection of closed subsets of $[0, 1]$, hence compact. After $k$ steps it is covered by $2^k$ intervals of length $3^{-k}$, so its measure is at most $(2/3)^k$. Its points are the numbers
$\sum_k a_k3^{-k}$ with $a_k \in \{0, 2\}$, and distinct sequences give distinct numbers, so it is in bijection with $\{0, 1\}^{\NN}$, which is uncountable (Theorem~\ref{sets:theorem-cantor}).
\end{solution}

\begin{exercise}\label{measure-theory:exercise-borel-cantelli}
(Borel--Cantelli.) If $\sum_k\mu(A_k) < \infty$, show that almost every point lies in only finitely many $A_k$.
\end{exercise}

\begin{solution}
By monotone convergence $\int\sum_k\mathbf{1}_{A_k} = \sum_k\mu(A_k) < \infty$, so $\sum_k\mathbf{1}_{A_k} < \infty$ a.e.
\end{solution}

\begin{exercise}\label{measure-theory:exercise-lp-inclusions}
Show that $L^q(\mu) \subseteq L^p(\mu)$ for $p \leq q$ if $\mu(X) < \infty$, and that there is no inclusion in either direction for Lebesgue measure on $\RR$.
\end{exercise}

\begin{solution}
By H\"older with exponents $q/p$ and its conjugate, $\int|f|^p \leq (\int|f|^q)^{p/q}\mu(X)^{1 - p/q}$. On $\RR$, $|x|^{-1/2}\mathbf{1}_{[-1,1]}$ is in $L^1$ but not $L^2$, and $(1 + |x|)^{-1}$ is in $L^2$ but not $L^1$.
\end{solution}

\begin{exercise}\label{measure-theory:exercise-convolution}
For $f \in L^1(\RR^n)$ and $g \in L^p(\RR^n)$, show using Fubini's theorem that $(f * g)(x) = \int f(x - y)g(y)\,dy$ is defined a.e. and $\|f * g\|_p \leq \|f\|_1\|g\|_p$.
\end{exercise}

\begin{exercise}\label{measure-theory:exercise-ball-volume}
Show that the volume of the unit ball in $\RR^n$ is $\omega_n = \pi^{n/2}/\Gamma(\frac n2 + 1)$, where $\Gamma(s) = \int_0^\infty t^{s-1}e^{-t}\,dt$, by computing $\int_{\RR^n}e^{-\|x\|^2}$ in two ways.
\end{exercise}

\begin{exercise}\label{measure-theory:exercise-ae-measurable}
Let $(X, \mathcal{A}, \mu)$ be a complete measure space, $f$ measurable and $g = f$ almost everywhere. Show that $g$ is measurable. Deduce that for Lebesgue measure, changing a measurable function on a null set
gives a measurable function, so that ``measurable function defined almost everywhere'' makes sense.
\end{exercise}

\begin{solution}
Let $N \in \mathcal{A}$ be a null set with $f = g$ outside $N$. For $a \in \RR$, $\{g > a\} = (\{f > a\} \setminus N) \cup (\{g > a\} \cap N)$. The first set is measurable, and the second is a subset of
the null set $N$, hence measurable by completeness. So $g$ is measurable (Proposition~\ref{measure-theory:proposition-measurable-functions}(1)). Lebesgue measure is complete
(Theorem~\ref{measure-theory:theorem-lebesgue-measure}).
\end{solution}

\begin{exercise}\label{measure-theory:exercise-inner-measure}
For $A \subseteq [0, 1]$ define the \emph{inner measure} $\lambda_*(A) = \sup\{\lambda(K) : K \subseteq A \text{ compact}\}$, approximating $A$ from inside instead of from outside. Show that
$\lambda_*(A) = 1 - \lambda^*([0, 1] \setminus A)$, that $\lambda_*(A) \leq \lambda^*(A)$, and that $A$ is Lebesgue measurable if and only if $\lambda_*(A) = \lambda^*(A)$. (This was Lebesgue's original
definition of measurability.)
\end{exercise}

\begin{solution}
Compact subsets $K \subseteq A$ correspond to relatively open sets $V = [0, 1] \setminus K \supseteq [0, 1] \setminus A$, with $\lambda(K) = 1 - \lambda(V)$; and $\lambda^*(B)$ is the infimum of $\lambda(U \cap [0, 1])$
over open $U \supseteq B$ for $B \subseteq [0, 1]$ (enlarge the covering boxes, as in the proof of Theorem~\ref{measure-theory:theorem-lebesgue-measure}). This gives the formula. If $K \subseteq A \subseteq U$
then $\lambda(K) \leq \lambda^*(A) \leq \lambda(U)$, so $\lambda_* \leq \lambda^*$. If $A$ is measurable, equality is inner regularity. Conversely, if $\lambda_*(A) = \lambda^*(A)$, choose compact $K_n \subseteq A$
and open $U_n \supseteq A$ with $\lambda(K_n) \to \lambda_*(A)$ and $\lambda(U_n) \to \lambda^*(A)$; then $F = \bigcup K_n \subseteq A \subseteq G = \bigcap U_n$ are Borel sets with
$\lambda(G \setminus F) = \lambda(G) - \lambda(F) = 0$, so $A = F \cup (A \setminus F)$ with $A \setminus F \subseteq G \setminus F$ null, and $A$ is measurable by completeness.
\end{solution}

\begin{exercise}\label{measure-theory:exercise-ball-semicontinuous}
Let $\mu$ be a Radon measure on $\RR^n$ and $r > 0$. Show that $x \mapsto \mu(B(x, r))$ is lower semicontinuous and $x \mapsto \mu(\overline{B}(x, r))$ is upper semicontinuous, and that for fixed $x$ the
sphere $\{y : \|y - x\| = r\}$ has $\mu$-measure $0$ for all but countably many $r$.
\end{exercise}

\begin{solution}
Let $x_k \to x$. If $\|y - x\| < r$, then $\|y - x_k\| < r$ for $k$ large, so $\mathbf{1}_{B(x, r)} \leq \liminf_k\mathbf{1}_{B(x_k, r)}$, and Fatou's lemma
(Corollary~\ref{measure-theory:corollary-integral-properties}) gives $\mu(B(x, r)) \leq \liminf_k\mu(B(x_k, r))$. If $\|y - x_k\| \leq r$ for infinitely many $k$, then $\|y - x\| \leq r$; so
$\limsup_k\mathbf{1}_{\overline{B}(x_k, r)} \leq \mathbf{1}_{\overline{B}(x, r)}$, and Fatou's lemma applied to the non-negative functions $\mathbf{1}_{\overline{B}(x, r + 1)} - \mathbf{1}_{\overline{B}(x_k, r)}$
($k$ large), whose integrals are finite, gives $\limsup_k\mu(\overline{B}(x_k, r)) \leq \mu(\overline{B}(x, r))$. The spheres of different radii are disjoint; those of radius less than $R$ with measure at
least $1/m$ number at most $m\mu(B(x, R))$, so only countably many spheres have positive measure.
\end{solution}

\begin{exercise}\label{measure-theory:exercise-dirichlet-approximation}
(\emph{Dirichlet's approximation theorem.}) Let $x \in \RR$ and $N \geq 1$. Using the pigeonhole principle for the $N + 1$ numbers $kx - \lfloor kx\rfloor \in [0, 1)$, $k = 0, \ldots, N$, show that there
are integers $p, q$ with $1 \leq q \leq N$ and $|qx - p| < 1/N$. Deduce that an irrational $x$ has infinitely many approximations $p/q$ with $|x - p/q| < 1/q^2$.
\end{exercise}

\begin{solution}
Divide $[0, 1)$ into the $N$ intervals $[\frac{i}{N}, \frac{i+1}{N})$. Two of the $N + 1$ numbers, for $k < l$ say, lie in the same interval, so $|(l - k)x - (\lfloor lx\rfloor - \lfloor kx\rfloor)| < 1/N$; take
$q = l - k$ and $p = \lfloor lx\rfloor - \lfloor kx\rfloor$. Then $|x - p/q| < 1/(qN) \leq 1/q^2$, and also $|x - p/q| < 1/N$. Writing $p/q$ in lowest terms $p^{\prime}/q^{\prime}$ with $q^{\prime} \leq q$ keeps
$|x - p^{\prime}/q^{\prime}| < 1/q^2 \leq 1/q^{\prime 2}$. If $x$ is irrational, no fraction equals $x$, so a finite set of fractions has positive minimal distance from $x$; since letting $N \to \infty$ produces
fractions arbitrarily close to $x$, we obtain infinitely many distinct ones.
\end{solution}

\begin{exercise}\label{measure-theory:exercise-decimal-blocks}
For $x \in [0, 1)$ let $e_k(x) = \lfloor 10^kx\rfloor - 10\lfloor 10^{k-1}x\rfloor$ be its decimal digits. Show that for almost every $x$:
\begin{enumerate}
\item every finite block of decimal digits occurs infinitely often among the $e_k(x)$;
\item the frequency $\frac1n\#\{k \leq n : e_k(x) = 7\}$ tends to $\frac1{10}$.
\end{enumerate}
\end{exercise}

\begin{solution}
The analogue of Lemma~\ref{measure-theory:lemma-binary-digits} holds with the same proof: prescribing the digits at the positions of a finite set $S$ defines a set of measure $10^{-|S|}$, a union of
decimal intervals of length $10^{-K}$, $K = \max S$, and events determined by disjoint sets of positions are independent, with the product formula for integrals. (1) is then proved as in
Example~\ref{measure-theory:example-infinite-monkey}, with probability $10^{-m}$ for a block of length $m$. For (2) let $\rho_k = \mathbf{1}_{\{e_k = 7\}} - \frac1{10}$, so $\int\rho_k = 0$ and $|\rho_k| \leq 1$, and
$S_n = \rho_1 + \cdots + \rho_n$. In $\int S_n^4 = \sum\int\rho_i\rho_j\rho_k\rho_l$, a term in which some index occurs exactly once vanishes by the product formula. The remaining terms have all indices equal
($n$ terms) or two values each occurring twice ($3n(n - 1)$ terms), and each is at most $1$ in absolute value; so $\int S_n^4 \leq 3n^2$, and Step 2 of the proof of
Theorem~\ref{measure-theory:theorem-borel-normal} applies verbatim.
\end{solution}

\begin{exercise}\label{measure-theory:exercise-steinmetz}
Show that the intersection of the two solid cylinders $\{x^2 + y^2 \leq R^2\}$ and $\{x^2 + z^2 \leq R^2\}$ in $\RR^3$ (the \emph{Steinmetz solid}) has volume $\frac{16}{3}R^3$, using Cavalieri's principle with
slices perpendicular to the $x$-axis.
\end{exercise}

\begin{solution}
For $|x| \leq R$ the slice is $\{(y, z) : |y| \leq \sqrt{R^2 - x^2}, |z| \leq \sqrt{R^2 - x^2}\}$, a square of side $2\sqrt{R^2 - x^2}$, and it is empty for $|x| > R$. By
Proposition~\ref{measure-theory:proposition-cavalieri} the volume is $\int_{-R}^R4(R^2 - x^2)\,dx = 4\bigl(2R^3 - \frac23R^3\bigr) = \frac{16}{3}R^3$. (Archimedes knew this too; compare with the ball inscribed in the
solid, whose slices are the inscribed discs, of area $\frac\pi4$ times that of the squares, and whose volume is $\frac\pi4 \cdot \frac{16}{3}R^3 = \frac43\pi R^3$.)
\end{solution}

\begin{exercise}\label{measure-theory:exercise-shells}
Compute the volume of the solid obtained by rotating the region $\{(r, z) : 0 < r \leq \pi, 0 \leq z \leq \sin r\}$ around the $z$-axis, using shells, and check the result with the Pappus--Guldin theorem.
\end{exercise}

\begin{solution}
By Proposition~\ref{measure-theory:proposition-revolution-volume}(2), with $a \to 0$ by monotone convergence, the volume is $2\pi\int_0^\pi r\sin r\,dr = 2\pi\bigl[\sin r - r\cos r\bigr]_0^\pi = 2\pi^2$. The region has area
$\int_0^\pi\sin r\,dr = 2$, and by symmetry about $r = \frac\pi2$ its centroid has $\bar r = \frac\pi2$; indeed $2\pi \cdot \frac\pi2 \cdot 2 = 2\pi^2$.
\end{solution}

\begin{exercise}\label{measure-theory:exercise-paraboloid}
Show that the solid paraboloid $\{(x, y, z) : x^2 + y^2 \leq z \leq h\}$ has half the volume of the circumscribed cylinder $\{x^2 + y^2 \leq h, 0 \leq z \leq h\}$, a result of Archimedes.
\end{exercise}

\begin{solution}
By the disc method with $f(z) = \sqrt z$ on $[0, h]$, the volume is $\pi\int_0^hz\,dz = \frac12\pi h^2$, and the cylinder of radius $\sqrt h$ and height $h$ has volume $\pi h \cdot h = \pi h^2$.
\end{solution}

\begin{exercise}\label{measure-theory:exercise-monotone-discontinuities}
Let $F \colon \RR \to \RR$ be nondecreasing. Show that $F$ has at most countably many points of discontinuity, using the Lebesgue--Stieltjes measure of $G(x) = F(x+) = \lim_{t \downarrow x}F(t)$.
\end{exercise}

\begin{solution}
$G$ is nondecreasing, and right-continuous: for $x < t < s$, $F(t) \leq G(t) \leq F(s)$, so $G(t) \to G(x)$ as $t \downarrow x$. At every $x$, $G(x-) = F(x-)$ (both are the limit of $F$ from the
left, as $F(t) \leq G(t) \leq F(t^{\prime})$ for $t < t^{\prime} < x$) and $G(x) = F(x+)$. A monotone function is continuous at $x$ if and only if $F(x-) = F(x+)$, so the discontinuities of $F$ are the
atoms of $\mu_G$ (Theorem~\ref{measure-theory:theorem-lebesgue-stieltjes}(4)). In $(-m, m]$ there are at most $k\,\mu_G((-m, m])$ atoms of mass $> 1/k$, a finite number; taking the union over $m$ and $k$
gives countably many atoms.
\end{solution}

\begin{exercise}\label{measure-theory:exercise-stieltjes-density}
Let $F \colon \RR \to \RR$ be nondecreasing and continuously differentiable. Show that $\mu_F = F^{\prime}\lambda$, that is, $\mu_F(E) = \int_EF^{\prime}\,d\lambda$ for every Borel set $E$.
\end{exercise}

\begin{solution}
$F^{\prime} \geq 0$ is continuous, so $E \mapsto \int_EF^{\prime}\,d\lambda$ is a Borel measure (Example~\ref{measure-theory:example-measures}(3)), finite on bounded intervals. For $a < b$,
$\int_{(a, b]}F^{\prime}\,d\lambda$ equals the Riemann integral (Proposition~\ref{measure-theory:proposition-riemann-lebesgue}; the point $a$ is null), which is $F(b) - F(a)$ by the fundamental theorem
of calculus. So the two measures agree on bounded h-intervals and are equal by Theorem~\ref{measure-theory:theorem-lebesgue-stieltjes}(1).
\end{solution}

\begin{exercise}\label{measure-theory:exercise-quadrants}
Show that two Borel probability measures $\mu$, $\nu$ on $\RR^2$ with $\mu((-\infty, x] \times (-\infty, y]) = \nu((-\infty, x] \times (-\infty, y])$ for all $x, y \in \RR$ are equal.
\end{exercise}

\begin{solution}
The quadrants $(-\infty, x] \times (-\infty, y]$, together with $\emptyset$, form a $\pi$-system: the intersection of two quadrants is the quadrant at $(\min(x, x^{\prime}), \min(y, y^{\prime}))$. It generates
$\mathcal{B}(\RR^2)$: the box $(a, b] \times (c, d]$ is the quadrant at $(b, d)$ minus the quadrants at $(a, d)$ and $(b, c)$, which overlap in the quadrant at $(a, c)$, so boxes lie in the generated
$\sigma$-algebra, and every open set is a countable union of such boxes with rational vertices. As $\mu(\RR^2) = \nu(\RR^2) = 1$, Lemma~\ref{measure-theory:lemma-dynkin} applies, with the quadrants at
$(m, m)$ increasing to $\RR^2$.
\end{solution}
